%% =====================================================================
%%  FINANZAS DE EMPRESAS — RESUMEN INTEGRAL DE LAS CINCO UNIDADES
%% ---------------------------------------------------------------------
%%  Archivo único: compilar con pdfLaTeX tres veces (índice, referencias
%%  cruzadas e índices de definiciones y fórmulas); Overleaf y latexmk
%%  repiten las pasadas automáticamente.
%%
%%  Imágenes: copiar los archivos (.jpg, .png o .pdf) con su nombre
%%  original en la misma carpeta en la que se compila este .tex (si no
%%  están ahí, también se buscan en images/ y «MFJQ images/»). Cada una
%%  se inserta con el código estándar:
%%      \begin{figure}[H]
%%        \centering
%%        \includegraphics[width=...,height=...,keepaspectratio]{nombre}
%%        \caption{Epígrafe}
%%      \end{figure}
%%  - El nombre va sin extensión, como en los apuntes originales.
%%  - Si falta un archivo, LaTeX lo avisa con un error que da su nombre.
%%
%%  Cajas: \begin{definicion}{Término}  -> bloque de diamante
%%         \begin{formulaBox}[Título]  -> bloque de oro
%%         \begin{ejemplo}[Título]     -> bloque de piedra, OCULTO:
%%           los ejemplos están en el archivo pero no se imprimen; para
%%           que vuelvan a verse, cambiar \mostrarejemplosfalse por
%%           \mostrarejemplostrue en el bloque COMPLETAR.
%%         \begin{ficha}               -> ficha de referencia (madera) al
%%           comienzo de cada unidad, OCULTA: para que vuelva a verse,
%%           cambiar \mostrarfichasfalse por \mostrarfichastrue.
%%         \begin{control}             -> control del resumen (esmeralda)
%%           al final de cada unidad, OCULTO: para que vuelva a verse,
%%           cambiar \mostrarcontrolesfalse por \mostrarcontrolestrue.
%% =====================================================================
%% =====================================================================
%%  PREÁMBULO
%% =====================================================================
\documentclass[11pt,a4paper,openany,oneside]{book}

% --- Codificación e idioma -------------------------------------------
\usepackage[utf8]{inputenc}
\usepackage[T1]{fontenc}
\usepackage[spanish,es-tabla,es-noshorthands]{babel}
\spanishplainpercent   % el espacio fino antes de % se escribe a mano: 35\,\%

% --- Tipografía: Palatino (newpx si está instalado; si no, mathpazo) ---
\IfFileExists{newpxtext.sty}
  {\usepackage{newpxtext}\usepackage{newpxmath}}
  {\usepackage[sc]{mathpazo}}
\usepackage[scaled=0.90]{helvet}
\usepackage{microtype}
\usepackage{icomma}            % coma decimal sin espacio: 0,20

% --- Página ------------------------------------------------------------
\usepackage[a4paper,top=2.5cm,bottom=2.3cm,left=2.2cm,right=2.2cm,
            headheight=24pt,headsep=14pt]{geometry}
\setlength{\parindent}{0pt}
\setlength{\parskip}{0.55em plus 0.15em minus 0.1em}
\linespread{1.04}

% --- Matemática y tablas -----------------------------------------------
\usepackage{amsmath,amssymb,mathtools}
\usepackage{enumitem}
\setlist{itemsep=2pt,topsep=3pt,parsep=1pt}
\usepackage{booktabs,array,tabularx,multirow}
\usepackage{float}
\usepackage[export]{adjustbox}  % incluye graphicx
\usepackage{caption}
\captionsetup{font=small,labelfont={bf,sf},justification=centering,
              skip=4pt}
\usepackage[table]{xcolor}
\usepackage[most]{tcolorbox}
\usepackage{tikz}
\usetikzlibrary{calc}

% --- Colores: diamante y oro (paletas propias, estilo bloque pixelado) --
\definecolor{diaX}{HTML}{073F3B}   % contorno
\definecolor{diaD}{HTML}{12746C}   % sombra
\definecolor{diaM}{HTML}{25B8A9}   % tono medio
\definecolor{diaL}{HTML}{7DF0E1}   % luz
\definecolor{diaB}{HTML}{D4FFF8}   % brillo
\definecolor{diaW}{HTML}{FFFFFF}
\definecolor{diaF}{HTML}{EFFFFC}   % fondo

\definecolor{oroX}{HTML}{5A3A00}
\definecolor{oroD}{HTML}{A8740A}
\definecolor{oroM}{HTML}{E0B020}
\definecolor{oroL}{HTML}{FFE36A}
\definecolor{oroB}{HTML}{FFF6C4}
\definecolor{oroW}{HTML}{FFFFFF}
\definecolor{oroF}{HTML}{FFFBEA}

\definecolor{piedraX}{HTML}{3A3A3A}
\definecolor{piedraD}{HTML}{6E6E6E}
\definecolor{piedraM}{HTML}{9C9C9C}
\definecolor{piedraL}{HTML}{C9C9C9}
\definecolor{piedraB}{HTML}{E6E6E6}
\definecolor{piedraF}{HTML}{F7F7F7}

\definecolor{titulos}{HTML}{0B4F4A}

% --- Arte en píxeles ----------------------------------------------------
%  \pixelarte[tamaño]{paleta}{fila1,fila2,...}
%  Cada carácter de una fila es un píxel: . = vacío; X, D, M, L, B, W =
%  colores de la paleta (por ejemplo diaX, diaM...).
\ExplSyntaxOn
\int_new:N \l__fin_fila_int
\int_new:N \l__fin_col_int
\NewDocumentCommand \pixelarte { O{1pt} m m }
  {
    \begin{tikzpicture}[x=#1,y=#1,
        baseline={([yshift=-0.55ex]current~bounding~box.center)}]
      \int_zero:N \l__fin_fila_int
      \clist_map_inline:nn {#3}
        {
          \int_zero:N \l__fin_col_int
          \tl_map_inline:nn {##1}
            {
              \str_if_eq:nnF {####1} {.}
                {
                  \fill[#2####1]
                    (\int_use:N \l__fin_col_int, -\int_use:N \l__fin_fila_int)
                    rectangle ++(1,1);
                }
              \int_incr:N \l__fin_col_int
            }
          \int_incr:N \l__fin_fila_int
        }
    \end{tikzpicture}
  }
\ExplSyntaxOff

% Gema facetada (diseño propio)
\newcommand{\gemapixel}[1][1.2pt]{\pixelarte[#1]{dia}{%
  ...XXXXX...,%
  ..XLLWMMDX..,%
  .XLWLLMMMDX.,%
  XLLLMMMMMDDX,%
  .XLMMMMMDDX.,%
  ..XMMMMDDX..,%
  ...XMMDDX...,%
  ....XMDX....,%
  .....XX.....}}
% Lingote (diseño propio)
\newcommand{\lingotepixel}[1][1.15pt]{\pixelarte[#1]{oro}{%
  ...XXXXXXXX..,%
  ..XLLWLLLMMX.,%
  .XLLLLLLMMMDX,%
  XMMMMMMMMMDDX,%
  XDDDDDDDDDDDX,%
  .XXXXXXXXXXX.}}
% Bloque de piedra (diseño propio)
\newcommand{\piedrapixel}[1][1.05pt]{\pixelarte[#1]{piedra}{%
  XXXXXXXXX,%
  XLLMMLMDX,%
  XLMMDMMDX,%
  XMMLMMDMX,%
  XMDMMLMDX,%
  XLMMDMMDX,%
  XDDDDDDDX,%
  XXXXXXXXX}}

% --- Marco pixelado con bisel (estilo bloque de videojuego) ------------
%  Cuatro filas de "píxeles" desde afuera hacia adentro: contorno oscuro
%  con esquinas escalonadas; banda media con destellos; banda clara con
%  textura; bisel interior (luz arriba-izquierda, sombra abajo-derecha).
\newlength{\finpix}\setlength{\finpix}{1.55pt}
\tikzset{pixtrazo/.style={line width=\finpix,line cap=butt,line join=miter}}
\newcommand{\marcopixel}[1]{%
  \fill[#1X] let \p1=(frame.south west), \p2=(frame.north east) in
    (\x1+\finpix,\y1) -- (\x2-\finpix,\y1) -- (\x2-\finpix,\y1+\finpix)
    -- (\x2,\y1+\finpix) -- (\x2,\y2-\finpix) -- (\x2-\finpix,\y2-\finpix)
    -- (\x2-\finpix,\y2) -- (\x1+\finpix,\y2) -- (\x1+\finpix,\y2-\finpix)
    -- (\x1,\y2-\finpix) -- (\x1,\y1+\finpix) -- (\x1+\finpix,\y1+\finpix)
    -- cycle;
  % fila 1: tono medio con destellos y luces
  \draw[pixtrazo,#1M] let \p1=(frame.south west), \p2=(frame.north east) in
    (\x1+1.5\finpix,\y1+1.5\finpix) rectangle (\x2-1.5\finpix,\y2-1.5\finpix);
  \draw[pixtrazo,#1B,dash pattern=on \finpix off 6\finpix]
    let \p1=(frame.south west), \p2=(frame.north east) in
    (\x1+1.5\finpix,\y1+1.5\finpix) rectangle (\x2-1.5\finpix,\y2-1.5\finpix);
  \draw[pixtrazo,#1L,dash pattern=on 2\finpix off 5\finpix,dash phase=4\finpix]
    let \p1=(frame.south west), \p2=(frame.north east) in
    (\x1+1.5\finpix,\y1+1.5\finpix) rectangle (\x2-1.5\finpix,\y2-1.5\finpix);
  % fila 2: tono claro con textura media
  \draw[pixtrazo,#1L] let \p1=(frame.south west), \p2=(frame.north east) in
    (\x1+2.5\finpix,\y1+2.5\finpix) rectangle (\x2-2.5\finpix,\y2-2.5\finpix);
  \draw[pixtrazo,#1M,dash pattern=on \finpix off 3\finpix,dash phase=2\finpix]
    let \p1=(frame.south west), \p2=(frame.north east) in
    (\x1+2.5\finpix,\y1+2.5\finpix) rectangle (\x2-2.5\finpix,\y2-2.5\finpix);
  % fila 3: bisel
  \draw[pixtrazo,#1B] let \p1=(frame.south west), \p2=(frame.north east) in
    (\x1+3.5\finpix,\y1+3\finpix) -- (\x1+3.5\finpix,\y2-3.5\finpix)
    -- (\x2-3\finpix,\y2-3.5\finpix);
  \draw[pixtrazo,#1D] let \p1=(frame.south west), \p2=(frame.north east) in
    (\x1+4\finpix,\y1+3.5\finpix) -- (\x2-3.5\finpix,\y1+3.5\finpix)
    -- (\x2-3.5\finpix,\y2-4\finpix);
}
\newcommand{\marcopixelsimple}[1]{% versión liviana (3 filas)
  \fill[#1X] let \p1=(frame.south west), \p2=(frame.north east) in
    (\x1+\finpix,\y1) -- (\x2-\finpix,\y1) -- (\x2-\finpix,\y1+\finpix)
    -- (\x2,\y1+\finpix) -- (\x2,\y2-\finpix) -- (\x2-\finpix,\y2-\finpix)
    -- (\x2-\finpix,\y2) -- (\x1+\finpix,\y2) -- (\x1+\finpix,\y2-\finpix)
    -- (\x1,\y2-\finpix) -- (\x1,\y1+\finpix) -- (\x1+\finpix,\y1+\finpix)
    -- cycle;
  \draw[pixtrazo,#1M] let \p1=(frame.south west), \p2=(frame.north east) in
    (\x1+1.5\finpix,\y1+1.5\finpix) rectangle (\x2-1.5\finpix,\y2-1.5\finpix);
  \draw[pixtrazo,#1L,dash pattern=on \finpix off 4\finpix]
    let \p1=(frame.south west), \p2=(frame.north east) in
    (\x1+1.5\finpix,\y1+1.5\finpix) rectangle (\x2-1.5\finpix,\y2-1.5\finpix);
  \draw[pixtrazo,#1B] let \p1=(frame.south west), \p2=(frame.north east) in
    (\x1+2.5\finpix,\y1+2\finpix) -- (\x1+2.5\finpix,\y2-2.5\finpix)
    -- (\x2-2\finpix,\y2-2.5\finpix);
  \draw[pixtrazo,#1D] let \p1=(frame.south west), \p2=(frame.north east) in
    (\x1+3\finpix,\y1+2.5\finpix) -- (\x2-2.5\finpix,\y1+2.5\finpix)
    -- (\x2-2.5\finpix,\y2-3\finpix);
}

% --- Estilo común de las cajas -------------------------------------------
%  pixelado={<código del marco>}{<código del título>}{<color de fondo>}
%  Aplica el marco pixelado también a los fragmentos de una caja que se
%  parte entre páginas (subpieles first/middle/last de tcolorbox).
\tcbset{
  bloque/.style={enhanced,breakable,
    boxrule=4\finpix,arc=0pt,outer arc=0pt,
    left=7pt,right=7pt,top=5pt,bottom=5pt,
    before skip=8pt,after skip=8pt,
    fonttitle=\sffamily\bfseries\small,
    toptitle=2.5pt,bottomtitle=3pt,titlerule=0pt},
  pixelado/.style n args={3}{%
    frame code={#1},
    interior code={\fill[#3] (interior.south west) rectangle (interior.north east);},
    title code={#2},
    skin first is subskin of={enhancedfirst}{frame code={#1},
      interior code={\fill[#3] (interior.south west) rectangle (interior.north east);},
      title code={#2}},
    skin middle is subskin of={enhancedmiddle}{frame code={#1},
      interior code={\fill[#3] (interior.south west) rectangle (interior.north east);}},
    skin last is subskin of={enhancedlast}{frame code={#1},
      interior code={\fill[#3] (interior.south west) rectangle (interior.north east);}}},
}

% Entradas de los índices de definiciones y de fórmulas (sin sangría)
\makeatletter
\renewcommand*\l@tcolorbox{\@dottedtocline{1}{0em}{1.2em}}
\makeatother

% Si el contenido de una caja empieza con \[ , se elimina la línea vacía
% que TeX deja antes de la fórmula.
\makeatletter
\newcommand{\fin@inicioformula}{%
  \setlength{\abovedisplayskip}{5pt plus 2pt minus 1pt}%
  \setlength{\abovedisplayshortskip}{5pt plus 2pt minus 1pt}%
  \setlength{\belowdisplayskip}{5pt plus 2pt minus 1pt}%
  \setlength{\belowdisplayshortskip}{5pt plus 2pt minus 1pt}%
  \@ifnextchar\[{\vspace*{-\baselineskip}}{}}
\makeatother

% Definiciones: bloque de diamante   (\begin{document}{Término})
\NewTColorBox[list inside=ldef]{definicion}{ O{} m }{bloque,
  colback=diaF,colframe=diaX,coltitle=white,colbacktitle=diaD,
  pixelado={\marcopixel{dia}}
           {\fill[diaD] (title.south west) rectangle (title.north east);
            \draw[pixtrazo,diaM] ([yshift=0.5\finpix]title.south west)
              -- ([yshift=0.5\finpix]title.south east);
            \draw[pixtrazo,diaL,dash pattern=on \finpix off 2\finpix]
              ([yshift=0.5\finpix]title.south west)
              -- ([yshift=0.5\finpix]title.south east);}
           {diaF},
  before upper*={\csname fin@inicioformula\endcsname},
  list entry={#2},
  title={\gemapixel\hspace{0.6em}#2},#1}

% Fórmulas: bloque de oro  (\begin{formulaBox}[título opcional])
\NewTColorBox[list inside=lfor]{formulaBox}{ o }{bloque,
  colback=oroF,colframe=oroX,coltitle=oroX,colbacktitle=oroL,
  pixelado={\marcopixel{oro}}
           {\fill[oroL] (title.south west) rectangle (title.north east);
            \draw[pixtrazo,oroM] ([yshift=0.5\finpix]title.south west)
              -- ([yshift=0.5\finpix]title.south east);
            \draw[pixtrazo,oroD,dash pattern=on \finpix off 2\finpix]
              ([yshift=0.5\finpix]title.south west)
              -- ([yshift=0.5\finpix]title.south east);}
           {oroF},
  before upper*={\csname fin@inicioformula\endcsname},
  IfValueT={#1}{title={\lingotepixel\hspace{0.55em}#1},
                list entry={#1}}}

% Ejemplos: bloque de piedra (discreto)  (\begin{ejemplo}[título opcional])
%  * Interruptor: con \mostrarejemplosfalse (el valor elegido en el bloque
%    COMPLETAR del comienzo del archivo) ningún ejemplo se imprime: el
%    entorno lee su contenido y lo descarta. Con \mostrarejemplostrue
%    vuelven a verse como bloques de piedra.
%  * El texto no depende de los ejemplos. Donde hacía falta un dato que
%    solo daba el ejemplo, se usa:
%      \sinejemplos{texto}  texto que solo se imprime con los ejemplos ocultos
%      \conejemplos{texto}  texto que solo se imprime con los ejemplos visibles
\newif\ifmostrarejemplos
\NewTColorBox{ejemploCaja}{ O{Ejemplo} }{bloque,
  boxrule=3\finpix,titlerule=0pt,
  colback=piedraF,colframe=piedraX,coltitle=piedraX,colbacktitle=piedraB,
  pixelado={\marcopixelsimple{piedra}}
           {\fill[piedraB] (title.south west) rectangle (title.north east);}
           {piedraF},
  before upper*={\csname fin@inicioformula\endcsname},
  fontupper=\small,
  title={\piedrapixel\hspace{0.5em}#1}}
\makeatletter
\NewDocumentEnvironment{ejemplo}{ O{Ejemplo} +b }
  {\ifmostrarejemplos\expandafter\@firstofone\else\expandafter\@gobble\fi
     {\begin{ejemploCaja}[#1]#2\end{ejemploCaja}}}
  {}
\newcommand{\sinejemplos}[1]{%
  \ifmostrarejemplos\expandafter\@gobble\else\expandafter\@firstofone\fi{#1}}
\newcommand{\conejemplos}[1]{%
  \ifmostrarejemplos\expandafter\@firstofone\else\expandafter\@gobble\fi{#1}}
\makeatother

% --- Figuras ------------------------------------------------------------
%  Cada imagen se inserta con el código estándar:
%      \begin{figure}[H]
%        \centering
%        \includegraphics[width=0.5\linewidth,height=0.42\textheight,keepaspectratio]{U1-1}
%        \caption{Epígrafe}
%      \end{figure}
%  * El nombre va sin extensión, como en los apuntes originales: LaTeX
%    prueba .pdf, .png, .jpg y .jpeg.
%  * Los archivos se buscan en la carpeta donde se compila este .tex (y,
%    si no están ahí, en las subcarpetas images/ y «MFJQ images/»).
%  * [H] (paquete float) deja cada figura exactamente en su lugar del
%    texto; width y height son los máximos y keepaspectratio conserva la
%    proporción de la imagen.
\graphicspath{{./}{./images/}{./MFJQ images/}}
\setlength{\intextsep}{8pt plus 3pt minus 2pt}
%% =====================================================================
%%  ESTILO: ficha de referencia, control del resumen, títulos, encabezados
%% =====================================================================
\usepackage{titlesec}
\usepackage{fancyhdr}
\usepackage{multicol}

% --- Paletas adicionales ----------------------------------------------
\definecolor{maderaX}{HTML}{4A2F12}
\definecolor{maderaD}{HTML}{7A5426}
\definecolor{maderaM}{HTML}{A87B43}
\definecolor{maderaL}{HTML}{C9A06A}
\definecolor{maderaB}{HTML}{E8D2A6}
\definecolor{maderaF}{HTML}{FCF8F0}

\definecolor{esmeX}{HTML}{0B3D1E}
\definecolor{esmeD}{HTML}{17803D}
\definecolor{esmeM}{HTML}{2BBF5A}
\definecolor{esmeL}{HTML}{7CEB9C}
\definecolor{esmeB}{HTML}{D2FADF}
\definecolor{esmeW}{HTML}{FFFFFF}
\definecolor{esmeF}{HTML}{F1FDF5}

% Libro (diseño propio)
\newcommand{\libropixel}[1][1.1pt]{\pixelarte[#1]{madera}{%
  .XXXXXXXXX.,%
  XLLLLXLLLLX,%
  XLBBLXLBBLX,%
  XLLLLXLLLLX,%
  XLBBLXLBBLX,%
  XLLLLXLLLLX,%
  XMMMMXMMMMX,%
  .XXXXDXXXX.}}
% Gema rómbica (diseño propio)
\newcommand{\esmeraldapixel}[1][1.2pt]{\pixelarte[#1]{esme}{%
  ....XX....,%
  ...XLWX...,%
  ..XLLMDX..,%
  .XLWMMMDX.,%
  XLLMMMMMDX,%
  .XMMMMMDX.,%
  ..XMMMDX..,%
  ...XMDX...,%
  ....XX....}}

% --- Ficha de referencia (madera) -------------------------------------
%  Interruptor: con \mostrarfichasfalse (el valor elegido en el bloque
%  COMPLETAR del comienzo del archivo) las fichas no se imprimen: el
%  entorno lee su contenido y lo descarta. Con \mostrarfichastrue vuelven
%  a verse al comienzo de cada unidad.
\newif\ifmostrarfichas
\NewTColorBox{fichaCaja}{ O{Ficha de referencia} }{bloque,
  boxrule=3\finpix,titlerule=0pt,
  colback=maderaF,colframe=maderaX,coltitle=maderaX,
  colbacktitle=maderaB,
  pixelado={\marcopixelsimple{madera}}
           {\fill[maderaB] (title.south west) rectangle (title.north east);}
           {maderaF},
  fontupper=\small,before skip=4pt,after skip=12pt,
  title={\libropixel\hspace{0.55em}#1}}
\makeatletter
\NewDocumentEnvironment{ficha}{ O{Ficha de referencia} +b }
  {\ifmostrarfichas\expandafter\@firstofone\else\expandafter\@gobble\fi
     {\begin{fichaCaja}[#1]#2\end{fichaCaja}}}
  {}
\makeatother
\newlength{\anchoetiqueta}\setlength{\anchoetiqueta}{3.3cm}
\newcommand{\fila}[2]{\par\noindent\hangindent=\anchoetiqueta\hangafter=1
  \makebox[\anchoetiqueta][l]{\sffamily\bfseries\color{maderaX}#1}#2\par
  \vspace{1pt}}
\newcommand{\Materia}{Finanzas de Empresas}
\newcommand*{\Curso}{}   % completar si se desea: \renewcommand*{\Curso}{...}
\newcommand*{\Alumno}{}  % completar si se desea: \renewcommand*{\Alumno}{...}
% Marcas por unidad en los índices de definiciones y de fórmulas
\newcommand{\listaunidad}[1]{\par\addvspace{5pt}%
  {\sffamily\bfseries\small\color{titulos}Unidad #1}\par\nopagebreak\vspace{1pt}}
\newcommand{\marcaunidad}{%
  \addtocontents{ldef}{\protect\listaunidad{\thechapter}}%
  \addtocontents{lfor}{\protect\listaunidad{\thechapter}}}
\newcommand{\fichabase}[1]{%
  \fila{Materia}{\Materia}%
  \ifx\Curso\empty\else\fila{Curso}{\Curso}\fi
  \fila{Unidad temática}{#1}}

% --- Control del resumen (esmeralda) ----------------------------------
%  Interruptor: con \mostrarcontrolesfalse (el valor elegido en el bloque
%  COMPLETAR del comienzo del archivo) los controles no se imprimen: el
%  entorno lee su contenido y lo descarta. Con \mostrarcontrolestrue
%  vuelven a verse al final de cada unidad.
\newif\ifmostrarcontroles
\NewTColorBox{controlCaja}{ O{Control del resumen} }{bloque,
  colback=esmeF,colframe=esmeX,coltitle=white,
  colbacktitle=esmeD,
  pixelado={\marcopixel{esme}}
           {\fill[esmeD] (title.south west) rectangle (title.north east);
            \draw[pixtrazo,esmeM] ([yshift=0.5\finpix]title.south west)
              -- ([yshift=0.5\finpix]title.south east);}
           {esmeF},
  before skip=14pt,
  title={\esmeraldapixel\hspace{0.55em}#1}}
\makeatletter
\NewDocumentEnvironment{control}{ O{Control del resumen} +b }
  {\ifmostrarcontroles\expandafter\@firstofone\else\expandafter\@gobble\fi
     {\begin{controlCaja}[#1]#2\end{controlCaja}}}
  {}
\makeatother
\newcommand{\pregunta}[1]{\par\smallskip\noindent{\sffamily\bfseries\color{esmeX}#1}\par\nopagebreak}

% --- Franja pixelada para títulos de capítulo ---------------------------
\ExplSyntaxOn
\tl_new:N \l__fin_franjaA_tl
\tl_new:N \l__fin_franjaB_tl
\tl_set:Nx \l__fin_franjaA_tl { \prg_replicate:nn {12} {XDMLBLMDXMLDDMLX} }
\tl_set:Nx \l__fin_franjaB_tl { \prg_replicate:nn {12} {DMLBWBLMDLBMMLBD} }
\tl_new:N \l__fin_franja_tl
\tl_set:Nx \l__fin_franja_tl
  { \exp_not:N \pixelarte [2pt] {dia} { \l__fin_franjaA_tl , \l__fin_franjaB_tl } }
\NewDocumentCommand \pixelfranja { }
  {
    \par\noindent
    \resizebox{\linewidth}{!}{ \tl_use:N \l__fin_franja_tl }
    \par
  }
\ExplSyntaxOff

% --- Títulos ---------------------------------------------------------------
\titleformat{\chapter}[display]
  {\normalfont\sffamily\bfseries\color{titulos}}
  {\pixelfranja\vspace{6pt}%
   {\large\color{diaD}\MakeUppercase{Unidad}~\thechapter}}
  {4pt}
  {\Huge\raggedright}
  [\vspace{6pt}\pixelfranja\marcaunidad]
\titleformat{name=\chapter,numberless}[display]
  {\normalfont\sffamily\bfseries\color{titulos}}
  {}
  {0pt}
  {\pixelfranja\vspace{8pt}\Huge\raggedright}
  [\vspace{6pt}\pixelfranja]
\titlespacing*{\chapter}{0pt}{-10pt}{18pt}
\titleformat{\section}
  {\normalfont\Large\sffamily\bfseries\color{titulos}}
  {\gemapixel[0.95pt]\hspace{0.45em}\thesection}{0.6em}{}
\titleformat{\subsection}
  {\normalfont\large\sffamily\bfseries\color{titulos}}
  {\thesubsection}{0.7em}{}
\titleformat{\subsubsection}
  {\normalfont\normalsize\sffamily\bfseries\color{diaD}}
  {}{0em}{}
\titlespacing*{\section}{0pt}{16pt plus 4pt minus 2pt}{8pt plus 2pt}
\titlespacing*{\subsection}{0pt}{12pt plus 3pt minus 2pt}{5pt plus 2pt}
\titlespacing*{\subsubsection}{0pt}{9pt plus 2pt minus 2pt}{3pt}
\setcounter{secnumdepth}{2}
\setcounter{tocdepth}{2}

% --- Encabezado y pie --------------------------------------------------------
\pagestyle{fancy}
\fancyhf{}
\renewcommand{\chaptermark}[1]{\markboth{Unidad \thechapter.\ #1}{}}
\renewcommand{\sectionmark}[1]{\markright{\thesection\ #1}}
\fancyhead[L]{\small\sffamily\color{diaD}Finanzas de Empresas · Resumen integral}
\fancyhead[R]{\small\sffamily\color{diaD}\nouppercase{\leftmark}}
\fancyfoot[C]{\small\sffamily\color{diaD}\thepage}
\renewcommand{\headrulewidth}{0.4pt}
\renewcommand{\headrule}{\hbox to\headwidth{\color{diaM}\leaders\hrule height \headrulewidth\hfill}}
\fancypagestyle{plain}{\fancyhf{}\fancyfoot[C]{\small\sffamily\color{diaD}\thepage}%
  \renewcommand{\headrulewidth}{0pt}}

% --- Referencias cruzadas ------------------------------------------------------
\newcommand{\vease}[1]{(véase §\,\ref{#1})}

% --- COMPLETAR (opcional) ----------------------------------------------------
% Si se completan, el curso y el nombre del alumno aparecen en la portada
% (y el curso, también en la ficha de cada unidad cuando las fichas se ven).
\renewcommand*{\Curso}{}
\renewcommand*{\Alumno}{}
% Ejemplos numéricos (bloques de piedra): ocultos. Para que vuelvan a verse,
% cambiar la línea siguiente por \mostrarejemplostrue
\mostrarejemplosfalse
% Fichas de referencia del comienzo de cada unidad: ocultas. Para que vuelvan
% a verse, cambiar la línea siguiente por \mostrarfichastrue
\mostrarfichasfalse
% Controles del resumen (las cinco preguntas al final de cada unidad): ocultos.
% Para que vuelvan a verse, cambiar la línea siguiente por \mostrarcontrolestrue
\mostrarcontrolesfalse
% ------------------------------------------------------------------------------

\usepackage[hidelinks]{hyperref}
\hypersetup{pdftitle={Finanzas de Empresas: resumen integral de las cinco unidades},
            pdfsubject={Finanzas de Empresas},
            pdfkeywords={finanzas corporativas, flujo de fondos, estructura de capital,
                         costo de capital, dividendos, valuación, riesgo, CAPM}}

\begin{document}

%% =====================================================================
%%  PORTADA Y GUÍA DE LECTURA
%% =====================================================================
\begin{titlepage}
\pagecolor{white}
\centering
\pixelfranja
\vspace*{1.6cm}

{\gemapixel[3.4pt]\hspace{1.4em}\lingotepixel[3.4pt]\hspace{1.4em}\esmeraldapixel[3.4pt]\hspace{1.4em}\libropixel[3.4pt]\hspace{1.4em}\piedrapixel[3.4pt]\par}
\vspace{1.4cm}

{\sffamily\bfseries\color{titulos}\fontsize{38}{44}\selectfont Finanzas de Empresas\par}
\vspace{0.5cm}
{\sffamily\Large\color{diaD}Resumen integral de las cinco unidades\par}
\vspace{0.35cm}
{\sffamily\normalsize\color{piedraD}\ifmostrarejemplos Teoría, fórmulas, ejemplos e imágenes\else Teoría, fórmulas e imágenes\fi{} de toda la bibliografía, sin repeticiones\par}
\vspace{1.2cm}

\begin{tcolorbox}[bloque,breakable=false,width=0.84\linewidth,
  colback=piedraF,colframe=piedraX,
  pixelado={\marcopixelsimple{piedra}}{}{piedraF},
  left=12pt,right=12pt,top=9pt,bottom=9pt,fontupper=\sffamily]
\renewcommand{\arraystretch}{1.55}
\begin{tabular}{@{}l@{\hspace{0.9em}}l@{}}
\gemapixel & \textbf{Unidad 1.} Empresa y modelos financieros\\
\lingotepixel & \textbf{Unidad 2.} Cash flow y planificación financiera\\
\esmeraldapixel & \textbf{Unidad 3.} Estructura de capital, costo de capital y punto de equilibrio\\
\libropixel & \textbf{Unidad 4.} Política de dividendos y valuación de bonos y acciones\\
\piedrapixel & \textbf{Unidad 5.} Decisiones de inversión, riesgo y costo del capital\\
\end{tabular}
\end{tcolorbox}

\vfill
{\sffamily\small\color{titulos}
\ifx\Curso\empty\else Curso: \Curso\par\smallskip\fi
\ifx\Alumno\empty\else \Alumno\par\smallskip\fi}
{\sffamily\footnotesize\color{piedraD}Elaborado a partir de la bibliografía de la cátedra: Fornero; Dumrauf; Ross, Westerfield y Jaffe;\\
Brealey, Myers y Allen; Gitman; apuntes de clase y material complementario.\par}
\vspace{0.8cm}
\pixelfranja
\end{titlepage}

%% ---------------------------------------------------------------------

\tableofcontents

%% =====================================================================
%%  UNIDAD 1
%% =====================================================================
\chapter{Empresa y modelos financieros}\label{cap:u1}

\begin{ficha}
\fichabase{1 — Empresa y modelos financieros}
\fila{Texto base}{R.\,A. Fornero, Parte~I «Empresa y modelos financieros», cap.~1 «Alcance de la función financiera» (texto de cátedra).}
\fila{}{Nota de clase «Flujo de fondos disponibles para la empresa (\emph{Free Cash Flow to the Firm})».}
\fila{Complementaria}{R. Brealey, S. Myers y F. Allen, \emph{Principios de finanzas corporativas}, McGraw-Hill: cap.~1 «Objetivos y gobernanza de la empresa» y cap.~2 «Cómo calcular los valores actuales».}
\fila{}{S. Ross, R. Westerfield y J. Jaffe, \emph{Finanzas corporativas}, McGraw-Hill: apuntes de repaso sobre formas de organización, mercados financieros y VPN, y valor del dinero en el tiempo.}
\fila{}{Apuntes de cátedra de matemática financiera básica (progresión geométrica, rentas, ecuación de Fisher).}
\end{ficha}
%% ---------------------------------------------------------------------
\section{Alcance de la función financiera}\label{sec:u1-alcance}

\subsection{Finanzas: capital y circulación de dinero}

Hablar de finanzas evoca el dinero, que puede gastarse o invertirse; en ambos casos, circula. Sin embargo, el dinero no es el elemento fundamental de las finanzas, sino solo su vehículo primario: la cantidad de dinero de una economía es pequeña frente a todo lo que en ella se gasta e invierte. Lo que en verdad permanece es la satisfacción por los bienes consumidos y los bienes de inversión, es decir, el \emph{capital}, que es potencial de consumo futuro.

Cada transacción de bienes tiene una contrapartida en dinero. En una economía de dos sectores, familias y empresas, el dinero recorre un circuito sin fin: las familias lo reciben por su trabajo y su capital y lo devuelven al comprar los bienes finales que producen las empresas. También circula dentro de cada sector, porque las empresas intercambian bienes intermedios y las familias se prestan servicios personales. Así se enlazan los flujos de todos los participantes: una embotelladora de agua vende la mayor parte de su producción no a sus empleados, sino a quienes trabajan en las empresas que fabrican sus botellas, etiquetas y máquinas, o los insumos de estas.

\begin{figure}[H]
  \centering
  \includegraphics[width=0.5\linewidth,height=0.42\textheight,keepaspectratio]{U1-1}
  \caption{Circuito de bienes (línea continua) y de dinero (línea cortada) entre los sectores de la economía.}
\end{figure}

A ese esquema le falta algo esencial: supone que familias y empresas gastan o invierten todo el dinero que obtienen. Cada participante está en un \emph{equilibrio financiero primario}, propio de una economía sin crédito ni títulos. La consecuencia es doble. Cada empresa solo puede aprovechar sus oportunidades de inversión con los excedentes que ella misma genera, y quien ahorra solo puede invertir en las oportunidades que él mismo advierte (por ejemplo, una familia que compra un inmueble para alquilarlo a una empresa). Por lo tanto, faltan recursos para las oportunidades más innovadoras y la economía funciona por debajo de su potencial de bienestar.

\subsection{El sector financiero y la función financiera en la empresa}

Algunos participantes tienen excedentes, es decir, ingresos mayores que su consumo e inversión; otros tienen oportunidades de inversión que superan sus fondos. Las transacciones financieras conectan a unos con otros.

El sector financiero elemental son los bancos. Quien tiene excedentes puede hacer una inversión financiera en un banco en lugar de invertir en bienes, y el banco destina esos recursos a empresas cuyas oportunidades requieren más dinero del que generan. Así, la empresa invierte en el momento que juzga oportuno y no cuando acumula el excedente; las familias pueden consumir más hoy; hay más transacciones y, en consecuencia, más excedente disponible para nuevas inversiones. Como la intermediación bancaria consiste sobre todo en dinero, sus flujos se dibujan con línea cortada.

\begin{figure}[H]
  \centering
  \includegraphics[width=0.5\linewidth,height=0.42\textheight,keepaspectratio]{U1-2}
  \caption{El circuito económico con el sector financiero como intermediario entre quienes tienen excedentes y quienes tienen oportunidades de inversión.}
\end{figure}

¿Por qué alguien prefiere un depósito a interés a un inmueble para alquilar? Quizá acumula fondos para una inversión futura, espera más consumo futuro o busca más libertad para decidir después. Toda inversión financiera se juzga, por eso, en dos dimensiones.

\begin{definicion}{Rendimiento y flexibilidad de una inversión financiera}
El \textbf{rendimiento} es la posibilidad de un consumo futuro mayor que el consumo actual al que se renuncia (renta o interés). La \textbf{flexibilidad} abarca la \emph{divisibilidad}, es decir, poder invertir importes variados, y la \emph{liquidez}, es decir, poder recuperar lo invertido con relativa rapidez.
\end{definicion}

Una empresa, en cambio, toma fondos prestados porque espera un rendimiento de sus actividades futuras. Ese rendimiento depende de su estrategia competitiva y de la calidad de su dirección, y una de esas cualidades es el modo en que desarrolla la función financiera.

\begin{definicion}{Función financiera}
Es la interacción de tres clases de decisiones de la empresa:
\begin{itemize}
\item de \textbf{inversión}: cuánto y cuándo invertir en los recursos que hacen funcionar el negocio (maquinaria, existencias de insumos y productos, posición comercial);
\item de \textbf{financiamiento}: cómo obtener el dinero para esos recursos. Sin sector financiero solo cabe decidir si los excedentes se reinvierten o se entregan a quienes aportaron el capital; con él, además, cuándo conviene endeudarse y hasta cuánto;
\item de \textbf{información estratégica}: qué se comunica a financiadores e inversores sobre qué es la empresa, su posición de mercado y sus perspectivas.
\end{itemize}
\end{definicion}

Estas decisiones condicionan el perfil competitivo de la empresa, su potencial de crecimiento y su capacidad para retribuir a inversores y prestamistas en la medida que estos requieren.

La literatura de finanzas corporativas usa una terminología paralela. Las \emph{decisiones de inversión} consisten en comprar activos reales —el presupuesto de capital o CAPEX— para obtener rendimientos futuros; las \emph{decisiones de financiamiento}, en conseguir los fondos para adquirirlos. Elegir entre deuda y capital propio define la \emph{estructura de capital}, y el capital propio se obtiene emitiendo acciones nuevas o reinvirtiendo el flujo que generan los activos existentes. Pagar dividendos o recomprar acciones es, por su parte, una \emph{decisión de distribución}.

\begin{definicion}{Gestor financiero y costo de oportunidad del capital}
Es \textbf{gestor financiero} toda persona responsable de una decisión de inversión o de financiamiento. Su papel es mediar entre los mercados financieros, de donde provienen los fondos, y las operaciones de la empresa, que los convierten en ganancias. El \textbf{costo de oportunidad del capital} es el rendimiento mínimo que exigen los inversores para financiar un proyecto: el que obtendrían en alternativas de riesgo comparable. Un proyecto conviene si su rendimiento lo supera.
\end{definicion}

El gestor crea valor por dos vías: comprando activos que generan más efectivo del que cuestan y emitiendo bonos, acciones u otros instrumentos que recaudan más efectivo del que cuestan. En síntesis, la empresa debe generar más efectivo del que utiliza: lo que paga a acreedores y accionistas debe superar lo que estos le aportaron.

\subsection{Mercados financieros, títulos y gobierno}

Para invertir más de lo que genera, la empresa tiene dos caminos: el préstamo bancario o la emisión de títulos de capital (acciones) o de crédito (por ejemplo, bonos). La emisión es una relación directa entre la empresa y quienes tienen excedentes.

\begin{definicion}{Título y mercado financiero}
Un \textbf{título} es una construcción contractual, una relación jurídica que otorga derechos sobre los bienes y los excedentes de la empresa emisora; por eso vale según el desempeño de esa empresa y de la economía en que crece. Un \textbf{mercado financiero} es el ámbito relativamente organizado en el que se relacionan los emisores de títulos (de capital o de deuda) con quienes desean invertir en ellos. A diferencia del sector financiero, no es un intermediario que preste con fondos propios o captados: solo vincula a las partes.
\end{definicion}

\begin{definicion}{Transacción primaria y transacción secundaria}
En la \textbf{transacción primaria} la empresa emite el título y recibe el dinero del inversor que lo suscribe. En la \textbf{transacción secundaria} el título pasa de un inversor a otro, sin intervención de la empresa.
\end{definicion}

Si un título no cotiza, liquidarlo exige encontrar a otro inversor interesado, lo que resulta complicado. Con mercados secundarios desarrollados, en cambio, el plazo durante el cual la empresa dispone de los fondos se independiza del plazo que desea cada inversor: acciones y bonos son financiamiento de largo plazo para la empresa aunque algunos inversores los conserven poco tiempo. Los mercados alcanzan así la máxima flexibilidad para la inversión financiera.

No todos los mercados ni todos los títulos son igualmente líquidos. Por eso las decisiones de financiamiento y de información de la empresa influyen en la flexibilidad de sus inversores y, en consecuencia, en las condiciones de su propio financiamiento. En las economías poco desarrolladas solo algunas empresas acceden al mercado; predominan el crédito bancario y la colocación negociada de títulos, como el aporte de capital de un nuevo socio.

La principal ventaja del mercado bursátil es precisamente permitir la negociación de títulos corporativos y, con ello, la obtención de capital. Los mercados secundarios adoptan dos formas:
\begin{itemize}
\item \textbf{de subasta}, organizados y con sede física, donde cotizan las acciones de las grandes empresas. El caso típico es la Bolsa de Nueva York (NYSE), que concentra más del 85\,\% de las acciones negociadas en esta modalidad, seguida por la American Stock Exchange y las bolsas regionales;
\item \textbf{de distribuidores} o extrabursátiles (\emph{over the counter}, OTC), sin sede física, donde se negocia la mayor parte de la deuda y algunas acciones. Su exponente es el NASDAQ, sistema electrónico de cotizaciones creado en 1971 que reúne unas dos veces más compañías que la NYSE, en general más pequeñas (con excepciones como Microsoft o Intel) y de valor total muy inferior.
\end{itemize}
Fuera de Estados Unidos sobresalen las bolsas de Tokio y Londres. Como el mercado OTC no tiene fronteras físicas, existe un enorme mercado internacional de deuda y, con la globalización, la negociación de muchos instrumentos nunca se detiene: simplemente viaja alrededor del mundo.

Al esquema económico le faltan el exterior y el gobierno. Dejando de lado el primero, el gobierno pesa mucho en lo financiero por dos motivos. Primero, provee el marco de regulaciones y servicios en que se desarrollan los negocios. Las transacciones financieras son relaciones jurídicas hechas de tiempo e información, que duran a veces largos períodos con participantes para quienes la empresa es bastante opaca; por eso requieren una regulación significativa, cuya efectividad se mide por la confiabilidad del sistema y la prevención de conductas disfuncionales. Segundo, configura los negocios mediante los impuestos. Los sistemas tributarios modernos son complejos y, sin llegar a la evasión, ofrecen muchas formas de atenuar o demorar su impacto; por eso son relevantes para la función financiera.

El diagrama completo incluye los mercados financieros —que no son un sector, sino un ámbito de transacciones regulado por el gobierno o autorregulado— y al gobierno, que compra bienes y factores a familias y empresas, realiza transferencias (previsionales, asistenciales) y se financia sobre todo con impuestos. En definitiva, cada sector y cada participante es un sistema de captación y distribución de dinero: lo recibe de unos por sus prestaciones y lo canaliza hacia otros en un circuito permanente.

\subsection{La empresa como función de transformación}

Decidir exige construir modelos del ámbito que la decisión afectará y comparar alternativas en un entorno virtual, para elegir la que mejor cumple el objetivo. Desde esta perspectiva, la empresa es una \emph{función de transformación} física e informacional cuyo objetivo es maximizar alguna medida de valor.

Lo visible es la transformación física: la planta, los almacenes, las máquinas, las instalaciones. Pero todo eso funciona porque existe «otro sistema empresa»: el de información y decisión, que en lenguaje simbólico identifica los cursos de acción y «decide» transformar insumos en productos, costos en beneficios y egresos en ingresos. La transformación se acepta si los sacrificios (insumos, costos, desembolsos) quedan compensados, o más que compensados, por las satisfacciones (productos, ingresos, cobros).

\begin{figure}[H]
  \centering
  \includegraphics[width=0.5\linewidth,height=0.42\textheight,keepaspectratio]{U1-3}
  \caption{Dimensiones de la empresa como función de transformación: sacrificios frente a satisfacciones.}
\end{figure}

Como ciertos elementos (materias primas, dinero) pierden su naturaleza para convertirse en productos, se necesitan medidas que anticipen si el valor $X$ de lo que «se pierde» será inferior al valor $Y$ de lo que se obtiene. Esa información se procesa en tres dimensiones.

\begin{definicion}{Las tres dimensiones de la medición en la empresa}
\begin{enumerate}[label=\arabic*)]
\item \textbf{Productividad}: relación física entre salidas e insumos (unidades adquiridas, procesadas, almacenadas, vendidas o cobradas); por ejemplo, la productividad de la tierra, del personal o de las máquinas.
\item \textbf{Rentabilidad}: las mismas entradas y salidas medidas en términos económicos, con cantidades ($Q$) y precios ($P$), es decir, como costos e ingresos.
\item \textbf{Disponibilidad de fondos}: costos e ingresos no coinciden en el tiempo con pagos y cobros; el desfase que crean los plazos exige una gestión de caja o \emph{cash flow}.
\end{enumerate}
\end{definicion}

Las tres pueden divergir. Una productividad creciente puede convivir con una rentabilidad decreciente si los precios evolucionan en contra de la empresa; y aun con buena productividad y buena rentabilidad pueden surgir problemas si no se pueden pagar a tiempo los compromisos.

\begin{figure}[H]
  \centering
  \includegraphics[width=0.5\linewidth,height=0.42\textheight,keepaspectratio]{U1-4}
\end{figure}
\begin{figure}[H]
  \centering
  \includegraphics[width=0.5\linewidth,height=0.42\textheight,keepaspectratio]{U1-5}
  \caption{Productividad, rentabilidad y disponibilidad de fondos: una evolución física favorable puede acompañarse de resultados económicos o financieros desfavorables.}
\end{figure}

Las tres mediciones son imprescindibles, pero la financiera es la que evita los tropiezos de una conversión inadecuada, insuficiente o inoportuna del dinero en el ciclo de adquisición, producción, almacenamiento, venta y cobro. La razón es que solo el dinero cancela deudas, y solo cuando un activo vuelve a ser dinero puede reasignarse. Por eso todo resultado expresado en otros términos, como la ganancia de un período, es provisorio hasta que se recupera lo invertido. La medida financiera del desempeño compara el dinero destinado a los recursos con el dinero recuperado: una inversión conviene si el segundo supera al primero, teniendo en cuenta el tiempo transcurrido y las oportunidades perdidas mientras tanto.

La gestión tiene, además, otras consecuencias: impacto ambiental, satisfacción de los clientes, capacidad de innovación y aprendizaje. Todas ellas, junto con las tres dimensiones, deben traducirse a un formato común que contemple el desarrollo de las actividades en el tiempo: el \textbf{valor del capital invertido por los propietarios}. Un sistema de información y decisión —hecho de métodos, modelos, algoritmos, preferencias, teorías y convenciones— traduce los distintos aspectos de la empresa a ese valor, que resume el efecto de las decisiones sobre las ganancias, la posición financiera, la exposición a riesgos (ineludibles si se pretende ganar), la ventaja competitiva, su duración y las oportunidades de crecimiento.

Maximizar ese valor implica elegir, y la elección se basa en expectativas, no en resultados ya observados. Se trabaja entonces con representaciones abstractas que anticipan los resultados de cada alternativa, una experimentación virtual del tipo «¿qué pasaría si…?». Se optimiza ese planteo, se ejecuta la decisión elegida y se controla su plan. Los resultados no se observan en el campo o en la planta, sino en modelos de las dimensiones relevantes para el valor, que terminan visualizándose en la pantalla de una computadora.

\begin{figure}[H]
  \centering
  \includegraphics[width=0.5\linewidth,height=0.42\textheight,keepaspectratio]{U1-6}
  \caption{Las decisiones se evalúan sobre modelos de las dimensiones relevantes para el valor y no sobre la observación directa de las transformaciones físicas.}
\end{figure}

\subsection{Objetivo financiero de la empresa}\label{sec:u1-objetivo}

\begin{definicion}{Objetivo financiero de la empresa}
Maximizar el valor del capital invertido por los propietarios, forma instrumental de considerar la riqueza que ellos obtienen por el capital que comprometieron. Maximizar significa procurar el mayor valor posible en un mundo incierto y con información parcial: el valor resulta de muchas variables de efectos a menudo contrapuestos y solo se observa con precisión en algunos momentos; aun así, debe orientar cada decisión hacia una supervivencia próspera.
\end{definicion}

El objetivo se precisa con tres preguntas.

\subsubsection{¿Por qué el valor del capital de los propietarios?}
En la empresa participan diversos interesados legítimos (\emph{stakeholders}): propietarios, prestamistas, directivos, personal, clientes, proveedores y gobierno. Cada uno interviene en el flujo de bienes, servicios y dinero buscando el intercambio más conveniente, dentro de relaciones jurídicas más o menos especificadas: un contrato de préstamo está muy detallado; la relación laboral, menos; otras surgen de disposiciones del gobierno o de reclamos contingentes, como los de un tercero perjudicado por un producto.

Los propietarios conservan el poder político —la asamblea de accionistas— y soportan el riesgo remanente del negocio. Por eso su valor es consecuencia de cómo se organizan los intercambios con los demás participantes, que actúan como restricciones. Ese valor depende de la lealtad de los clientes, la realización del personal, el cumplimiento de las regulaciones (productos, ambiente, impuestos), las relaciones fluidas de financiamiento y la capacidad de innovar: es un valor de largo plazo. Esta perspectiva no excluye que, cuando el éxito depende del capital humano, parte de la retribución sea una participación en el valor, como la compensación con acciones a directivos o empleados.

\begin{figure}[H]
  \centering
  \includegraphics[width=0.5\linewidth,height=0.42\textheight,keepaspectratio]{U1-7}
  \caption{Los interesados en la empresa (\emph{stakeholders}) y el valor del capital invertido por los propietarios.}
\end{figure}

\subsubsection{¿Qué representa ese valor?}
El capital invertido es el dinero que los propietarios aportaron y mantienen en la empresa; el valor es cuánto debería reconocérseles por la propiedad. La diferencia, el \emph{valor diferencial}, mide lo que generaron la estrategia de negocios y la forma de dirigir la empresa: los lineamientos competitivos y las decisiones de inversión, financiamiento e información.

\begin{ejemplo}[Capital invertido y valor para los propietarios]
Los propietarios aportaron y mantienen \$\,2~millones y la propiedad vale \$\,3,5~millones: el valor diferencial creado es de \$\,1,5~millones.
\end{ejemplo}

\subsubsection{¿Cómo se manifiesta?}
El valor tiene sentido en dinero: es lo que alguien pagaría por participar en la propiedad. Si las acciones cotizan, su precio es una medida instantánea que cambia con las expectativas de los inversores. Sin embargo, el valor relevante es el que se reconoce como moneda de cambio en las transacciones de capital: la emisión de acciones o el ingreso de un socio, una fusión, la toma de control, una sociedad conjunta o una compra pagada con acciones. Son transacciones muy esporádicas, pero en ellas se manifiesta plenamente el valor creado; en este sentido, la empresa es una preparación continua para ese momento.

\begin{definicion}{Valor fundamental}
Es el valor de la empresa por su potencial competitivo y de crecimiento. Como criterio de decisión no depende de que las acciones coticen: rige igual en una sociedad cerrada, familiar o aún no constituida como sociedad anónima, porque deriva del enfoque con que se dirige la empresa. Pierde interés solo en dos casos: las empresas de subsistencia, sin atributos diferenciales y basadas en el trabajo directo del dueño, y las que, aun con ventajas competitivas, se autofinancian sin recurrir al mercado financiero (sociedades cerradas con autonomía financiera o cadenas de cooperación integradas).
\end{definicion}

Adoptar este objetivo, por lo tanto, no consiste en «acertar al número» en un momento dado, sino en una forma de concebir el desarrollo del negocio.

Conviene distinguir, además, las medidas con que se expresa el valor.

\begin{definicion}{Valor de mercado, valor contable y valor de reposición}
El \textbf{precio de mercado} es aquel al que se puede transar; el \textbf{valor de mercado}, en sentido teórico, es el valor subyacente verdadero. El \textbf{valor contable} o de libros del patrimonio suma el valor nominal de las acciones, el superávit de capital (lo aportado por encima del nominal) y las ganancias retenidas acumuladas. El \textbf{valor de reposición} es lo que costaría reemplazar los activos; se usa en la $Q$ de Tobin \vease{sec:u2-ratios}.
\end{definicion}

\begin{formulaBox}[Capitalización bursátil]
\[
\text{Valor de mercado del capital} = \text{precio de la acción}\times\text{número de acciones en circulación}
\]
\end{formulaBox}

\subsection{Formas de organización y gobierno}\label{sec:u1-agencia}

La forma jurídica condiciona la capacidad de la empresa para reunir fondos.

\begin{definicion}{Formas de organización de la empresa}
\begin{itemize}
\item \textbf{Propiedad unipersonal}: es la más barata de constituir; no paga impuesto corporativo, porque sus ganancias tributan como renta del dueño; tiene responsabilidad ilimitada y no separa el patrimonio personal del empresarial.
\item \textbf{Sociedad de personas}: barata y fácil de formar, pero los socios generales responden en forma ilimitada, la sociedad se disuelve con la muerte o el retiro de uno de ellos y es difícil transferir la propiedad. Esas tres desventajas producen una cuarta: la dificultad para reunir efectivo.
\item \textbf{Corporación} (sociedad anónima): requiere un estatuto —nombre, duración, objeto, acciones autorizadas con sus derechos por clase, derechos de los accionistas y directores iniciales— y un reglamento interno. Sus ventajas son la responsabilidad limitada al aporte, la fácil transferencia de la propiedad y la vida ilimitada, que juntas refuerzan su capacidad de captar fondos. Su desventaja es impositiva: la renta corporativa tributa y, además, los accionistas tributan por los dividendos (doble imposición).
\end{itemize}
\end{definicion}

En las sociedades cerradas suelen coincidir accionistas, directores y gerentes; en las grandes corporaciones son grupos distintos. Esa separación entre propiedad y control da permanencia a la empresa, cuyo futuro no depende de dueños ni administradores particulares. Su contracara es el problema de agencia: los accionistas delegan la maximización del valor en un consejo de administración, y este delega la gestión diaria en los gerentes, cuyos objetivos personales pueden ser otros.

\begin{definicion}{Problema y costos de agencia}
El \textbf{problema de agencia} es el conflicto entre el objetivo de maximizar el valor del dinero invertido por los accionistas y los objetivos personales de los gerentes. Los \textbf{costos de agencia} son los que genera ese conflicto: la pérdida de valor cuando los gerentes no maximizan —las \emph{pérdidas residuales}, riqueza que los accionistas pierden por el comportamiento divergente de la gerencia— y lo que se gasta en supervisarlos.
\end{definicion}

G. Donaldson observó que los gerentes se mueven por dos motivaciones básicas: la supervivencia y la independencia o autosuficiencia, es decir, decidir sin depender de terceros ni de los mercados financieros. Por eso prefieren no emitir acciones y apoyarse en el flujo generado internamente; su objetivo real sería maximizar la \emph{riqueza corporativa}, la que ellos controlan, que no coincide necesariamente con la riqueza de los accionistas.

Los accionistas disponen de cuatro mecanismos de control: eligen a los directores, que designan a la gerencia; diseñan contratos y compensaciones —como las opciones sobre acciones— que alinean los incentivos; el temor a una adquisición hostil empuja a maximizar el precio de la acción; y la competencia en el mercado laboral de gerentes disciplina a los administradores. La evidencia es coherente con la idea de que estos controles orientan a la gerencia hacia el valor para el accionista.

Tampoco basta con que la alta dirección decida todas las inversiones: tendría que analizar demasiados proyectos; el diseño de cada proyecto incluye decisiones que ella no percibe; muchas inversiones no aparecen en el presupuesto de capital; las pequeñas decisiones diarias se acumulan; y el propio director financiero enfrenta las mismas tentaciones que los niveles inferiores. Sin un sistema de incentivos, los gerentes tienden a cinco conductas que destruyen valor:
\begin{itemize}
\item \textbf{menor esfuerzo}, porque encontrar y ejecutar proyectos valiosos es exigente;
\item \textbf{beneficios adicionales} no monetarios, que compensan la falta de bonificaciones;
\item \textbf{creación de imperios}, es decir, buscar tamaño en lugar de rentabilidad;
\item \textbf{inversiones de afianzamiento}, diseñadas para requerir las habilidades del gerente actual; la tentación de sobreinvertir es mayor cuando sobra efectivo y faltan oportunidades (el \emph{problema del flujo de efectivo libre});
\item \textbf{evitar el riesgo}, porque con un salario fijo no hay motivo para encarar proyectos riesgosos aunque sean muy rentables.
\end{itemize}

Los costos de agencia se reducen supervisando, pero la supervisión cuesta tiempo y dinero y tiene rendimientos decrecientes. Es responsabilidad última de los accionistas, aunque suelen delegarla en el directorio; cuando son muchos, cada uno tiende a dejar la tarea a los demás (\emph{problema del parasitismo}). Sobre la remuneración gerencial hay dos visiones: la de contratos pactados en un mercado competitivo por el talento y la de una mala gobernanza, con directorios débiles que toleran pagos excesivos. Más que el monto importa la estructura del paquete, que debe alentar la maximización de la riqueza del accionista; como el esfuerzo es difícil de medir, la remuneración se ata al desempeño verificable.

\subsection{Función financiera y gerencia financiera}

La función financiera, en sentido genérico, no se confunde con el ámbito de competencia de la gerencia financiera. Esta es un área de la empresa con competencia exclusiva sobre actividades bien definidas: la relación con los bancos, la gestión de cobros y pagos, y el mantenimiento de una liquidez adecuada sin recursos ociosos.

Otras actividades —la definición del presupuesto, la ampliación de la planta, el control de gestión, las fusiones y adquisiciones, la información al mercado de títulos— exceden esa competencia y son responsabilidad de la dirección, como órgano político que representa a la asamblea de accionistas. Aun en ellas la gerencia financiera es relevante: en el presupuesto, por ejemplo, asesora, reúne y consolida la información, produce informes y facilita los ajustes que pide la dirección. Funciona, así, como la \emph{secretaría ejecutiva} de la función financiera. Algo similar ocurre con el área comercial: participa en las decisiones de precios o en una campaña institucional, pero quien decide es la dirección. Es la misma diferencia que existe entre estrategia y táctica, entre dirección y ejecución, entre objetivos y metas.

\subsection{Instancias de decisión}\label{sec:u1-instancias}

La visión más difundida asigna a la función financiera las decisiones de inversión, de financiamiento y de dividendos. Como los dividendos están muy ligados al financiamiento, resulta más adecuado entenderla como la interacción de las decisiones de inversión, de financiamiento y de información estratégica. Los textos suelen ordenar esas decisiones según el horizonte temporal: primero las inversiones de largo plazo y la estructura de financiamiento, después los recursos y el financiamiento corrientes (capital de trabajo) y, por último, su integración en el presupuesto anual.

Este enfoque las organiza, en cambio, en tres instancias, que son cortes transversales de las decisiones de inversión, financiamiento e información, cada una con su naturaleza y sus herramientas propias:

\begin{definicion}{Instancias de decisión financiera}
\begin{itemize}
\item \textbf{Decisiones estratégicas o estructurales}: definen la estructura de activos y de financiamiento y, con ella, la relación entre costos fijos y variables. En la inversión, esa relación es la de recursos fijos (capacidad operativa y comercial) frente a recursos variables (necesarios para usar la capacidad); en el financiamiento, la de las fuentes permanentes (deuda y capital), que también generan costos fijos y variables.
\item \textbf{Decisiones corrientes o de evolución}: modifican los recursos corrientes y el financiamiento de corto plazo según el modo de usar la estructura (plazos de venta y de pago, política de existencias). Como alteran las inversiones netas, deben decidirse con el mismo modelo financiero.
\item \textbf{Programación anual o plan del negocio}: compromisos con horizonte de un año en los que el uso de la estructura se adapta al contexto del período, se armonizan las decisiones estratégicas y corrientes y se ejercen las opciones que abrieron las primeras.
\end{itemize}
\end{definicion}

Este planteo se acerca al enfoque sistémico de Moag, Carleton y Lerner (1967). En él, el sistema básico de la empresa es el conjunto de sus transformaciones físicas, y el metasistema es un conjunto de preferencias, información y modelos que produce reglas, medidas y normas de rendimiento; esas reglas convierten las entradas físicas en salidas. La función financiera diseña el metasistema y controla el sistema básico en términos del objetivo de maximizar el valor del capital de los propietarios.

Las instancias se asocian, además, con los momentos del ciclo de vida de la organización. La estructura puede revisarse en cualquier momento, pero no de manera continua; cuando se redefine, se habla de reingeniería, reorganización o relanzamiento, es decir, de una nueva vida para la empresa.

Una forma equivalente de ordenar el campo parte del \emph{modelo de balance} de la empresa: del lado izquierdo, el valor de los activos; del derecho, el valor de la empresa para sus inversores, que muestra cómo se reparte ese valor. De ese modelo surgen las tres preguntas de las finanzas corporativas:

\begin{definicion}{Las tres preguntas de las finanzas corporativas}
\begin{enumerate}
\item \textbf{Presupuesto de capital}: ¿en qué activos de largo plazo debe invertir la empresa?
\item \textbf{Estructura de capital}: ¿cómo debe obtener el efectivo para esas inversiones?
\item \textbf{Capital de trabajo neto} (activo corriente menos pasivo corriente): ¿cómo deben administrarse los flujos operativos de corto plazo?
\end{enumerate}
\end{definicion}

%% ---------------------------------------------------------------------
\section{Cómo calcular los valores actuales}\label{sec:u1-va}

Comparar importes de distintos momentos es la operación básica de toda valuación. Las herramientas de esta sección se aplican después a los flujos de fondos \vease{sec:u1-fcff}, a los bonos y las acciones (capítulo~\ref{cap:u4}) y a los proyectos de inversión (capítulo~\ref{cap:u5}).

\subsection{Mercados financieros y VPN}

Los mercados financieros existen porque las personas desean acomodar su consumo en el tiempo, prestando y pidiendo prestado. Además, ofrecen un punto de referencia: la tasa de interés es la prueba que toda inversión propuesta debe superar, y ese estándar vale igual para individuos, empresas y gobiernos. En un gráfico que relaciona el consumo de hoy (fecha 0) con el de dentro de un año (fecha 1), la línea de mercado tiene pendiente $-(1+r)$: con $r = 10\,\%$, renunciar a \$\,1 hoy permite consumir \$\,1,10 mañana, y \$\,1 de mañana equivale a $1/1{,}10$ de hoy.

Cualquiera sea la impaciencia del inversor, la decisión correcta depende solo de comparar la inversión con las alternativas del mercado. Y los mercados no solo indican qué inversiones aceptar: también proveen los fondos para realizarlas. Sirven, en síntesis, para decidir qué hacer y cómo hacerlo.

\subsection{Valores futuros y valores presentes}

El principio más básico de las finanzas es que un peso hoy vale más que un peso mañana, porque puede invertirse y ganar intereses. El crecimiento del dinero invertido, en el que también los intereses ganan intereses, se llama \emph{interés compuesto}.

\begin{definicion}{Valor futuro, valor presente y tasa de descuento}
El \textbf{valor futuro} ($VF$) es el monto en que se convertirá un capital invertido durante $t$ períodos a la tasa $r$. El \textbf{valor presente} o actual ($VP$) es lo que hay que invertir hoy para obtener un monto futuro dado. La tasa $r$ con que se descuenta se llama \textbf{tasa de descuento} o costo de oportunidad del capital. En un mercado que funciona bien, el valor actual de un bien coincide con su precio de mercado.
\end{definicion}

\begin{formulaBox}[Valor futuro y valor presente]
\[
VF = VP\,(1+r)^{t}
\qquad\qquad
VP = \frac{C_t}{(1+r)^{t}} = C_t \cdot \underbrace{\frac{1}{(1+r)^{t}}}_{\text{factor de descuento}}
\]
\end{formulaBox}

\begin{ejemplo}[Capitalización y descuento]
Un depósito de \$\,100 al 10\,\% anual vale \$\,110 al cabo de un año, \$\,121 a los dos y \$\,133 a los tres: el crecimiento es del 10\,\% de cada saldo, no del depósito original. En diez años el factor es $(1{,}10)^{10} = 2{,}594$, y los \$\,100 llegan a \$\,259,40. A la inversa, con una tasa segura del 9\,\%, \$\,1 dentro de dos años vale hoy $1/(1{,}09)^2 = \$\,0{,}84$: prestar \$\,0,84 y renovar el préstamo devuelve exactamente \$\,1, de modo que el mercado considera equivalentes ambas sumas.
\end{ejemplo}

\begin{definicion}{Valor actual neto (VAN o VPN)}
Es el valor actual de una inversión menos lo que cuesta. Si una inversión vale más de lo que cuesta, aporta un valor neto y aumenta la riqueza de quien la realiza.
\end{definicion}

\begin{formulaBox}[Valor actual neto]
\[
VAN = VP - \text{inversión requerida}
= C_0 + \sum_{t=1}^{T}\frac{C_t}{(1+r)^{t}}
\qquad (C_0 < 0 \text{ es la inversión})
\]
\end{formulaBox}

De aquí surgen dos reglas equivalentes de decisión: aceptar las inversiones con VAN positivo, o aceptar las que rinden más que el costo de oportunidad del capital. Su discusión y la de los criterios alternativos se desarrollan en el capítulo~\ref{cap:u5} \vease{sec:u5-vpn}.

\subsection{Perpetuidades y anualidades}

Muchos flujos repiten un mismo importe a intervalos regulares, y su valor actual tiene fórmulas cerradas que derivan de la suma de una progresión geométrica.

\begin{formulaBox}[Suma de una progresión geométrica]
Si $A = a + aq + aq^2 + \dots + aq^n$, multiplicando por $q$ y restando, con primer término $A_1 = a$ y último término $A_t = aq^n$:
\[
A = \frac{A_1 - q\,A_t}{1-q}
\qquad\qquad
\text{si } |q|<1 \text{ y } n\to\infty:\quad A = \frac{A_1}{1-q}
\]
\end{formulaBox}

\begin{definicion}{Perpetuidad}
Flujo constante que se paga indefinidamente; el ejemplo clásico son los \emph{consols} británicos, cuyo tenedor cobra intereses del gobierno para siempre. Aunque la serie tiene infinitos términos, su suma es finita porque cada término es una fracción del anterior ($q = 1/(1+r)$). La intuición lo confirma: para retirar $C$ por año para siempre, hoy se necesita el capital cuyo interés anual es $C$.
\end{definicion}

\begin{definicion}{Anualidad ordinaria y anualidad anticipada}
Una \textbf{anualidad} es una serie de pagos iguales a intervalos regulares (anuales, semestrales, mensuales) durante un número determinado de períodos: jubilaciones, cuotas hipotecarias, aportes a un fondo, pagos de préstamos. En la \textbf{ordinaria} (vencida) el primer pago ocurre al final del primer período, en la fecha~1; en la \textbf{anticipada} ocurre al comienzo, en la fecha~0, y su valor es $(1+r)$ veces el de la ordinaria.
\end{definicion}

La anualidad es una perpetuidad que se interrumpe después de $N$ pagos: su valor presente es el de una perpetuidad que empieza hoy menos el de otra que empieza en $N$. La anticipada equivale a un pago inmediato más una anualidad ordinaria de $N-1$ pagos. Los valores futuros se obtienen capitalizando los valores presentes $N$ períodos.

\begin{formulaBox}[Perpetuidad y anualidades]
\[
\begin{aligned}
VP_{\text{perpetuidad}} &= \frac{C}{r} \qquad\Longleftrightarrow\qquad r = \frac{C}{VP} \\[3pt]
VP_{\text{anualidad}} &= \frac{C}{r}\left[1-\frac{1}{(1+r)^{N}}\right]
&\qquad
VF_{\text{anualidad}} &= \frac{C}{r}\left[(1+r)^{N}-1\right] \\[3pt]
VP_{\text{anticipada}} &= C + \frac{C}{r}\left[1-\frac{1}{(1+r)^{N-1}}\right]
&\qquad
VF_{\text{anticipada}} &= VP_{\text{anticipada}}\,(1+r)^{N}
\end{aligned}
\]
\end{formulaBox}

\subsection{Perpetuidades y anualidades crecientes}

Cuando el flujo crece a una tasa constante $g$, la progresión tiene razón $q = (1+g)/(1+r)$, menor que 1 si $g<r$.

\begin{formulaBox}[Perpetuidad creciente y anualidad creciente]
\[
VP = \frac{C_1}{1+r}+\frac{C_1(1+g)}{(1+r)^2}+\dots = \frac{C_1}{r-g}
\qquad\qquad
VP = \frac{C_1}{r-g}\left[1-\left(\frac{1+g}{1+r}\right)^{N}\right]
\]
\end{formulaBox}

\subsection{Cómo se pagan y cotizan los intereses}

Los intereses pueden liquidarse en distintos intervalos, de modo que para comparar alternativas hay que llevarlas a una misma medida.

\begin{definicion}{Tasa porcentual anual y tasa efectiva anual}
La \textbf{tasa porcentual anual} (TPA o tasa nominal) es la tasa que se anuncia; la \textbf{tasa efectiva anual} (TEA) es la que realmente se obtiene al capitalizar los intereses dentro del año. Una misma TPA repartida en distintos subperíodos genera TEA distintas; si la capitalización es continua, el número de subperíodos es infinito y la inversión crece en forma exponencial.
\end{definicion}

\begin{formulaBox}[Tasa efectiva y capitalización continua]
\[
TEA = \left(1+\frac{TPA}{m}\right)^{m}-1
\qquad
\xrightarrow{\;m\to\infty\;}
\qquad
TEA = e^{TPA}-1,
\qquad
VF = VP\,e^{r t},\quad e \approx 2{,}718
\]
\end{formulaBox}

\subsection{Tasas de interés e inflación}

Una tasa nominal compensa el paso del tiempo y la pérdida de poder adquisitivo; la tasa real, solo lo primero.

\begin{formulaBox}[Ecuación de Fisher]
\[
1+r_{\text{nominal}} = (1+r_{\text{real}})(1+\pi)
\qquad\Longrightarrow\qquad
r_{\text{real}} = \frac{1+r_{\text{nominal}}}{1+\pi}-1 \;\approx\; r_{\text{nominal}}-\pi
\]
\end{formulaBox}

%% ---------------------------------------------------------------------
\section{Flujo de fondos disponibles para la empresa}\label{sec:u1-fcff}

Para analizar un curso de acción o un plan de negocios (un proyecto) hay que calcular el valor que genera. Eso exige dos insumos: la proyección de su flujo de fondos y un estándar financiero, la tasa de descuento, fijado según el riesgo de esos flujos. Esta sección se ocupa del primero.

\subsection{Flujo de fondos periódicos}

El flujo de fondos proyectado es el monto neto de dinero generado durante un período (semana, mes, trimestre, año): la suma de los ingresos de fondos menos los egresos del período. No coincide con la ganancia contable, que reconoce ingresos y costos cuando se devengan y no cuando se cobran o se pagan: una venta a crédito genera ganancia, pero ningún ingreso de fondos hasta que se cobra.

\begin{ejemplo}[Ganancia contable frente a flujo de efectivo]
Una empresa vende oro en \$\,10, que le había costado \$\,9, y todavía no cobró. La contabilidad registra una ganancia de \$\,1; las finanzas registran una entrada de efectivo nula y una salida de \$\,9.
\end{ejemplo}

\subsection{¿Flujos para quién?}

Se distinguen dos medidas, según a quién quede disponible el flujo.

\begin{definicion}{Flujo de fondos disponible para la empresa (FCFF)}
El \emph{free cash flow to the firm} es el flujo operativo que genera el negocio: el que queda disponible en cada período para remunerar a quienes aportaron fondos, sea como deuda financiera (bancos y otros acreedores financieros) o como capital (accionistas, propietarios). Mide la generación del negocio al margen de cómo está financiado, por lo que excluye los flujos derivados de la decisión de financiamiento.
\end{definicion}

\begin{definicion}{Flujo de fondos disponible para el accionista (FCFE)}
El \emph{free cash flow to equity} mide cuánto queda para remunerar el capital que los propietarios mantienen inmovilizado en el negocio, después de restar no solo costos, gastos e inversiones operativas, sino también los pagos de intereses y de capital de la deuda financiera. Su relación con el FCFF se desarrolla en el capítulo~\ref{cap:u2} \vease{sec:u2-flujos}.
\end{definicion}

\subsection{Flujos operativos}

El FCFF depende sobre todo de cuánto traccionen las ventas, de los costos y gastos operativos y de las inversiones netas que el negocio requiera en activos de corto y largo plazo. Su cálculo parte de un estado de resultados proyectado sin resultados financieros y sigue cuatro pasos.

\paragraph{1. Utilidad operativa e impuesto.} La utilidad operativa resulta de restar a las ventas los costos directos, los gastos operativos (gastos de estructura u \emph{overhead}, informalmente «la mochila» de costos fijos) y las amortizaciones y depreciaciones. Estas últimas reflejan el desgaste del capital productivo e impiden considerar distribuible un resultado que no prevé mantener los activos del negocio. A ese nivel no se restan intereses, porque se busca un flujo operativo. El impuesto a las ganancias sí se resta, pero el que corresponde a la utilidad operativa, sin el ahorro que generan los intereses deducibles, que es un efecto de la financiación.

\paragraph{2. Las depreciaciones se suman de nuevo.} Las depreciaciones restan contablemente pero no son una salida de caja. Su único efecto financiero es reducir el impuesto: el \emph{ahorro fiscal} equivale a la alícuota por el cargo, siempre que el fisco admita ese método de amortización. Una vez computado ese efecto, el monto de las depreciaciones se restituye al flujo.

\paragraph{3. Se resta la inversión incremental en capital de trabajo neto.} Para vender, el negocio inmoviliza recursos en caja operativa, existencias y créditos a clientes. No es dinero perdido, sino capital que se recupera más tarde; por eso el flujo debe reflejar el momento de cada inmovilización y de su recupero (un crédito a 60 días resta cuando se otorga y suma cuando se cobra). Se mide la variación del período en los activos corrientes operativos menos la de los pasivos corrientes espontáneos (deuda comercial, que no devenga intereses). La deuda bancaria se excluye: pertenece a la decisión de financiamiento.

\paragraph{4. Se resta la inversión incremental en activos fijos.} Se resta aunque la inversión se financie íntegramente con un préstamo, porque el FCFF mide la generación de caja de la operación antes de ver cómo se financia. Si se venden activos, el recupero de capital inmovilizado aumenta el flujo. El cálculo es siempre incremental (respecto del período anterior) y en términos netos (altas menos bajas).

\begin{definicion}{Inversión incremental}
Es la inversión en activos netos realizada en cada período: el aumento de los activos corrientes operativos netos de los pasivos corrientes no financieros (capital de trabajo neto operativo) más el aumento de los activos fijos. Los aumentos reducen el flujo disponible; las disminuciones (liquidaciones, consumos, ventas) lo incrementan.
\end{definicion}

\begin{ejemplo}[Cálculo del FCFF de 2019 (alícuota del impuesto: 35\,\%)]
\begin{center}
\begin{tabular}{@{}lr@{}}
\toprule
Ventas & 15.000\\
Costo de ventas & $-9.000$\\
\midrule
Utilidad bruta & 6.000\\
Gastos operativos & $-2.000$\\
Amortizaciones y depreciaciones & $-500$\\
\midrule
Utilidad operativa (EBIT) & 3.500\\
Impuesto sobre la utilidad operativa ($3.500\times 35\,\%$) & $-1.225$\\
\midrule
Utilidad operativa neta de impuestos & 2.275\\
más amortizaciones y depreciaciones & 500\\
\midrule
Flujo antes de inversiones incrementales & 2.775\\
menos inversión incremental en capital de trabajo neto & $-55$\\
menos inversión incremental en activos fijos & $-1.000$\\
\midrule
\textbf{Flujo de fondos disponible para la empresa} & \textbf{1.720}\\
\bottomrule
\end{tabular}
\end{center}
\textbf{Ahorro fiscal de las depreciaciones.} Sin ellas, la utilidad operativa sería 4.000, el impuesto 1.400 y la utilidad neta 2.600: las depreciaciones ahorran $500\times 35\,\% = 175$ de impuesto.

\textbf{Capital de trabajo neto.} Entre diciembre de 2018 y diciembre de 2019 el activo corriente operativo pasa de 3.730 a 3.825 ($+95$): caja $+20$ (de 180 a 200), bienes de cambio $+150$ (de 1.600 a 1.750) y créditos por ventas $-75$ (de 1.950 a 1.875, quizá por el cobro de una cuenta atrasada). La deuda comercial pasa de 860 a 900 ($+40$), es decir, los proveedores financian parte del aumento de existencias. La inversión neta es $95-40 = 55$. La deuda bancaria de 700 no interviene.

\textbf{Activos fijos.} Se incorporan equipos por 1.000, financiados con un préstamo. Si además se vendieran activos fijos por 1.400 (sin efecto en resultados), la inversión neta sería $1.000-1.400 = -400$, una desinversión, y el FCFF subiría a $2.775-55+400 = 3.120$.
\end{ejemplo}

Una advertencia práctica: el balance informa los bienes de uso netos de depreciación acumulada, de modo que su variación no es la inversión del período: si pasan de 5.000 a 5.500 ($+500$) en un año con una depreciación de 500, lo invertido fue 1.000. Sin revaluaciones, vale la siguiente relación.

\begin{formulaBox}[Inversión neta en bienes de uso a partir del balance]
\[
\text{BsUso}_t = \text{BsUso}_{t-1} - \text{Dep}_t + \text{Inv}_t
\qquad\Longrightarrow\qquad
\text{Inv}_t = \text{BsUso}_t - \text{BsUso}_{t-1} + \text{Dep}_t
\]
donde $\text{BsUso}_t$ son los bienes de uso al cierre de $t$, $\text{Dep}_t$ la depreciación del período (no la acumulada) e $\text{Inv}_t$ lo invertido durante $t$.
\end{formulaBox}

\subsection{Síntesis}

\begin{formulaBox}[Flujo de fondos disponible para la empresa]
\[
\begin{aligned}
FCFF &= EBIT\,(1-t) + \text{Dep} - \Delta CTN - \text{Inv. neta en activos fijos}\\
     &= EBITDA - t\cdot EBIT - \Delta CTN - \text{Inv. neta en activos fijos}
\end{aligned}
\]
donde $EBITDA = EBIT + \text{Dep}$ y $\Delta CTN$ es la inversión incremental en capital de trabajo neto operativo.
\end{formulaBox}

El EBITDA (utilidad antes de intereses, impuestos, depreciaciones y amortizaciones) es una medida muy distinta del flujo operativo: los requerimientos de activos de corto y largo plazo suelen ser una porción muy significativa del flujo, no un dato marginal.

Proyectar y monitorear el flujo operativo es una práctica central de la gestión y de la evaluación de proyectos. Exige información sobre ventas, márgenes e inversiones, de modo que intervienen todas las áreas, no solo la gerencia financiera. Todo proyecto de envergadura —lanzar un producto, entrar en un mercado, cambiar la tecnología, adquirir equipos o discontinuar una línea— obliga a sus impulsores a traducir sus expectativas en el flujo operativo esperado. En cambio, no se les pide proyectar deuda, costo de las fuentes ni forma de financiamiento: eso corresponde a la gerencia financiera, que financia a la empresa en su conjunto y no proyecto por proyecto. Por eso el flujo operativo, que excluye resultados y activos financieros, es el más empleado para analizar proyectos. Su descuento a la tasa apropiada permite medir cuánto valor crea o destruye cada uno (capítulos~\ref{cap:u3} y~\ref{cap:u5}); con una buena estimación del flujo operativo, como suele decirse, ya está resuelto el 80\,\% del problema.

%% ---------------------------------------------------------------------
\begin{control}
\pregunta{1. ¿Cuál es el tema principal?}
La función financiera de la empresa: qué decisiones abarca, qué objetivo persigue y con qué herramientas mide el valor.

\pregunta{2. ¿Dónde y cómo se desarrolla?}
Parte del circuito de bienes y dinero de la economía, incorpora bancos, mercados y gobierno, y se traslada a la empresa como función de transformación que se decide sobre modelos. Luego aporta las herramientas de valor actual y el cálculo del flujo de fondos operativo.

\pregunta{3. ¿Cuáles son los temas secundarios?}
Rendimiento y flexibilidad de las inversiones financieras; transacciones primarias y secundarias; productividad, rentabilidad y disponibilidad de fondos; \emph{stakeholders} y valor fundamental; formas de organización y costos de agencia; instancias de decisión; perpetuidades, anualidades, TEA y Fisher; inversiones incrementales en capital de trabajo y activos fijos.

\pregunta{4. ¿Las ideas principales y secundarias están relacionadas entre sí?}
Sí. El sector financiero y los mercados amplían las oportunidades de inversión; la empresa las aprovecha con decisiones de inversión, financiamiento e información, que se evalúan por su efecto en el valor del capital de los propietarios. Ese valor se calcula descontando flujos de fondos, y el primero que se estima es el operativo.

\pregunta{5. ¿Cuál es la conclusión o la idea general?}
La empresa debe crear más valor del dinero que consume. Como solo el dinero cancela deudas y permite reasignar recursos, el desempeño se mide en flujos de fondos llevados a valor actual, y el flujo disponible para la empresa es la medida operativa de esa generación, independiente de cómo se financie.
\end{control}

%% =====================================================================
%%  UNIDAD 2
%% =====================================================================
\chapter{Cash flow y planificación financiera}\label{cap:u2}

\begin{ficha}
\fichabase{2 — Cash flow y planificación financiera}
\fila{Texto base}{G.\,L. Dumrauf, \emph{Finanzas corporativas: un enfoque latinoamericano}, 3.ª ed., Alfaomega: cap.~4 «Cash flow y planificación financiera».}
\fila{Complementaria}{G.\,L. Dumrauf, ob. cit., cap.~2 «Panorama de los estados financieros, los impuestos y el flujo de efectivo».}
\fila{}{S. Ross, R. Westerfield y J. Jaffe, \emph{Finanzas corporativas}, McGraw-Hill: cap.~2 «Estados contables y flujo de efectivo» y cap.~3 «Análisis de estados financieros y modelos financieros».}
\fila{}{Apuntes de repaso: esquema de los estados contables, índices financieros, financiamiento y crecimiento.}
\end{ficha}

%% ---------------------------------------------------------------------
\section{Panorama de los estados financieros, los impuestos y el flujo de efectivo}\label{sec:u2-eeff}

Las empresas que cotizan deben presentar cuatro estados básicos: el \textbf{balance general} (estado de situación patrimonial), el \textbf{estado de resultados}, el \textbf{estado de ganancias retenidas} (evolución del patrimonio neto), que explica los cambios en las utilidades acumuladas, y el \textbf{estado de flujos de efectivo}, que explica la variación del efectivo y sus equivalentes. Las notas y los anexos abren la composición de cada rubro.

\subsection{El balance general}

\begin{definicion}{Balance general}
Es una fotografía contable del valor de la empresa en un instante. El lado izquierdo muestra en qué invirtió la empresa (activo) y el derecho con qué recursos lo financió (pasivo y patrimonio neto). Se presenta anualmente y, en las empresas que cotizan, también cada trimestre.
\end{definicion}

\begin{formulaBox}[Identidad del balance]
\[
\text{Activo} = \text{Pasivo} + \text{Patrimonio neto}
\qquad\Longleftrightarrow\qquad
\text{Patrimonio neto} = \text{Activo} - \text{Pasivo}
\]
\end{formulaBox}

El \textbf{activo} reúne bienes y derechos ordenados por liquidez decreciente: caja y bancos, inversiones transitorias, créditos por ventas, inventarios y otros activos corrientes (créditos fiscales, gastos pagados por adelantado); luego, en el largo plazo, bienes de uso, intangibles (patentes, marcas) y créditos de realización diferida. El límite entre corriente y no corriente es el plazo de realización: dentro del año o más allá. Las empresas mantienen saldos mínimos de caja porque el dinero ocioso no rinde.

\begin{definicion}{Liquidez}
Es la facilidad y rapidez con que un activo se convierte en efectivo sin perder valor. Cuanto más líquida es una empresa, menor es el riesgo de incumplir sus obligaciones de corto plazo; pero los activos líquidos rinden menos que los fijos, de modo que sostener liquidez implica resignar rentabilidad.
\end{definicion}

El \textbf{pasivo} muestra cómo se financian las inversiones y se ordena por exigibilidad. La deuda obliga a pagar aunque no haya utilidades y su incumplimiento tiene consecuencias legales; el capital propio, en cambio, limita la responsabilidad al monto invertido. Los pasivos pueden ser:
\begin{itemize}
\item \textbf{comerciales}: la deuda espontánea con proveedores, que nace automáticamente con cada compra;
\item \textbf{fiscales y sociales}: impuesto a las ganancias, IVA y otros tributos, cargas patronales y sueldos a pagar;
\item \textbf{financieros}: deudas con bancos y con el público (obligaciones negociables); los giros y préstamos cortos financian el capital de trabajo, y los préstamos largos u obligaciones, las nuevas inversiones;
\item \textbf{otros}: juicios laborales y compromisos varios. El \emph{leasing} (alquiler con opción de compra) compromete pagos fijos y por eso funciona como un pasivo, aunque suele informarse en notas.
\end{itemize}

\begin{definicion}{Patrimonio neto (capital contable)}
Es un derecho residual, no fijo, sobre los activos: lo que queda después de atender a los acreedores, que tienen prioridad sobre los flujos de la empresa. Equivale a lo que recibirían los accionistas si todo se liquidara a valores de libros, y se entiende mejor como el saldo inicial más la ganancia neta menos los dividendos distribuidos; su evolución revela la política de dividendos.
\end{definicion}

Dentro del patrimonio se distinguen dos clases de acciones. Las \textbf{comunes} u ordinarias representan la propiedad verdadera: sus tenedores aportan el capital y asumen el riesgo del negocio; las clases (A, B, C) diferencian dividendos, votos o restricciones, y como no hay obligación legal de pagar dividendos, el accionista soporta más riesgo que el acreedor. Las \textbf{preferidas} ocupan un lugar intermedio entre las comunes y la deuda: cobran un dividendo fijo (acumulable si es acumulativo), no se benefician si las ganancias suben y, en una quiebra, cobran después de los obligacionistas y antes que los accionistas comunes. Como obligan a un desembolso fijo, generan apalancamiento y los analistas las asimilan a deuda.

Las \textbf{utilidades retenidas} se acumulan como reservas, resultados no distribuidos o dividendos en acciones. Constituir una reserva solo traslada saldos entre cuentas: es un cambio cualitativo, no cuantitativo, del patrimonio. El saldo funciona como una cuenta corriente histórica, que crece con la reinversión y cae con los dividendos en efectivo; su composición muestra cuánto capital aportaron los accionistas y cuánto proviene de reinvertir.

\begin{center}
\small
\begin{tabular}{@{}ll@{}}
\toprule
\textbf{Activo} & \textbf{Pasivo y patrimonio neto}\\
\midrule
\emph{Corriente}: efectivo y equivalentes, & \emph{Corriente}: cuentas y documentos por pagar,\\
cuentas por cobrar, inventarios & gastos acumulados\\
\emph{No corriente}: propiedad, planta y equipo & \emph{No corriente}: impuestos diferidos,\\
(neto de depreciación acumulada), & deuda de largo plazo\\
intangibles y otros & \emph{Patrimonio}: acciones preferidas y comunes,\\
 & superávit de capital, ganancias retenidas,\\
 & menos acciones en cartera (de tesorería)\\
\bottomrule
\end{tabular}
\end{center}

\begin{definicion}{Valor de libros y valor de mercado}
El valor contable registra costos históricos; el de mercado es el precio al que compradores y vendedores están dispuestos a negociar. Que coincidan sería casual. Como la meta del administrador financiero es aumentar el valor de las acciones, la referencia relevante es siempre el valor de mercado, que puede duplicar o más al registrado en libros.
\end{definicion}

\subsection{El estado de resultados y el criterio de devengado}

\begin{definicion}{Estado de resultados}
Mide el desempeño a lo largo de un período: ventas, costos, gastos, intereses e impuestos acumulados en el intervalo. Sigue el \textbf{criterio de devengado}: computa lo que corresponde al período aunque no se haya cobrado ni pagado, y el \textbf{criterio de igualación}, que aparea los ingresos con los costos y gastos necesarios para generarlos. Por eso no es un flujo de efectivo.
\end{definicion}

\begin{formulaBox}[Del ingreso a la utilidad]
\[
\begin{gathered}
\text{Ingresos} - \text{Gastos} = \text{Utilidad}
\\[4pt]
CMV = EI + \text{Compras} - EF
\\[4pt]
UPA = \frac{\text{Utilidad neta}}{\text{Acciones en circulación}}
\\[4pt]
DPA = \frac{\text{Dividendos}}{\text{Acciones en circulación}}
\end{gathered}
\]
\end{formulaBox}

Las ventas menos el costo de la mercadería vendida (CMV: materias primas y materiales, mano de obra y gastos de fabricación) dan la utilidad bruta; restando los gastos de administración y comercialización se llega al \textbf{resultado operativo} o EBIT (UAII), el renglón más importante, porque muestra lo que rinden los activos con independencia de cómo se financiaron. Después vienen:
\begin{itemize}
\item los \textbf{resultados financieros}, generados por pasivos (intereses de deudas) y por activos (intereses de colocaciones o créditos), que incluyen diferencias de cambio y de cotización, resultados por tenencia y ajuste por inflación;
\item \textbf{otros ingresos y egresos} sin encuadre específico (comisiones, alquileres ganados) y los \textbf{resultados extraordinarios} (venta de un activo fijo, indemnizaciones, siniestros);
\item el \textbf{impuesto a las ganancias} y, al final, la \textbf{utilidad neta}, que la asamblea distribuye como dividendo en efectivo o retiene (capitalizándola en acciones, formando reservas o dejándola en resultados acumulados); en todos esos casos es utilidad reinvertida.
\end{itemize}

Para leer el estado hay que considerar cuatro cuestiones:
\begin{itemize}
\item \textbf{Los principios contables.} El ingreso se reconoce cuando su proceso de generación está prácticamente concluido y hubo intercambio; una revaluación no realizada no es utilidad, lo que permite suavizar resultados vendiendo propiedades revaluadas en el momento conveniente.
\item \textbf{Las partidas que no son efectivo.} La \emph{depreciación} estima el desgaste del equipo y reduce la utilidad sin drenar caja. Los \emph{impuestos diferidos} surgen de la brecha entre utilidad contable y gravable: si hoy se tributa menos, se tributará más después, y el balance los expone como pasivo, aunque no son una salida de caja del período.
\item \textbf{El horizonte temporal.} En el corto plazo ciertos equipos y compromisos están fijos y la producción varía con más mano de obra y materias primas: son fijos, por ejemplo, los intereses de los bonos, los gastos indirectos y el impuesto inmobiliario; son variables las materias primas y los salarios de línea. En el largo plazo todos los costos son variables.
\item \textbf{La naturaleza de los costos.} La contabilidad no separa fijos de variables: los reúne como \emph{costos del producto} y llama \emph{costos del período} a los de venta, generales y de administración.
\end{itemize}

\subsection{El capital de trabajo}

\begin{definicion}{Capital de trabajo neto}
Contablemente es el activo corriente menos el pasivo corriente; una definición más depurada toma solo los activos operativos de corto plazo menos las deudas operativas de corto plazo. Invertir en él se llama \emph{cambio en el capital de trabajo neto} y suele ser positivo en las empresas que crecen.
\end{definicion}

El capital de trabajo muta en forma permanente, porque las ventas generan sin cesar cuentas a cobrar, inventarios y cuentas a pagar, que crecen solos cuando crecen las ventas. Su aumento suele leerse como buena señal, pero conviene ser cauto. Un capital de trabajo negativo puede ser ventajoso cuando la empresa cobra rápido y paga tarde: es el caso de una telefónica con inventarios mínimos y cobranzas que preceden a los pagos, o de los supermercados, que cobran las tarjetas en días y pagan a proveedores a largo plazo. Esa posición dominante les permite financiar con pasivos cortos incluso activos ajenos al capital de trabajo. En la mayoría de las empresas, sin embargo, es positivo, porque aun con el proveedor financiando toda la compra queda por cubrir el período de cobranza. Su tamaño depende del ciclo de efectivo \vease{sec:u2-ratios}, y los bancos siguen su evolución al dar crédito, porque mide a la vez desempeño y liquidez.

En una empresa ideal, que pudiera predecir con exactitud sus usos y fuentes de fondos, el capital de trabajo neto podría ser cero; en la realidad, el efectivo y el capital de trabajo funcionan como amortiguador para cumplir las obligaciones. Se mantiene efectivo por dos razones: el \emph{motivo transacción} (pagar las operaciones) y los \emph{saldos compensatorios} que exigen los bancos. Pero el exceso de efectivo tiene un costo de oportunidad —los accionistas podrían invertirlo mejor— y preocupa porque puede conducir a inversiones mal pensadas.

\subsection{Rubros que exigen depuración}

Antes de analizar, el analista debe depurar varios rubros:
\begin{itemize}
\item \textbf{Activos fijos tangibles}: figuran al costo histórico menos depreciaciones calculadas con criterios (línea recta, cuotas crecientes o decrecientes) que pueden alejarse de la depreciación económica; el terreno, además, puede revalorizarse.
\item \textbf{Bienes de cambio}: con FIFO las existencias finales quedan valuadas a los últimos precios, presumiblemente más altos; baja el CMV, sube la utilidad y se paga más impuesto. Con LIFO ocurre lo contrario, de modo que migrar a LIFO mejora el flujo por el ahorro fiscal. El techo es el valor recuperable, aunque se ven valuaciones discutibles, como computar fruta aún no cosechada.
\item \textbf{Otros créditos}: pueden esconder préstamos a los dueños, créditos impositivos abultados o créditos a empresas del grupo que se cancelarían al consolidar.
\item \textbf{Inversiones transitorias}: plazos fijos, títulos y fondos comunes, cuyo valor de libros puede diferir mucho del de mercado (por ejemplo, ante un \emph{default} soberano).
\item \textbf{Créditos por ventas}: el crédito comercial impulsa las ventas pero puede ocultar incobrables; conviene revisar el período medio de cobranza y la morosidad.
\end{itemize}

\paragraph{Métodos de depreciación.} Los activos fijos pierden valor por uso, tiempo u obsolescencia, y ese costo debe imputarse a resultados repartiéndolo a lo largo de la vida útil, aunque en sectores de alta tecnología la brecha entre depreciación económica y contable puede ser enorme. La depreciación acelerada castiga los primeros ejercicios y ahorra impuestos antes; aunque después se tribute más, siempre conviene pagar más tarde. El método elegido afecta a la vez el valor de los bienes de uso y el cargo a resultados.

\paragraph{Investigación y desarrollo.} Investigación, desarrollo y publicidad se imputan al período, pero se parecen más a inversiones, porque buscan ventas futuras; deberían capitalizarse, no porque siempre creen valor, sino porque se espera que lo hagan. Como los inversores pagan lo que los activos valen, las empresas con mucho gasto en I+D muestran precios de mercado muy alejados de sus valores de libros.

\paragraph{Inflación y devaluación.} La inflación distorsiona ambos estados, sobre todo los activos fijos. A veces favorece: un pasivo en moneda local sin ajuste se licúa. Otras perjudica: vender a 60 días con una inflación del 10\,\% en ese lapso permite recuperar cerca del 90\,\% de lo vendido. Con inflación anual del 20\,\%, ventas y gastos mezclan precios de comienzo y de fin de año, y el EBIT del cierre no representa el actual: el analista debe actuar como un detective. La devaluación, concentrada en lapsos breves, transfiere ingresos con fuerza: beneficia a un exportador endeudado en moneda local y perjudica a un importador o a quien se endeudó en moneda extranjera, generando un resultado por exposición que alcanza activos, pasivos y resultados.

\subsection{Impuestos y planeamiento impositivo}

\begin{definicion}{Tasa promedio y tasa marginal del impuesto}
La \textbf{tasa promedio} es el impuesto dividido por el ingreso gravable, es decir, la porción de la renta que va al fisco. La \textbf{tasa marginal} es el impuesto que se pagaría por un peso adicional de ganancia. Con una alícuota fija ambas coinciden; en los sistemas progresivos, con escalas, la marginal supera a la promedio y la promedio sube con el ingreso. Para las decisiones importa la marginal; los sistemas corporativos tienden a una tasa fija para las rentas altas.
\end{definicion}

El \emph{tax planning} usa figuras legales para reducir el impacto impositivo y crear valor para el accionista; es lícito y la propia ley lo promueve con exenciones y desgravaciones. La estrategia más difundida es endeudarse, porque los intereses son deducibles: una deuda de \$\,50 al 10\,\%, con alícuota del 40\,\%, reduce el impuesto en \$\,2. Otra vía es aprovechar los escudos fiscales dentro del ejercicio: cuando las existencias cerrarían valuadas a precios muy altos —como el cereal de una exportadora tras una devaluación—, canjear parte de ellas antes del cierre por bienes de uso amortizables reduce el impuesto.

\begin{ejemplo}[Planeamiento impositivo en una cerealera exportadora]
El trigo, que sigue al dólar, pasa de \$\,104,60 por tonelada en diciembre de 2001 a \$\,400 en junio de 2002 y a \$\,540 en diciembre. Sin planificar, el impuesto llega a \$\,76.195, porque el cereal en existencia cierra valuado a precios muy altos. Si en junio se canjea cereal por maquinaria por \$\,200.000, amortizable en línea recta (escudo de \$\,20.000 en el ejercicio), el impuesto baja a \$\,44.695; si el canje se difiere a diciembre, por \$\,270.000, el impuesto queda entre ambos valores.
\end{ejemplo}

Cuatro reglas prácticas resumen el planeamiento: por efecto de la inflación conviene reconocer beneficios al inicio del ejercicio y pérdidas al final; conviene cerrar el ejercicio con activos de baja valuación impositiva; el endeudamiento reduce la imposición total; y la fecha de cierre debe elegirse según el negocio.

También son materia de planeamiento los \textbf{quebrantos}, deducibles solo en los cinco períodos fiscales siguientes, y la imputación de ciertos desembolsos: tratar una reparación como gasto castiga el primer ejercicio (y en un negocio nuevo puede desperdiciarse por falta de ganancias), mientras que activarla la amortiza en varios años. Las diferencias entre balance contable y fiscal provienen de la amortización acelerada, los criterios de valuación, las exenciones, los límites a la deducción de gastos y el arrastre de quebrantos; el \textbf{impuesto diferido} reconoce el efecto de esas diferencias temporarias imputándolo al ejercicio en que se devenga. Los impuestos sobre activos o patrimonios suelen aplicar tasas proporcionales.

\subsection{Flujo de efectivo financiero}\label{sec:u2-feA}

El dato más valioso que puede extraerse de los estados contables es el flujo de efectivo, porque en finanzas el valor de una empresa es su capacidad de generarlo; no debe confundirse con el capital de trabajo neto. La identidad básica dice que el efectivo que producen los activos se reparte entre acreedores y accionistas.

\begin{formulaBox}[Identidad del flujo de efectivo]
\[
FE(A) = FE(B) + FE(S)
\]
\begin{center}
\small
\begin{tabular}{@{}l@{\qquad}l@{}}
\toprule
\textbf{Flujo de los activos, $FE(A)$} & \textbf{Flujo a los inversores}\\
\midrule
Flujo operativo: $EBIT + \text{Dep} - \text{impuestos}$ & $FE(B)$, acreedores: intereses $+$ amortización \\
$-$ gastos de capital (altas $-$ bajas de activos fijos) & \quad de deuda $-$ nueva deuda\\
$-$ aumento del capital de trabajo neto & $FE(S)$, accionistas: dividendos $+$ recompras\\
 & \quad $-$ nuevas emisiones\\
\bottomrule
\end{tabular}
\end{center}
\end{formulaBox}

El \textbf{flujo operativo} refleja el pago de impuestos pero excluye el financiamiento, los gastos de capital y la inversión en capital de trabajo. Los \textbf{gastos de capital} son las adquisiciones menos las ventas de activos fijos y se obtienen del balance como la variación de los activos fijos netos más la depreciación, la misma relación vista en el capítulo~\ref{cap:u1} \vease{sec:u1-fcff}. A los acreedores se les paga el servicio de la deuda (intereses y principal) neto de los nuevos préstamos; a los accionistas, los dividendos y recompras netos de nuevas emisiones.

\subsection{El estado contable de flujos de efectivo}

El estado oficial explica la variación del efectivo con tres bloques:
\begin{enumerate}
\item \textbf{actividades de operación}: parte de la utilidad neta, reincorpora los gastos que no son efectivo (depreciación, impuestos diferidos) y ajusta los cambios en activos y pasivos corrientes, salvo el efectivo y los documentos por pagar, que van a financiamiento;
\item \textbf{actividades de inversión}: compras y ventas de activos fijos;
\item \textbf{actividades de financiamiento}: pagos netos a acreedores y propietarios, excluidos los intereses.
\end{enumerate}

La suma de los tres reproduce la variación de caja del balance. Pero hay un desfase conceptual con el flujo financiero: los intereses son un pago a los acreedores y deberían figurar en financiamiento, mientras que la contabilidad los deduce al calcular la utilidad neta. Esa es la diferencia de fondo entre el flujo contable y el flujo financiero.

Del análisis surgen cuatro observaciones. El flujo operativo mide lo que generan las operaciones sin considerar gastos de capital ni capital de trabajo. El flujo total sí los incorpora y con frecuencia es negativo: cuando la empresa crece rápido, lo que invierte en inventarios y activos fijos supera lo que genera. La utilidad neta no es un flujo de efectivo, y para juzgar la situación de la empresa el flujo es más revelador. Por último, el \emph{free cash flow} es el efectivo que la empresa puede repartir libremente entre acreedores y accionistas porque no está comprometido en capital de trabajo ni en activos fijos.

El flujo se popularizó como indicador porque deja menos margen a la subjetividad que los criterios contables. Aun así, varias empresas inflaron su flujo operativo reclasificando partidas: imputaron como operativos flujos de inversión (o a la inversa), registraron compras de cuentas como financiamiento y sus cobros como operación, armaron operaciones de ida y vuelta o capitalizaron mano de obra como activo fijo (casos Avis Budget, Vail Resorts, Tyco, Dynegy y Adelphia). Ninguna de esas maniobras altera el flujo \emph{total} de la firma; por eso conviene vigilar esa cifra y no solo el flujo operativo.

%% ---------------------------------------------------------------------
\section{Análisis de estados financieros y modelos financieros}\label{sec:u2-analisis}

\subsection{Estandarización de los estados financieros}

Para comparar empresas de distinto tamaño se trabaja con porcentajes en lugar de moneda. En los \textbf{estados porcentuales}, cada rubro del balance se expresa como porcentaje del activo total y cada renglón del estado de resultados como porcentaje de las ventas, lo que muestra qué ocurre con cada peso vendido. La otra vía para neutralizar el tamaño son las razones financieras.

\begin{definicion}{Métodos de análisis de estados financieros}
\begin{itemize}
\item \textbf{Análisis vertical}: sobre estados porcentuales de un mismo período.
\item \textbf{Análisis horizontal o de tendencia}: compara la evolución temporal de cada rubro o razón. Es administración por excepción: una razón corriente de 2,4 que cae de modo sostenido durante diez años puede indicar deterioro de la liquidez, pero también un uso más eficiente del activo corriente o un cambio en el negocio; no es necesariamente mala, pero debe investigarse.
\item \textbf{Análisis de grupo de pares}: compara con empresas que compiten en los mismos mercados y operan de modo semejante (promedio de la industria, las diez mejores o peores). Como no hay dos empresas idénticas, la elección del grupo es en última instancia subjetiva.
\end{itemize}
\end{definicion}

\subsection{Razones financieras}\label{sec:u2-ratios}

Las razones eliminan el problema del tamaño y expresan la información como porcentajes, múltiplos o plazos. Se agrupan en cinco familias; las de valor de mercado solo pueden calcularse directamente en empresas que cotizan.

\paragraph{Liquidez o solvencia de corto plazo.} Miden la capacidad de pagar las cuentas inmediatas sin presiones. Se basan en activos y pasivos corrientes, que cambian con rapidez y por eso no siempre anticipan bien el futuro.

\begin{formulaBox}[Razones de liquidez]
\[
\begin{gathered}
\text{Razón corriente} = \frac{\text{Activo corriente}}{\text{Pasivo corriente}}
\\[4pt]
\text{Prueba ácida} = \frac{\text{Activo corriente}-\text{Inventario}}{\text{Pasivo corriente}}
\\[4pt]
\text{Razón de efectivo} = \frac{\text{Efectivo}}{\text{Pasivo corriente}}
\end{gathered}
\]
\end{formulaBox}

Para el acreedor de corto plazo, cuanto más alta la razón corriente, mejor; para la empresa, una razón alta indica liquidez pero puede delatar un uso ineficiente del efectivo. Se espera al menos 1, porque por debajo el capital de trabajo neto es negativo, aunque una razón baja no es grave si la firma conserva capacidad de endeudamiento. La prueba ácida excluye el inventario, el activo corriente menos líquido, y la razón de efectivo interesa al acreedor de plazo muy breve.

\paragraph{Solvencia de largo plazo o apalancamiento.} Evalúan la capacidad de atender las obligaciones de largo plazo.

\begin{formulaBox}[Razones de endeudamiento y cobertura]
\[
\begin{gathered}
\text{Razón de deuda total} = \frac{\text{Activo total}-\text{Patrimonio}}{\text{Activo total}}
\\[4pt]
\text{Razón deuda/capital} = \frac{\text{Deuda}}{\text{Patrimonio}}
\\[4pt]
\text{Multiplicador del capital} = \frac{\text{Activo total}}{\text{Patrimonio}} = 1+\frac{\text{Deuda}}{\text{Patrimonio}}
\\[4pt]
\text{Cobertura de intereses (TIE)} = \frac{EBIT}{\text{Intereses}}
\\[4pt]
\text{Cobertura de efectivo} = \frac{EBIT+\text{Dep. y amort.}}{\text{Intereses}}
\\[4pt]
\frac{\text{Deuda que devenga intereses}}{EBITDA}
\end{gathered}
\]
\end{formulaBox}

La cobertura de intereses usa el EBIT, que no es efectivo disponible porque ya dedujo depreciaciones; de allí la cobertura de efectivo. La relación entre la deuda que devenga intereses (documentos por pagar y deuda de largo plazo) y el EBITDA es cada vez más usada en los contratos de crédito: valores inferiores a 1 se consideran muy sólidos y los superiores a 5, débiles.

\paragraph{Administración o rotación de activos.} Describen con qué eficiencia e intensidad se usan los activos para vender.

\begin{formulaBox}[Razones de rotación y plazos]
\[
\begin{gathered}
\text{Rotación de inventario} = \frac{CMV}{\text{Inventario}}
\\[4pt]
\text{Días de inventario} = \frac{365}{\text{Rotación de inventario}}
\\[4pt]
\text{Rotación de créditos} = \frac{\text{Ventas}}{\text{Cuentas por cobrar}}
\\[4pt]
\text{Período medio de cobranza} = \frac{365}{\text{Rotación de créditos}}
\\[4pt]
\text{Rotación del activo total} = \frac{\text{Ventas}}{\text{Activo total}}
\end{gathered}
\]
\end{formulaBox}

Mientras la empresa no se quede sin existencias ni pierda ventas, una rotación de inventario más alta indica mejor administración; la rotación del activo total da la visión global. El capital de trabajo ocupa gran parte del tiempo del gerente financiero, y su manejo admite tres políticas:

\begin{definicion}{Políticas de capital de trabajo}
\begin{itemize}
\item \textbf{Conservadora}: altos saldos de efectivo e inventarios financiados con deuda de largo plazo; reduce el riesgo de quedarse sin stock o sin liquidez, pero baja la rentabilidad.
\item \textbf{Agresiva}: saldos mínimos financiados con deuda bancaria de corto plazo; mejora la rentabilidad, pero expone a la empresa a la tasa de interés y al desabastecimiento.
\item \textbf{Moderada o de sincronización}: calza el plazo del financiamiento con la duración de cada activo; la porción permanente del capital de trabajo se financia a largo plazo y la estacional, a corto.
\end{itemize}
Un aumento de las cuentas por cobrar no es una política «conservadora»: suele reflejar una relajación del crédito para vender más, que aumenta el riesgo.
\end{definicion}

\begin{definicion}{Ciclo operativo y ciclo de efectivo}
El \textbf{ciclo operativo} va desde la llegada de las existencias hasta el cobro: días de inventario más días de cobranza. El \textbf{ciclo de efectivo} o ciclo dinero–mercancía–dinero va desde el pago de los materiales hasta el cobro: el ciclo operativo menos los días de financiamiento de los proveedores.
\end{definicion}

Cuanto más corto el ciclo de efectivo, menor el capital inmovilizado y, en igualdad de condiciones, mayor la rentabilidad. La elección de la política de capital de trabajo depende también de cómo afecte cada escenario (recesión, normalidad o expansión) a las ventas y al costo del financiamiento de corto plazo.

\begin{ejemplo}[Ciclos]
Con 45 días de inventario y 45 de cobranza, el ciclo operativo es de 90 días; si los proveedores financian la compra a 30 días, el ciclo de efectivo se reduce a 60.
\end{ejemplo}

\paragraph{Rentabilidad.} Son las razones más difundidas: miden la eficacia en el uso de los activos y la eficiencia de las operaciones.

\begin{formulaBox}[Razones de rentabilidad y de distribución]
\[
\begin{gathered}
\text{Margen neto} = \frac{\text{Utilidad neta}}{\text{Ventas}}
\\[4pt]
\text{Margen EBITDA} = \frac{EBITDA}{\text{Ventas}}
\\[4pt]
ROA = \frac{\text{Utilidad neta}}{\text{Activo total}}
\\[4pt]
ROE = \frac{\text{Utilidad neta}}{\text{Patrimonio}}
\\[4pt]
\text{Índice de pago } (d) = \frac{\text{Dividendos}}{\text{Utilidad neta}}
\\[4pt]
\text{Índice de retención } (b) = \frac{\text{Utilidades retenidas}}{\text{Utilidad neta}} = 1-d
\end{gathered}
\]
\end{formulaBox}

El margen EBITDA se acerca más al flujo operativo, porque suma las partidas que no son efectivo y queda libre del efecto de impuestos e intereses, es decir, de la estructura de capital. El ROE mide en términos contables cómo les fue a los accionistas, y el ROA, la utilidad por unidad de activo; también suele calcularse con EBIT (ROA bruto) y sobre saldos promedio. Estas medidas tienen tres limitaciones: las ganancias actuales pueden reflejar mal la rentabilidad futura, ignoran el riesgo y carecen de un punto de referencia propio.

\paragraph{Valor de mercado.} Requieren un dato ajeno a los estados contables: el precio de la acción.

\begin{formulaBox}[Razones de valor de mercado]
\[
\begin{gathered}
\frac{P}{U} = \frac{\text{Precio por acción}}{UPA}
\\[4pt]
\text{Rendimiento por dividendos} = \frac{DPA}{\text{Precio}}
\\[4pt]
\frac{\text{Valor de mercado}}{\text{Valor de libros}} = \frac{\text{Precio por acción}}{\text{Patrimonio}/\text{Acciones en circulación}}
\\[4pt]
VE = \text{Capitalización bursátil} + \text{Deuda que devenga intereses (a mercado)} - \text{Efectivo}
\\[4pt]
Q \text{ de Tobin} = \frac{\text{Valor de mercado de deuda y capital}}{\text{Valor de reposición de los activos}}
\end{gathered}
\]
\end{formulaBox}

\begin{ejemplo}[Razón precio-utilidad]
Si una acción cotiza a \$\,24 y la utilidad por acción de los últimos doce meses es \$\,3, la razón P/U es 8: el comprador paga ocho pesos por cada peso de utilidad corriente.
\end{ejemplo}

La razón P/U suele leerse como expectativa de crecimiento: las acciones con crecimiento más alto o más cierto muestran P/U elevadas, y las de crecimiento menor o más riesgoso, bajas. Con pérdidas queda indefinida (aunque a veces se exhiba negativa) y con utilidades mínimas se dispara. Si dos acciones fueran iguales en todo, sobre todo en crecimiento, la de P/U doble sería menos atractiva; como rara vez lo son, la comparación entre industrias o períodos engaña y la razón sirve sobre todo dentro de un mismo sector. Sus determinantes se analizan en el capítulo~\ref{cap:u4}.

El rendimiento por dividendos es más bajo en las empresas con altas perspectivas de crecimiento. La razón valor de mercado a valor de libros compara el valor de las inversiones con su costo histórico: por debajo de 1 sugiere que la firma no creó valor. El \textbf{valor de la empresa} (VE) estima cuánto costaría comprar todas las acciones y cancelar la deuda; se resta el efectivo porque el comprador lo usaría de inmediato para recomprar deuda o pagar dividendos, y sus múltiplos se usan para valuar el negocio completo. Una $Q$ de Tobin mayor que 1 indica un incentivo a invertir y suele asociarse con oportunidades atractivas o con una ventaja competitiva significativa.

\subsection{La identidad de Du Pont}

El ROE puede descomponerse multiplicando y dividiendo por el activo y por las ventas:

\begin{formulaBox}[Identidad de Du Pont]
\[
ROE = \frac{\text{Utilidad neta}}{\text{Patrimonio}}
= \underbrace{\frac{\text{Utilidad neta}}{\text{Ventas}}}_{\text{margen}}
\times \underbrace{\frac{\text{Ventas}}{\text{Activo}}}_{\text{rotación}}
\times \underbrace{\frac{\text{Activo}}{\text{Patrimonio}}}_{\text{multiplicador}}
\]
\[
ROA = \text{margen}\times\text{rotación}
\]
\end{formulaBox}

El rendimiento del capital depende así de tres palancas: la eficiencia operativa (margen), la eficiencia en el uso de los activos (rotación) y el apalancamiento financiero (multiplicador del capital). Si el ROE es insatisfactorio, la identidad indica dónde buscar la causa.

%% ---------------------------------------------------------------------
\section{Cash flow y planificación financiera}\label{sec:u2-cashflow}

\subsection{Usos del cash flow}

El análisis del flujo de efectivo y su proyección tienen tres usuarios: los gerentes lo emplean como indicador de la gestión financiera de corto y mediano plazo; los consultores, para estimar el valor de una empresa o el VAN de un proyecto; y los analistas bancarios, para juzgar la capacidad de repago de una deuda.

\subsection{El modelo indirecto}

\begin{definicion}{Cash flow indirecto}
Integra el balance, el estado de resultados y el de evolución del patrimonio neto. Son \textbf{ingresos de fondos} las disminuciones de activos, los aumentos de pasivos, los resultados positivos y los aportes de los socios; son \textbf{egresos} los aumentos de activos, las disminuciones de pasivos, los resultados negativos y los dividendos en efectivo. La caja es el resultado del modelo, por lo que su propio aumento no se computa como egreso.
\end{definicion}

Ninguna operación con efecto en caja queda sin reflejar: una venta a crédito aumenta las cuentas a cobrar (egreso) y el EBIT (ingreso), con efecto neto nulo hasta que se cobra; lo mismo sucede con las compras a crédito o con un resultado por exposición a la inflación. El modelo tiene cuatro bloques:
\begin{enumerate}
\item \textbf{Operaciones}: del EBIT se pasa al EBITDA sumando depreciaciones y amortizaciones; se ajustan los cambios en el capital de trabajo (los aumentos de créditos e inventarios restan; los de deudas comerciales, fiscales y sociales suman) y se resta el impuesto a las ganancias.
\item \textbf{Inversiones}: se restan las altas de bienes de uso y se suman las bajas; conviene separar las inversiones permanentes en otras empresas.
\item \textbf{Financiamiento}: se suman los aumentos de deuda financiera y los aportes de capital; se restan las cancelaciones y los intereses pagados, y se suman los intereses cobrados.
\item \textbf{Accionista y cash flow neto}: el flujo del accionista se destina a dividendos o a inversiones transitorias, y lo que resta, sumado a la caja inicial, da la caja final.
\end{enumerate}

Un formato equivalente es el \textbf{estado de origen y aplicación de fondos}. Son orígenes las disminuciones de activos (salvo caja y bienes de uso), las bajas brutas de activos fijos, los aumentos de pasivos, las emisiones de acciones y los fondos de las operaciones; son aplicaciones los movimientos inversos, las recompras y los dividendos. La depreciación no es fuente de fondos: con pérdidas operativas no hay fondos, por grande que sea. Ubicar los dividendos frente a la utilidad y las altas de activos fijos frente a la depreciación facilita la lectura. Estos estados alertan sobre problemas (inventarios que crecen más que las ventas), muestran qué parte del crecimiento se financió interna o externamente y si la empresa creció demasiado rápido.

\subsection{Free Cash Flow, cash flow del accionista y Capital Cash Flow}\label{sec:u2-flujos}

El \textbf{free cash flow} (FCF) es el flujo de fondos disponible para la empresa definido en el capítulo~\ref{cap:u1} \vease{sec:u1-fcff}: el que genera la firma con independencia de cómo se financia, como si lo hiciera solo con capital propio. Es la medida más usada para valuar empresas y proyectos y también como indicador de desempeño. Para aislar el financiamiento se calcula el impuesto que habrían pagado las operaciones ($t\times EBIT$), se excluyen los intereses y no se computan los cambios de deuda ni las cuotas de \emph{leasing}. A partir de él se definen los otros dos flujos.

\begin{definicion}{Cash flow del accionista y Capital Cash Flow}
El \textbf{cash flow del accionista} (ECF) incorpora todos los efectos financieros: resta los intereses pagados, suma los cobrados, suma o resta los cambios de deuda y agrega el ahorro fiscal de los intereses, que compensa el impuesto «hipotético» más alto usado en el FCF. Es un residuo: el accionista cobra después de proveedores, empleados, bancos y Estado. El \textbf{Capital Cash Flow} (CCF) es lo que la empresa destina a todos sus inversores, accionistas y acreedores.
\end{definicion}

\begin{formulaBox}[Relación entre los flujos]
\[
\begin{gathered}
ECF = FCF - \text{Intereses}\,(1-t) \pm \Delta\,\text{Deuda financiera}
\qquad\qquad
\text{Ahorro fiscal} = t\times\text{Intereses}
\\[4pt]
CCF = ECF + \text{Flujo de los acreedores} = FCF + t\times\text{Intereses}
\end{gathered}
\]
El flujo de los activos $FE(A)$ de la identidad contable \vease{sec:u2-feA} usa los impuestos efectivamente pagados, que ya reflejan la deducción de intereses: coincide, por lo tanto, con el CCF.
\end{formulaBox}

\begin{figure}[H]
  \centering
  \includegraphics[width=0.5\linewidth,height=0.42\textheight,keepaspectratio]{Finanzas20corporativas203ed20Dumrauf-2}
  \caption{Free Cash Flow, cash flow del accionista y Capital Cash Flow.}
\end{figure}

\paragraph{Interpretación del cash flow histórico.} El cash flow integra las variaciones patrimoniales con los resultados. Si las ventas crecen al ritmo del PIB, créditos, inventarios y deudas comerciales crecen a un ritmo similar; en una recesión se contraen. Como un solo año produce un «efecto foto», conviene analizar el flujo acumulado de varios.

\begin{ejemplo}[Cash flow acumulado de cinco años]
Una empresa generó en cinco años un EBITDA de \$\,20.456 millones; pagó impuestos e invirtió unos \$\,9.500 millones en activos fijos, y obtuvo un FCF de \$\,8.401 millones, algo más del 40\,\% del EBITDA. El financiamiento le restó \$\,4.277 millones, y el flujo del accionista, \$\,4.124 millones, se destinó en partes similares a dividendos y a inversiones transitorias.
\end{ejemplo}

Suele decirse que el verdadero flujo del accionista son los dividendos. Para la teoría de la valuación, sin embargo, lo no distribuido son dividendos potenciales invertidos en títulos de VAN esperado nulo, que conservan su valor: es solo una diferencia de tiempo. Por eso dividendos e inversiones transitorias se exponen como destinos del flujo del accionista, y valuar descontando ese flujo no difiere en esencia de descontar dividendos.

\subsection{Planificación financiera}

\begin{definicion}{Planificación de corto y de largo plazo}
La \textbf{planificación de corto plazo} abarca menos de un año, con proyecciones mensuales o diarias: es el terreno de la tesorería, donde se prevé cuándo pagar sueldos, impuestos y deudas y se colocan los excedentes en inversiones seguras y líquidas. Su herramienta es el \textbf{presupuesto de efectivo}, que proyecta ingresos y egresos por período: mensual si hay estacionalidad, más frecuente si los flujos son volátiles, y tan bueno como sus supuestos. La \textbf{planificación de largo plazo} abarca cinco años o más, combina todos los proyectos para determinar la inversión total y se arma por años, con menos detalle y escenarios alternativos; sirve para el plan de negocios, los proyectos de inversión y la valuación.
\end{definicion}

Una proyección aislada del flujo no es consistente: debe integrar estado de resultados, balance, evolución del patrimonio y cash flow. Para lograrlo conviene seguir siete reglas:
\begin{enumerate}
\item usar premisas consistentes y defendibles ante terceros, basadas en datos y no en opiniones (si crecen las ventas deben crecer los activos fijos; si se supone una devaluación, deben ajustarse precios y salarios);
\item analizar el desempeño histórico: evolución de las ventas, relación EBIT/ventas, ciclo de efectivo y relación entre ventas y activos fijos;
\item «exprimir» los datos con controles cruzados;
\item formular las preguntas pertinentes: qué activos exige el crecimiento, cuándo reponerlos y cómo financiarlos;
\item incluir las variables macroeconómicas (inflación, tipo de cambio, tasa de interés, nivel de actividad) con escenarios ponderados;
\item considerar hechos futuros, como juicios o el vencimiento de patentes;
\item ejercitar el pesimismo, porque en las proyecciones el optimismo es la regla.
\end{enumerate}
Proyectar obliga a pensar el futuro, mejora el control de inventarios y créditos, alerta cuando no se alcanzan las metas y muestra al banco las necesidades de capital de trabajo y la capacidad de pago.

\subsection{Proyección del estado de resultados}

Es siempre el primer paso, porque define las ventas y el costo de ventas, que luego determinan créditos, inventarios, deudas comerciales y los activos fijos necesarios.

\paragraph{Proyección de las ventas.} Las ventas de la mayoría de las empresas siguen el ciclo económico: el PIB crece en el largo plazo (en Sudamérica, cerca del 3\,\% anual) y la inflación aumenta las ventas nominales. Una forma eficaz de comunicar esta hipótesis es comparar la evolución de una unidad de ventas con una unidad de PIB a precios corrientes, partiendo ambas de 1 y multiplicando cada año por el cociente entre el valor del año y el del anterior.

\begin{figure}[H]
  \centering
  \includegraphics[width=0.5\linewidth,height=0.42\textheight,keepaspectratio]{Finanzas20corporativas203ed20Dumrauf-3}
\end{figure}
\begin{figure}[H]
  \centering
  \includegraphics[width=0.5\linewidth,height=0.42\textheight,keepaspectratio]{Finanzas20corporativas203ed20Dumrauf-5}
  \caption{Ventas y PIB a precios corrientes: evolución de una unidad de cada uno en empresas cotizantes y relación entre ambos en una empresa de telecomunicaciones.}
\end{figure}

Dos métodos estadísticos apoyan la proyección.

\begin{definicion}{Regresión lineal de las ventas}
Establece una relación causal entre las ventas (variable dependiente) y el PIB a precios corrientes (variable independiente). La pendiente es significativa si su estadístico $t$ (pendiente sobre su desvío estándar) tiene un valor $p$ menor que el 5\,\%, porque entonces se rechaza la hipótesis de que no hay relación. La ecuación permite predecir el valor medio de las ventas, con una \emph{banda de confianza}, o un valor particular, con un \emph{intervalo de pronóstico} más amplio porque suma la incertidumbre del término de error. Ambos se ensanchan a medida que el PIB se aleja de su media.
\end{definicion}

\begin{formulaBox}[Proyección por regresión]
\[
\text{PIB}_{t+1} = \text{PIB}_{t}\,(1+g_{t+1})(1+\pi_{t+1})
\qquad\qquad
\text{Ventas} = a + b\cdot\text{PIB}
\]
\[
s_{\hat y} = s_e\sqrt{\frac{1}{n}+\frac{(x_0-\bar x)^2}{\sum (x_i-\bar x)^2}}
\qquad\qquad
s_{\text{pronóstico}} = s_e\sqrt{1+\frac{1}{n}+\frac{(x_0-\bar x)^2}{\sum (x_i-\bar x)^2}}
\]
donde $g$ es el crecimiento real del PIB, $\pi$ la inflación implícita, $s_e$ el error estándar de la regresión y $x_0$ el PIB proyectado.
\end{formulaBox}

\begin{figure}[H]
  \centering
  \includegraphics[width=0.5\linewidth,height=0.42\textheight,keepaspectratio]{Finanzas20corporativas203ed20Dumrauf-12}
  \caption{Bandas de confianza e intervalos de pronóstico para las ventas.}
\end{figure}

\begin{definicion}{Series de tiempo}
Abandonan la causalidad y usan los valores pasados de la misma variable. La serie se descompone en estacionalidad, aleatoriedad, tendencia y ciclo ($V = E\times A\times T\times C$). El promedio móvil de cuatro trimestres elimina estacionalidad y aleatoriedad y deja $T\times C$; la tendencia se estima por mínimos cuadrados; el ciclo es el promedio móvil dividido por la tendencia; y los coeficientes estacionales surgen de promediar $V/(T\times C)$ para cada trimestre, ajustados para que sumen 4. El pronóstico multiplica la tendencia por el factor estacional y por el último ciclo conocido; la aleatoriedad no se proyecta.
\end{definicion}

\begin{figure}[H]
  \centering
  \includegraphics[width=0.5\linewidth,height=0.42\textheight,keepaspectratio]{Finanzas20corporativas203ed20Dumrauf-16}
\end{figure}
\begin{figure}[H]
  \centering
  \includegraphics[width=0.5\linewidth,height=0.42\textheight,keepaspectratio]{Finanzas20corporativas203ed20Dumrauf-19}
  \caption{Ventas trimestrales: promedio móvil y tendencia lineal.}
\end{figure}

Hay otras alternativas: la regresión múltiple (para un exportador, sobre el tipo de cambio y el ingreso del país vecino), los modelos autorregresivos de corto plazo y, para productos sin historia, las encuestas, con sus problemas de muestra, costo y sesgo. El pronóstico de ventas es el punto más importante de toda la proyección, porque determina resultados, impuestos, cash flow e inversiones.

\paragraph{Costos, resultados financieros e impuestos.} Los costos operativos (costo de ventas, gastos de comercialización y de administración) se proyectan como porcentaje de las ventas, con el promedio histórico o con la tendencia reciente si el promedio está afectado por valores extremos, o bien con regresiones contra las ventas, que dan resultados similares.

\begin{figure}[H]
  \centering
  \includegraphics[width=0.5\linewidth,height=0.42\textheight,keepaspectratio]{Finanzas20corporativas203ed20Dumrauf-20}
  \caption{Regresiones lineales entre las ventas y los costos operativos.}
\end{figure}

Los resultados financieros (intereses, diferencias de cambio y de cotización, exposición a la inflación, resultados por tenencia) son más difíciles de pronosticar, porque dependen de precios relativos que la dirección no controla. Los intereses pueden estimarse con las tasas y los sistemas de amortización pactados; el resto es casi imposible de predecir, y una devaluación aguda puede dominar el resultado financiero de dos décadas. Además, proyectar una diferencia de cambio negativa supone poder predecir el tipo de cambio, en cuyo caso convendría cubrirse. Por eso la valuación trabaja con el FCF, que excluye estos resultados y considera la tasa de la deuda solo en el costo de capital.

\begin{figure}[H]
  \centering
  \includegraphics[width=0.5\linewidth,height=0.42\textheight,keepaspectratio]{Finanzas20corporativas203ed20Dumrauf-21}
  \caption{Intereses generados por pasivos y por activos.}
\end{figure}

La tasa efectiva del impuesto difiere de la legal por las diferencias entre balance contable y fiscal y por los quebrantos previos; para el año siguiente se consulta a los asesores impositivos, y en proyecciones largas se usa la tasa legal.

\subsection{Proyección del balance}

\paragraph{Rubros discrecionales.} Dependen de decisiones de la dirección. Los \textbf{bienes de uso} no guardan una relación estable con las ventas: la relación baja con la depreciación y salta con cada inversión. El principio es comparar las ventas proyectadas en unidades con la capacidad instalada y prever la inversión cuando se la supera; en la práctica se mantienen relaciones como bienes de uso/ventas y capex/depreciación cercanas a las recientes. Los \textbf{intangibles} siguen su calendario de amortización, y las \textbf{deudas financieras}, su cronograma y la política de endeudamiento.

\paragraph{Rubros de generación espontánea.} Cuando crecen las ventas aumentan solos los créditos por ventas, los inventarios y las deudas comerciales. Su proyección es la parte más sensible del cash flow, sobre todo la de los créditos, y se hace con los días de antigüedad como dato de entrada; por eso antes se proyecta el estado de resultados.

\begin{formulaBox}[Rubros espontáneos a partir de los días]
\[
\text{Cuentas a cobrar} = \frac{\text{Días de cobranza}}{365}\times\text{Ventas}
\qquad
\text{Inventarios} = \frac{\text{Días de inv.}}{365}\times CMV
\]
\[
\text{Compras} = CMV - EI + EF
\qquad
\text{Deudas comerciales} = \frac{\text{Días de pago}}{365}\times\text{Compras}
\]
\end{formulaBox}

\begin{ejemplo}[Días y ciclo de efectivo negativo]
Cobrar a 45 días implica mantener créditos por cerca del 12,5\,\% de las ventas anuales ($45/365$). Una empresa que cobra a 45 días, vende su stock en 20 y paga a 140 tiene un ciclo de efectivo negativo ($45+20-140=-75$ días): cobra antes de pagar, y el crecimiento le genera fondos.
\end{ejemplo}

\begin{figure}[H]
  \centering
  \includegraphics[width=0.5\linewidth,height=0.42\textheight,keepaspectratio]{Finanzas20corporativas203ed20Dumrauf-23}
  \caption{Ratios de actividad y ciclo de efectivo.}
\end{figure}

\paragraph{Rubros semiespontáneos.} Los créditos y deudas fiscales (anticipos y retenciones del impuesto a las ganancias, saldos de IVA) tienen una gestión propia. Un exportador acumula IVA a favor que solo recupera con trámites ante el fisco: se proyecta por días de recupero sobre el monto anual, y el excedente demorado puede tratarse como activo no operativo. Los demás créditos y deudas de corto plazo se proyectan como porcentaje de las ventas.

\paragraph{Integración y análisis de sensibilidad.} Proyectados el estado de resultados y el balance, el cash flow surge de su interacción: el patrimonio neto es el del año anterior más la utilidad no distribuida, y la caja excedente se acumula en inversiones transitorias, que son dividendos potenciales. El \textbf{análisis de sensibilidad} modifica de a una las premisas para ver su efecto; el FCF suele ser muy sensible a los días de cobranza y menos al resto. Algunas variables (endeudarse, adelantar gastos, atrasar deudas fiscales) son decisiones discrecionales que nunca se proyectan con rigor, y los nuevos proyectos deben incorporarse al modelo.

\subsection{Financiamiento externo necesario y crecimiento}

\paragraph{El método del porcentaje de ventas.} Separa las cuentas que varían en proporción directa con las ventas (activos, deudas comerciales espontáneas) de las que no lo hacen (deuda bancaria y de largo plazo, capital aportado). Las utilidades retenidas se calculan con la utilidad proyectada y la tasa de retención $b = 1-d$, donde $d$ es el \emph{payout}. La relación activos/ventas es la \emph{intensidad de capital}, recíproca de la rotación.

\begin{definicion}{Financiamiento externo necesario (EFN o FER)}
Es la diferencia entre el aumento previsto de los activos y el aumento previsto de los pasivos espontáneos más las utilidades retenidas: el financiamiento externo que requiere el crecimiento proyectado. Es uno de los resultados clave del pronóstico, y se obtiene construyendo los estados proyectados y calculando la brecha, o directamente con la fórmula.
\end{definicion}

\begin{formulaBox}[Financiamiento externo necesario]
\[
\begin{aligned}
EFN ={}& \frac{\text{Activos}}{\text{Ventas}}\,\Delta\text{Ventas}
- \frac{\text{Pasivos espontáneos}}{\text{Ventas}}\,\Delta\text{Ventas}\\
&- PM\times\text{Ventas proyectadas}\times(1-d)
\end{aligned}
\]
donde $PM$ es el margen neto y $d$ el índice de pago de dividendos.
\end{formulaBox}

\begin{definicion}{Tasa de crecimiento interna y tasa de crecimiento sostenible}
La \textbf{tasa de crecimiento interna} es la máxima alcanzable sin ningún financiamiento externo: el aumento de activos se cubre exactamente con utilidades retenidas y el EFN es cero. La \textbf{tasa de crecimiento sostenible} (SGR) es la máxima posible sin emitir acciones y manteniendo constante la relación deuda/capital, es decir, sin aumentar el apalancamiento.
\end{definicion}

\begin{formulaBox}[Tasas de crecimiento]
\[
g_{\text{interna}} = \frac{ROA\times b}{1-ROA\times b}
\qquad\qquad
g_{\text{sostenible}} = \frac{ROE\times b}{1-ROE\times b}
\]
Si el ROE se mide sobre el patrimonio inicial, la tasa sostenible se reduce a $ROE\times b$.
\[
SGR = \frac{b\,(NP/S)\,(1+D/E)}{(A/S)-b\,(NP/S)\,(1+D/E)}
\]
con $A/S$ la relación activos/ventas, $NP/S$ el margen neto, $b$ la retención y $D/E$ la relación deuda/capital objetivo (versión de estado estable).
\end{formulaBox}

Por la identidad de Du Pont, el crecimiento sostenible depende de cuatro factores: el margen de utilidad, la rotación de activos, la política de dividendos y la política financiera. Con esos factores fijos y sin emitir acciones, solo hay una tasa de crecimiento posible; si las ventas deben crecer más, hay que mejorar el margen o la rotación, endeudarse más, retener más o emitir acciones. Si el crecimiento real supera la SGR, algo está cambiando, casi siempre el endeudamiento. Cuando las variables cambian año a año, la SGR de un año puede no repetirse; conociendo cinco de las seis variables puede despejarse la sexta, y la simulación muestra la sensibilidad del crecimiento a cada una. Para crecer en forma balanceada, el capital debe crecer en proporción a las ventas.

%% ---------------------------------------------------------------------
\begin{control}
\pregunta{1. ¿Cuál es el tema principal?}
La obtención del flujo de fondos a partir de los estados contables y su proyección para planificar las necesidades financieras de la empresa.

\pregunta{2. ¿Dónde y cómo se desarrolla?}
Primero se describen y depuran los estados contables y los efectos de impuestos, inflación y devaluación; luego se analizan con razones y con la identidad de Du Pont; por último, se construye el cash flow por el método indirecto, se distinguen FCF, ECF y CCF y se proyectan resultados, balance y financiamiento externo.

\pregunta{3. ¿Cuáles son los temas secundarios?}
Devengado y partidas no monetarias; capital de trabajo y ciclos; planeamiento impositivo; manipulación del flujo; políticas de capital de trabajo; regresión y series de tiempo para las ventas; rubros discrecionales, espontáneos y semiespontáneos; EFN y tasas de crecimiento.

\pregunta{4. ¿Las ideas principales y secundarias están relacionadas entre sí?}
Sí. Los estados contables son la materia prima; las razones diagnostican su calidad; el cash flow los traduce a efectivo; y la proyección integrada permite anticipar cuánto financiamiento exige el crecimiento y hasta dónde puede crecer la empresa sin cambiar su política financiera.

\pregunta{5. ¿Cuál es la conclusión o la idea general?}
El valor depende del efectivo y no de la utilidad contable. Proyectar estados consistentes, empezando por las ventas, permite estimar el FCF para valuar, el flujo del accionista para distribuir y el financiamiento externo necesario para crecer de manera sostenible.
\end{control}

%% =====================================================================
%%  UNIDAD 3 (parte A)
%% =====================================================================
\chapter{Estructura de capital, costo de capital y punto de equilibrio}\label{cap:u3}

\begin{ficha}
\fichabase{3 — Estructura de capital, costo de capital y punto de equilibrio}
\fila{Texto base}{S. Ross, R. Westerfield y J. Jaffe, \emph{Finanzas corporativas}, 9.ª ed., McGraw-Hill: cap.~16 «Estructura de capital: conceptos básicos» y cap.~17 «Estructura de capital: límites al uso de la deuda» (con su esquema de repaso); cap.~7 «Análisis de riesgos, opciones reales y presupuesto de capital» (análisis de escenarios y del punto de equilibrio).}
\fila{}{L.\,J. Gitman, \emph{Principios de administración financiera}, Pearson: cap.~9 «El costo de capital», apartado 9.5 «Costo de capital promedio ponderado» y resumen.}
\fila{}{Texto de cátedra «Punto de equilibrio», apartados 7.1 a 7.3 (elaboración y evaluación de proyectos de inversión).}
\fila{Complementaria}{R. Brealey, S. Myers y F. Allen, \emph{Principios de finanzas corporativas}, McGraw-Hill: caps. «¿Importa la política de endeudamiento?», «¿Cuánto debe endeudarse una empresa?» y «Financiamiento y valuación».}
\fila{}{G.\,L. Dumrauf, \emph{Finanzas corporativas: un enfoque latinoamericano}, 3.ª ed., Alfaomega: estructura de capital, métodos de valuación por descuento de flujos, punto de equilibrio y \emph{leverage} (apuntes).}
\fila{}{Ross, Westerfield y Jaffe, ob. cit.: resúmenes de «Financiamiento a largo plazo: una introducción» y «Valoración y presupuesto de capital para la empresa apalancada»; preguntas de repaso.}
\fila{Fórmulas previas}{Valor actual y perpetuidades \vease{sec:u1-va}; flujos FCF, ECF y CCF \vease{sec:u2-flujos}.}
\end{ficha}

%% ---------------------------------------------------------------------
\section{Financiamiento a largo plazo: una introducción}\label{sec:u3-flp}

\subsection{Acciones comunes}

Cada acción tiene un valor declarado o \textbf{valor nominal}; el número de acciones emitidas por ese valor forma el \emph{capital dedicado} de la sociedad. Lo aportado por encima del nominal es el \textbf{superávit de capital} y, en la mayoría de las jurisdicciones, queda bloqueado: solo se distribuye en la liquidación.

\begin{definicion}{Acciones autorizadas, emitidas y en cartera}
El estatuto de una sociedad fija el número de acciones comunes que está \textbf{autorizada} a emitir; las \textbf{emitidas} son las que efectivamente colocó. Las acciones \textbf{en cartera} o \textbf{de tesorería} son acciones ya emitidas que la propia empresa recompró.
\end{definicion}

Los accionistas votan para elegir el directorio. La \emph{votación acumulativa} permite que las minorías tengan representación; la \emph{votación directa} puede excluirlas, por lo que muchas jurisdicciones imponen la primera. El dividendo queda a discreción del directorio, no es un gasto del negocio, sale de las ganancias después de impuestos y, para el accionista individual, es renta gravada. La evidencia muestra que los gerentes prefieren acciones con votos superiores, que cotizan alrededor de un 5\,\% por encima de acciones idénticas con voto inferior.

\subsection{Deuda corporativa de largo plazo}

El interés es un costo del negocio y es deducible del impuesto, de modo que el Estado subsidia la deuda frente al capital; por eso las empresas a veces diseñan títulos que son deuda en apariencia y capital en sustancia, para obtener la ventaja fiscal sin cargar con los costos de quiebra.

\begin{definicion}{Instrumentos de deuda de largo plazo}
La deuda de largo plazo suele tener un valor nominal igual al facial, y su precio se expresa como porcentaje de ese valor. Un \textbf{\emph{debenture}} es deuda no garantizada; un \textbf{bono}, en sentido estricto, está garantizado con una hipoteca sobre bienes de la sociedad, aunque en el uso corriente la palabra designa ambos; un \textbf{pagaré} es deuda no garantizada de plazo más corto (quizá menos de diez años). La devolución en cuotas regulares se llama \textbf{amortización}, y el contrato que fija vencimiento, tasa y demás condiciones es la \textbf{escritura de emisión}.
\end{definicion}

\subsection{Acciones preferidas}

Tienen un valor de liquidación declarado, por lo general \$\,100, y su dividendo se expresa en pesos por acción. El directorio puede no pagarlo; según las condiciones, el dividendo es \emph{acumulativo} o no, y los atrasados y los corrientes deben cancelarse antes de pagar algo a los accionistas comunes. A diferencia de los intereses, el dividendo preferido no es deducible. Para el inversor individual es renta gravada, pero las empresas que las poseen tienen exento el 70\,\% del dividendo; por eso la mayoría de las preferidas está en manos de otras empresas.

\subsection{Arrendamiento}

\begin{definicion}{Arrendamiento operativo y arrendamiento financiero}
Un \textbf{arrendamiento} (\emph{leasing}) es un contrato por el cual el arrendatario usa un activo a cambio de pagos periódicos al arrendador, su propietario; por eso la decisión de comprar o arrendar es una decisión de financiamiento. El \textbf{operativo} no está totalmente amortizado, el arrendador mantiene y asegura el bien y, en su origen, incluía un operador junto con el equipo. El \textbf{financiero} es lo opuesto: no prevé mantenimiento del arrendador, está totalmente amortizado, el arrendatario suele poder renovarlo y, en general, no puede cancelarse.
\end{definicion}

\subsection{Patrones del financiamiento}

Históricamente, las empresas destinaron alrededor del 80\,\% de su flujo de efectivo a gastos de capital y el 20\,\% a capital de trabajo. El \textbf{financiamiento interno} (utilidad neta más depreciación menos dividendos) dominó como fuente y cubrió entre el 70 y el 90\,\% del financiamiento de largo plazo; el déficit se cubrió con préstamos y emisiones, aunque las emisiones netas de acciones fueron pequeñas e incluso negativas a fines de los años ochenta y noventa. Entre 1984 y 1990, y otra vez después de 1993, las empresas estadounidenses emitieron mucha deuda para retirar acciones: a valores contables la estructura de capital cambió, pero a valores de mercado el ratio de deuda se mantuvo notablemente estable, razón por la cual los economistas prefieren medirlo a valores de mercado.

El exceso en cualquiera de las dos fuentes tiene costos. Cuanta más deuda, mayor la probabilidad de incumplir los contratos; cuantas más acciones nuevas, mayor la dilución del porcentaje de propiedad y de los votos de los accionistas existentes, que por eso no siempre reciben bien las emisiones.

%% ---------------------------------------------------------------------
\section{Estructura de capital: conceptos básicos}\label{sec:u3-conceptos}

\subsection{La cuestión de la estructura de capital y la teoría del pastel}

¿Cómo debe elegir una empresa su razón de deuda a capital? El \emph{modelo del pastel} responde que la empresa es la suma de los derechos financieros sobre ella. El tamaño del pastel depende primero de las decisiones de inversión, que determinan el valor de los activos; la estructura de capital decide cómo se reparte.

\begin{formulaBox}[Valor de la empresa]
\[
V = B + S
\]
donde $B$ es el valor de mercado de la deuda y $S$ el valor de mercado del capital accionario. Si la meta es que la empresa valga lo más posible, debe elegirse la razón deuda/capital que haga más grande el pastel.
\end{formulaBox}

\begin{figure}[H]
  \centering
  \includegraphics[width=0.5\linewidth,height=0.42\textheight,keepaspectratio]{U3-24}
  \caption{Dos modelos de pastel de la estructura de capital: 40\,\% de deuda y 60\,\% de capital, y a la inversa (figura 16.1).}
\end{figure}

\begin{definicion}{Balance a valor de mercado}
Los contadores valúan las partidas a costo histórico; los financistas, a valor de mercado. El balance a valor de mercado muestra del lado izquierdo los activos (incluso los intangibles, como el escudo fiscal que se verá luego) y del derecho la deuda y el capital a sus valores de mercado.
\end{definicion}

\subsection{Maximización del valor de la empresa frente a maximización de los intereses de los accionistas}

¿Por qué los accionistas deberían interesarse en maximizar el valor de toda la empresa y no solo el de sus acciones? Porque, mientras la deuda no tenga riesgo considerable de incumplimiento, ambos objetivos coinciden.

\begin{definicion}{Principio de la estructura de capital}
Los cambios en la estructura de capital benefician a los accionistas si y solo si aumenta el valor de la empresa. Por lo tanto, los administradores deben elegir la estructura que consideren de mayor valor para la empresa, porque es la más beneficiosa para los accionistas.
\end{definicion}

Cuando el valor se reparte entre deuda y capital, las acciones son \emph{capital propio apalancado}: todo aumento o caída de valor causado por un cambio de estructura recae sobre los accionistas. Bajo dos supuestos —que la política de dividendos puede ignorarse y que el valor de la deuda permanece constante— la política que maximiza el valor de mercado de la empresa es también la mejor para ellos. Queda por responder la segunda pregunta: qué razón deuda/capital maximiza ese valor.

\subsection{Apalancamiento financiero y valor de la empresa: un ejemplo}

\paragraph{Apalancamiento y rendimientos de los accionistas.} Trans Am Corporation no tiene deuda, sus activos valen \$\,8.000 y tiene 400 acciones de \$\,20. Analiza emitir \$\,4.000 de deuda al 10\,\% para recomprar la mitad de las acciones. El efecto depende de la utilidad antes de intereses (UAII), que se estima en \$\,400 en recesión, \$\,1.200 en el escenario esperado y \$\,2.000 en expansión.

\begin{ejemplo}[Trans Am: la estructura actual frente a la propuesta]
\begin{center}
\begin{tabular}{@{}lrrr@{}}
\toprule
 & Recesión & Esperado & Expansión\\
\midrule
Rendimiento de los activos (ROA) & 5\,\% & 15\,\% & 25\,\%\\
Utilidad antes de intereses (UAII) & 400 & 1.200 & 2.000\\
\midrule
\multicolumn{4}{@{}l}{\emph{Sin deuda (capital 8.000; 400 acciones)}}\\
ROE & 5\,\% & 15\,\% & 25\,\%\\
UPA & 1,00 & 3,00 & 5,00\\
\midrule
\multicolumn{4}{@{}l}{\emph{Con deuda de 4.000 al 10\,\% (capital 4.000; 200 acciones)}}\\
Intereses & $-400$ & $-400$ & $-400$\\
Utilidad después de intereses & 0 & 800 & 1.600\\
ROE & 0\,\% & 20\,\% & 40\,\%\\
UPA & 0,00 & 4,00 & 8,00\\
\bottomrule
\end{tabular}
\end{center}
El ROA es igual en ambas estructuras porque se calcula antes de intereses.
\end{ejemplo}

\begin{figure}[H]
  \centering
  \includegraphics[width=0.5\linewidth,height=0.42\textheight,keepaspectratio]{U3-25}
  \caption{Apalancamiento financiero: UPA y UAII de Trans Am con y sin deuda (figura 16.2).}
\end{figure}

La recta sin deuda parte del origen; la recta con deuda corta el eje en una UPA negativa, porque los \$\,400 de intereses se pagan cualquiera sea la utilidad, pero tiene mayor pendiente, porque cada peso adicional de utilidad se reparte entre menos acciones. Por eso se cruzan en un \textbf{punto de indiferencia}: con una UAII de \$\,800 ambas estructuras dan una UPA de \$\,2. Por encima de ese nivel conviene la deuda; por debajo, el capital.

\paragraph{La visión ingenua y el análisis EBIT-EPS.} El caso parece decir que el apalancamiento es bueno, porque la UPA esperada sube de 3 a 4 (\$\,1.200 entre 400 acciones sin deuda; \$\,800, después de intereses, entre 200 acciones con ella). Esa es la \emph{visión ingenua} del \emph{leverage}: un negocio que rinde 20\,\% antes de impuestos financiado solo con capital pasa a rendir 30\,\% sobre el capital si la mitad se financia con deuda al 10\,\% (\$\,15 de utilidad sobre \$\,50 invertidos), y parece conveniente seguir sustituyendo acciones por deuda. Se olvida que el mayor rendimiento viene con mayor riesgo, por el que los accionistas exigen compensación.

\begin{definicion}{Análisis EBIT-EPS}
Calcula la ganancia por acción (EPS o UPA) que resulta de cada alternativa de financiamiento en distintos escenarios económicos. Si el rendimiento de los activos supera el costo de la deuda, la palanca opera a favor y la deuda eleva la UPA; si es igual, la elección es indiferente; si es menor, la palanca opera en contra.
\end{definicion}

\begin{ejemplo}[Tres formas de financiar \$\,100 que llevan el activo a \$\,200 ($k_d = 10\,\%$)]
\begin{center}
\begin{tabular}{@{}lccc@{}}
\toprule
Rendimiento del activo (EBIT) & Solo acciones & Mitad y mitad & Solo deuda\\
\midrule
20\,\% (EBIT 40) & 0,20 & 0,23 & 0,30\\
10\,\% (EBIT 20) & 0,10 & 0,10 & 0,10\\
5\,\% (EBIT 10) & 0,05 & 0,03 & 0,00\\
\bottomrule
\end{tabular}
\end{center}
Las cifras son ganancias por acción (200, 150 y 100 acciones de \$\,1, respectivamente). Del mismo modo, con intereses fijos de 5, un EBIT de 10, 20 o 30 ($\pm 50\,\%$) deja un resultado de 5, 15 o 25 ($\pm 66{,}7\,\%$).
\end{ejemplo}

Comparadas las alternativas de financiamiento, la más endeudada ofrece el mejor resultado en expansión y el peor en recesión: como los intereses son fijos, la deuda amplifica la variabilidad de la ganancia por acción \vease{sec:u3-leverage}. Aun así, el análisis EBIT-EPS tiene dos limitaciones que lo vuelven inadecuado para elegir la estructura óptima: no considera el riesgo financiero que genera la deuda ni su efecto en el costo del capital, y persigue maximizar la ganancia por acción, que no equivale necesariamente a maximizar la riqueza del accionista.

\paragraph{La elección entre deuda y capital.} Modigliani y Miller (MM) mostraron que la empresa no puede cambiar el valor total de sus títulos modificando sus proporciones. Su argumento compara dos estrategias para un inversor de Trans Am:
\begin{itemize}
\item \textbf{Estrategia A}: comprar 100 acciones de la empresa apalancada por \$\,2.000; obtiene 0, 400 u 800 según el escenario.
\item \textbf{Estrategia B} (\emph{apalancamiento casero}): pedir prestados \$\,2.000 al 10\,\% y, con \$\,2.000 propios, comprar 200 acciones de la empresa sin deuda; obtiene 200, 600 o 1.000 menos 200 de intereses, es decir, 0, 400 u 800, con el mismo costo neto de \$\,2.000.
\end{itemize}

Como ambas estrategias tienen el mismo costo y el mismo resultado, la empresa no ayuda ni perjudica a sus accionistas al endeudarse: el inversor no recibe del apalancamiento corporativo nada que no pueda obtener por su cuenta. Si la empresa apalancada valiera más, nadie compraría sus acciones, porque B sería más barata; el arbitraje bajaría su valor y subiría el de la no apalancada hasta igualarlos.

\begin{definicion}{Apalancamiento casero}
Es la sustitución del endeudamiento de la empresa por el endeudamiento personal del inversor. Mientras los individuos puedan pedir y otorgar préstamos en los mismos términos que las empresas, pueden replicar o deshacer por sí mismos cualquier apalancamiento corporativo.
\end{definicion}

\begin{formulaBox}[Proposición I de MM (sin impuestos)]
\[
V_L = V_U
\qquad\qquad
V = E + D = \frac{EBIT}{k_u}
\quad\Longleftrightarrow\quad
k_u = \frac{EBIT}{V}
\]
El valor de la empresa apalancada ($L$) es igual al de la no apalancada ($U$): resulta de descontar el resultado operativo a la tasa $k_u$ exigida a una empresa sin deuda de la misma clase de riesgo.
\end{formulaBox}

La Proposición I es una aplicación de la \emph{ley de conservación del valor}: el valor de un activo se conserva cualquiera sea la naturaleza de los derechos sobre él. El valor lo determinan los activos reales del lado izquierdo del balance, no su forma de financiarlos. Se la considera el punto de partida de las finanzas corporativas modernas.

\paragraph{Un supuesto fundamental.} El resultado depende de que los individuos se endeuden tan barato como las empresas. Es razonable: un inversor puede abrir una cuenta de margen con su corredor, que le presta parte del precio de las acciones (por ejemplo, compra \$\,10.000 de acciones con \$\,6.000 propios y \$\,4.000 prestados). Como el corredor conserva las acciones en garantía y exige reponer el margen cuando el precio cae, su riesgo es pequeño y cobra tasas cercanas a la libre de riesgo. Las empresas, en cambio, ofrecen en garantía activos ilíquidos y generan costos de negociación, supervisión y litigio; por eso es difícil sostener que los individuos se endeuden más caro.

\subsection{Modigliani y Miller: Proposición II (sin impuestos)}

\paragraph{El riesgo para los accionistas aumenta con el apalancamiento.} En Trans Am la UPA sin deuda varía entre 1 y 5, y con deuda entre 0 y 8: los accionistas apalancados ganan más en los buenos tiempos y pierden más en los malos. La mayor pendiente de su recta mide, justamente, su mayor sensibilidad a la UAII, es decir, su mayor riesgo; por eso exigen un rendimiento mayor (20\,\% en lugar de 15\,\%).

\begin{definicion}{Costo promedio ponderado del capital (WACC)}
Es el promedio de los costos de la deuda y del capital propio, ponderados por la proporción de cada fuente en la estructura de capital a valores de mercado. Mide el rendimiento esperado de los activos, es decir, la utilidad operativa esperada sobre el valor de mercado de todos los títulos, y por eso es el costo de capital de la empresa.
\end{definicion}

\begin{formulaBox}[WACC sin impuestos]
\[
R_{WACC} = \frac{B}{B+S}\,R_B + \frac{S}{B+S}\,R_S
\qquad\Longleftrightarrow\qquad
WACC = k_e\frac{E}{E+D} + k_d\frac{D}{E+D}
\]
La segunda es la notación de Dumrauf: $E$ y $D$ son los valores de mercado de las acciones y de la deuda.
\end{formulaBox}

\begin{figure}[H]
  \centering
  \includegraphics[width=0.5\linewidth,height=0.42\textheight,keepaspectratio]{U3-26}
\end{figure}
\begin{figure}[H]
  \centering
  \includegraphics[width=0.5\linewidth,height=0.42\textheight,keepaspectratio]{U3-27}
\end{figure}
\begin{figure}[H]
  \centering
  \includegraphics[width=0.5\linewidth,height=0.42\textheight,keepaspectratio]{U3-28}
  \caption{Costo promedio ponderado del capital de Trans Am con y sin deuda: 15\,\% en ambos casos (tabla 16.5).}
\end{figure}

Una implicación de la Proposición I es que el WACC de una empresa es constante, cualquiera sea su estructura; en ausencia de impuestos es igual a $R_0$, el costo de capital de la empresa financiada solo con acciones ($R_0 = 1.200/8.000 = 15\,\%$ en Trans Am).

\begin{formulaBox}[Proposición II de MM (sin impuestos)]
\[
R_S = R_0 + \frac{B}{S}\,(R_0 - R_B)
\qquad\Longleftrightarrow\qquad
k_e = k_u + (k_u - k_d)\,\frac{D}{E}
\]
\emph{Derivación}: igualando el WACC a $R_0$ y multiplicando por $(B+S)/S$ se obtiene $\frac{B}{S}R_B + R_S = \frac{B}{S}R_0 + R_0$. Del mismo modo, $k_e = \dfrac{EBIT - k_dD}{E} = \dfrac{k_u(E+D) - k_dD}{E}$. En Trans Am apalancada: $R_S = 15\,\% + \frac{4.000}{4.000}(15\,\% - 10\,\%) = 20\,\%$.
\end{formulaBox}

El rendimiento requerido del capital es una función lineal de la razón deuda/capital, con pendiente $R_0 - R_B$, positiva porque aun el capital sin deuda es riesgoso. Cada peso de capital se apalanca con más deuda, su riesgo aumenta y también su rendimiento exigido; el WACC, en cambio, no varía. En el gráfico, $R_0$ es un solo punto, mientras que el WACC es una línea horizontal entera.

\begin{figure}[H]
  \centering
  \includegraphics[width=0.5\linewidth,height=0.42\textheight,keepaspectratio]{U3-29}
\end{figure}
\begin{figure}[H]
  \centering
  \includegraphics[width=0.5\linewidth,height=0.42\textheight,keepaspectratio]{U3-30}
\end{figure}
\begin{figure}[H]
  \centering
  \includegraphics[width=0.5\linewidth,height=0.42\textheight,keepaspectratio]{U3-31}
\end{figure}
\begin{figure}[H]
  \centering
  \includegraphics[width=0.5\linewidth,height=0.42\textheight,keepaspectratio]{U3-32}
\end{figure}
\begin{figure}[H]
  \centering
  \includegraphics[width=0.5\linewidth,height=0.42\textheight,keepaspectratio]{U3-33}
\end{figure}
\begin{figure}[H]
  \centering
  \includegraphics[width=0.5\linewidth,height=0.42\textheight,keepaspectratio]{U3-34}
\end{figure}
\begin{figure}[H]
  \centering
  \includegraphics[width=0.5\linewidth,height=0.42\textheight,keepaspectratio]{U3-35}
\end{figure}
\begin{figure}[H]
  \centering
  \includegraphics[width=0.5\linewidth,height=0.42\textheight,keepaspectratio]{U3-36}
  \caption{Proposición II sin impuestos: verificación para Trans Am y costo del capital, costo de la deuda y WACC en función de $B/S$ (figura 16.3).}
\end{figure}

\begin{figure}[H]
  \centering
  \includegraphics[width=0.5\linewidth,height=0.42\textheight,keepaspectratio]{U3-37}
\end{figure}
\begin{figure}[H]
  \centering
  \includegraphics[width=0.5\linewidth,height=0.42\textheight,keepaspectratio]{U3-38}
  \caption{Síntesis de las proposiciones de MM sin impuestos.}
\end{figure}

\paragraph{Supuestos y proposiciones.} La tesis de MM, publicada en 1958 para hallar la tasa de corte de las inversiones financiadas con deuda, acciones o ambas, se hizo célebre por afirmar que en mercados eficientes la combinación de deuda y acciones es irrelevante. Descansa en cinco supuestos:
\begin{enumerate}
\item mercados de capitales perfectos, sin costos de transacción, con inversores racionales que maximizan su riqueza;
\item un resultado operativo futuro aleatorio sobre cuyo valor esperado todos coinciden;
\item empresas agrupables en clases de rendimiento equivalente, con rendimientos proporcionales dentro de cada clase;
\item deuda libre de riesgo, a cuya tasa se endeudan empresas e individuos;
\item ausencia de impuestos corporativos y personales y de costos de quiebra.
\end{enumerate}
A las dos proposiciones se agrega una tercera.

\begin{definicion}{Proposición III de MM}
Una empresa que actúa en beneficio de sus accionistas solo debe emprender las inversiones cuyo rendimiento supere el costo de capital total de la empresa ($k_u$), cualquiera sea la fuente con que las financie.
\end{definicion}

\begin{ejemplo}[Una inversión financiada con deuda barata que destruye valor]
Una empresa sin deuda tiene un EBIT perpetuo de 100 y $k_u = 10\,\%$: $V_0 = E_0 = 1.000$. Invierte 100, financiados con deuda al 4\,\%, en un proyecto que rinde 8\,\%. Entonces $V_1 = (100 + 8)/0{,}10 = 1.080$ y $E_1 = 1.080 - 100 = 980$: los accionistas pierden 20, aunque el proyecto rinde el doble que la deuda, porque rinde menos que $k_u$. Lo confirma la Proposición II: $k_e = 0{,}10 + (0{,}10 - 0{,}04)\,\frac{100}{980} = 10{,}61\,\%$ y $E = (108 - 4)/0{,}1061 = 980$.
\end{ejemplo}

\paragraph{MM: una interpretación.} Los administradores no pueden cambiar el valor de la empresa «reempacando» sus títulos. Aunque la deuda parezca más barata, al sustituir capital por deuda el capital restante se vuelve más riesgoso y más caro, y ambos efectos se compensan con exactitud. MM lo ilustran con un productor de leche: vender leche entera o separarla en crema y leche descremada no cambia el ingreso total, porque si la segunda opción valiera más los arbitrajistas comprarían leche entera, la descremarían y venderían las partes. El valor de la leche no depende de cómo se envasa, ni el tamaño del pastel de cómo se corta. Según la teoría, las razones deuda/capital serían caprichosas, fruto de accidentes históricos; en la realidad, sin embargo, cada industria tiene una razón típica (alta en la banca, baja en la farmacéutica), lo que sugiere que la teoría omite factores del mundo real: los impuestos y los costos de quiebra y de agencia.

\paragraph{La estructura de capital y el beta.} Accionistas y acreedores comparten los flujos y el riesgo de la empresa. El beta de los activos es el de una cartera con toda la deuda y todo el capital, es decir, el promedio ponderado de los betas de ambos. El apalancamiento no cambia el riesgo ni el rendimiento esperado de los activos, pero sí el riesgo de las acciones, y por eso los accionistas exigen más rendimiento \vease{sec:u3-hamada}.

\paragraph{La visión tradicional.} Frente a MM, autores como David Durand y Ely Schwartz sostuvieron que, con endeudamiento moderado, el costo del capital propio y el de la deuda se mantienen más o menos estables y solo se disparan pasado cierto límite. Así, el encarecimiento del capital propio no alcanza a anular la economía de usar deuda: el WACC primero baja, toca un mínimo en un endeudamiento intermedio y luego sube, en forma de U. Minimizar el WACC equivaldría a maximizar el valor si el resultado operativo no cambia con la deuda. La postura se apoya en dos argumentos: los inversores podrían no percibir el riesgo de un endeudamiento moderado, y los mercados reales son imperfectos aunque MM tenga razón en los perfectos. De allí que los gerentes busquen una \emph{clientela insatisfecha}: inversores que desean un instrumento que las imperfecciones del mercado, muchas veces creadas por los gobiernos, les impiden obtener a bajo costo.

Dos advertencias completan el cuadro. Los accionistas quieren enriquecerse, no ser dueños de una empresa con WACC bajo. Y tratar de minimizar el WACC invita a un «cortocircuito lógico»: suponer que la deuda es una fuente barata y que endeudarse más reduce el costo total, cuando los accionistas exigirán más por el mayor riesgo.

\subsection{Impuestos}\label{sec:u3-impuestos}

\paragraph{Percepción básica.} Con impuestos corporativos el pastel tiene una porción más: la del fisco. Si el tamaño total no cambia, el valor de deuda más capital es máximo con la estructura que paga menos impuestos, es decir, la que «menos le gusta al fisco». Como los intereses son deducibles y los dividendos no, la empresa apalancada paga menos: lo que reciben en conjunto sus accionistas y acreedores supera lo que reciben los accionistas de una empresa igual sin deuda exactamente en el impuesto ahorrado, la alícuota por los intereses. Que sus accionistas cobren menos en total no los perjudica: si la deuda reemplaza acciones, quedan menos acciones y la UPA es mayor.

\begin{figure}[H]
  \centering
  \includegraphics[width=0.5\linewidth,height=0.42\textheight,keepaspectratio]{U3-39}
  \caption{Dos modelos de pastel de la estructura de capital en presencia de impuestos corporativos (figura 16.4).}
\end{figure}

\begin{ejemplo}[Water Products Company: plan I sin deuda y plan II con deuda]
Con una UAII de \$\,1.000.000 y una tasa del 35\,\%, el plan I paga \$\,350.000 de impuestos y entrega \$\,650.000 a los accionistas. El plan II tiene \$\,4.000.000 de deuda al 10\,\%: paga \$\,400.000 de intereses, \$\,210.000 de impuestos y entrega \$\,390.000 a los accionistas, de modo que accionistas y acreedores reciben juntos \$\,790.000. La diferencia de \$\,140.000 es exactamente el menor impuesto ($0{,}35\times 400.000$).
\end{ejemplo}

\begin{figure}[H]
  \centering
  \includegraphics[width=0.5\linewidth,height=0.42\textheight,keepaspectratio]{U3-40}
  \caption{Flujos para accionistas, acreedores y fisco en una empresa sin deuda (plan I) y con deuda (plan II), caso Water Products Company.}
\end{figure}

\begin{definicion}{Escudo fiscal de la deuda}
Es el ahorro de impuestos que generan los intereses deducibles: el fisco subsidia a la empresa que se endeuda. Es un monto anual igual a la alícuota por los intereses. Si la empresa espera tener ganancias gravables, el escudo tiene el mismo riesgo que los intereses y se descuenta a la tasa de la deuda.
\end{definicion}

\begin{formulaBox}[Escudo fiscal y su valor presente]
\[
\text{Escudo anual} = t_c\,R_B\,B
\qquad\qquad
VP(\text{escudo perpetuo}) = \frac{t_c\,R_B\,B}{R_B} = t_c\,B
\]
Si la deuda tiene vida finita, el valor del escudo es la suma de los escudos de cada año, que bajan a medida que se amortiza el capital, descontados a $R_B$; esa suma es lo que la empresa endeudada vale de más frente a una igual sin deuda.
\end{formulaBox}

\begin{ejemplo}[Deuda con vida finita]
Maxwell Company toma \$\,1.000.000 al 8\,\% que amortiza en dos cuotas iguales; $R_B = 10\,\%$ y $t_c = 35\,\%$. Los intereses son \$\,80.000 y \$\,40.000, y los escudos, \$\,28.000 y \$\,14.000. Su valor presente es $\frac{0{,}35\times 80.000}{1{,}10} + \frac{0{,}35\times 40.000}{1{,}10^2} = 25.454{,}54 + 11.570{,}25 = \$\,37.024{,}79$.
\end{ejemplo}

\paragraph{Valor de la empresa apalancada.} El flujo anual de la empresa sin deuda es $UAII\,(1-t_c)$, y su valor, ese flujo descontado a $R_0$. Endeudarse agrega el valor del escudo.

\begin{figure}[H]
  \centering
  \includegraphics[width=0.5\linewidth,height=0.42\textheight,keepaspectratio]{U3-41}
\end{figure}
\begin{figure}[H]
  \centering
  \includegraphics[width=0.5\linewidth,height=0.42\textheight,keepaspectratio]{U3-42}
\end{figure}
\begin{figure}[H]
  \centering
  \includegraphics[width=0.5\linewidth,height=0.42\textheight,keepaspectratio]{U3-43}
\end{figure}
\begin{figure}[H]
  \centering
  \includegraphics[width=0.5\linewidth,height=0.42\textheight,keepaspectratio]{U3-44}
\end{figure}
\begin{figure}[H]
  \centering
  \includegraphics[width=0.5\linewidth,height=0.42\textheight,keepaspectratio]{U3-45}
  \caption{Valor de la empresa no apalancada y de la apalancada con impuestos corporativos: elementos de la Proposición I.}
\end{figure}

\begin{formulaBox}[Proposición I de MM con impuestos corporativos]
\[
V_U = \frac{UAII\,(1-t_c)}{R_0}
\qquad\qquad
V_L = \frac{UAII\,(1-t_c)}{R_0} + \frac{t_c R_B B}{R_B} = V_U + t_c\,B
\]
Vale con deuda constante en el tiempo; si lo constante fuera la razón deuda/capital, la fórmula sería otra.
\end{formulaBox}

Como el escudo crece con la deuda, la empresa aumenta su flujo total y su valor sustituyendo capital por deuda: la recta de $V_L$ crece con pendiente $t_c$. Las figuras siguientes lo aplican a Divided Airlines, una empresa sin deuda que evalúa emitir \$\,200 de deuda al 10\,\% para recomprar acciones.

\begin{ejemplo}[Divided Airlines]
Espera una UAII perpetua de \$\,153,85; con $t_c = 35\,\%$ gana \$\,100 después de impuestos y los reparte. Las empresas sin deuda del sector tienen $R_0 = 20\,\%$: $V_U = 100/0{,}20 = \$\,500$. Con la emisión, $V_L = 500 + 0{,}35\times 200 = \$\,570$, y el capital vale $570 - 200 = \$\,370$.
\end{ejemplo}

\begin{figure}[H]
  \centering
  \includegraphics[width=0.5\linewidth,height=0.42\textheight,keepaspectratio]{U3-46}
\end{figure}
\begin{figure}[H]
  \centering
  \includegraphics[width=0.5\linewidth,height=0.42\textheight,keepaspectratio]{U3-47}
\end{figure}
\begin{figure}[H]
  \centering
  \includegraphics[width=0.5\linewidth,height=0.42\textheight,keepaspectratio]{U3-48}
  \caption{Proposición I con impuestos: el valor crece con la deuda; aplicación a Divided Airlines.}
\end{figure}

\paragraph{Rendimiento esperado y apalancamiento bajo impuestos corporativos.} La intuición de la Proposición II se mantiene: el riesgo del capital crece con la deuda. Como el escudo absorbe parte de ese riesgo, la pendiente se reduce por el factor $(1-t_c)$.

\begin{formulaBox}[Proposición II de MM con impuestos y valor del capital apalancado]
\[
R_S = R_0 + \frac{B}{S}\,(1-t_c)\,(R_0 - R_B)
\qquad\qquad
S = \frac{(UAII - R_B B)(1-t_c)}{R_S}
\]
\emph{Derivación}: en el balance a valor de mercado, $V_U + t_cB = B + S$. Igualando el flujo esperado de los activos, $V_UR_0 + t_cBR_B$, con el que reciben accionistas y acreedores, $SR_S + BR_B$, y usando $V_U = S + (1-t_c)B$, se despeja $R_S$.
\end{formulaBox}

\begin{ejemplo}[Divided Airlines: costo del capital y WACC]
$R_S = 0{,}20 + \frac{200}{370}\times 0{,}65\times(0{,}20-0{,}10) = 23{,}51\,\%$. Comprobación: $S = (153{,}85-20)\times 0{,}65/0{,}2351 = \$\,370$. El WACC con impuestos es $\frac{370}{570}\times 0{,}2351 + \frac{200}{570}\times 0{,}10\times 0{,}65 = 17{,}54\,\%$, menor que el 20\,\% sin deuda, y $V_L = 100/0{,}1754 = \$\,570$.
\end{ejemplo}

\begin{figure}[H]
  \centering
  \includegraphics[width=0.5\linewidth,height=0.42\textheight,keepaspectratio]{U3-49}
\end{figure}
\begin{figure}[H]
  \centering
  \includegraphics[width=0.5\linewidth,height=0.42\textheight,keepaspectratio]{U3-50}
  \caption{Proposición II con impuestos corporativos aplicada a Divided Airlines.}
\end{figure}

\begin{figure}[H]
  \centering
  \includegraphics[width=0.5\linewidth,height=0.42\textheight,keepaspectratio]{U3-51}
\end{figure}
\begin{figure}[H]
  \centering
  \includegraphics[width=0.5\linewidth,height=0.42\textheight,keepaspectratio]{U3-52}
\end{figure}
\begin{figure}[H]
  \centering
  \includegraphics[width=0.5\linewidth,height=0.42\textheight,keepaspectratio]{U3-53}
\end{figure}
\begin{figure}[H]
  \centering
  \includegraphics[width=0.5\linewidth,height=0.42\textheight,keepaspectratio]{U3-54}
\end{figure}
\begin{figure}[H]
  \centering
  \includegraphics[width=0.5\linewidth,height=0.42\textheight,keepaspectratio]{U3-55}
\end{figure}
\begin{figure}[H]
  \centering
  \includegraphics[width=0.5\linewidth,height=0.42\textheight,keepaspectratio]{U3-56}
\end{figure}
\begin{figure}[H]
  \centering
  \includegraphics[width=0.5\linewidth,height=0.42\textheight,keepaspectratio]{U3-57}
\end{figure}
\begin{figure}[H]
  \centering
  \includegraphics[width=0.5\linewidth,height=0.42\textheight,keepaspectratio]{U3-58}
  \caption{Valor del capital apalancado y efecto del apalancamiento sobre el costo de la deuda, el costo del capital y el WACC (figura 16.6).}
\end{figure}

\begin{formulaBox}[WACC con impuestos corporativos]
\[
R_{WACC} = \frac{S}{V_L}\,R_S + \frac{B}{V_L}\,R_B\,(1-t_c)
\qquad\qquad
V_L = \frac{UAII\,(1-t_c)}{R_{WACC}}
\]
El costo de la deuda se multiplica por $(1-t_c)$ porque los intereses son deducibles; el del capital no, porque los dividendos no lo son. Con impuestos, el WACC disminuye con el apalancamiento.
\end{formulaBox}

\begin{figure}[H]
  \centering
  \includegraphics[width=0.5\linewidth,height=0.42\textheight,keepaspectratio]{U3-59}
\end{figure}
\begin{figure}[H]
  \centering
  \includegraphics[width=0.5\linewidth,height=0.42\textheight,keepaspectratio]{U3-60}
\end{figure}
\begin{figure}[H]
  \centering
  \includegraphics[width=0.5\linewidth,height=0.42\textheight,keepaspectratio]{U3-61}
\end{figure}
\begin{figure}[H]
  \centering
  \includegraphics[width=0.5\linewidth,height=0.42\textheight,keepaspectratio]{U3-62}
\end{figure}
\begin{figure}[H]
  \centering
  \includegraphics[width=0.5\linewidth,height=0.42\textheight,keepaspectratio]{U3-63}
\end{figure}
\begin{figure}[H]
  \centering
  \includegraphics[width=0.5\linewidth,height=0.42\textheight,keepaspectratio]{U3-64}
  \caption{Valuación de Divided Airlines con el WACC y balance a valor de mercado de la empresa sin deuda.}
\end{figure}

\paragraph{El precio de las acciones y el apalancamiento bajo impuestos corporativos.} ¿Por qué maximizar el valor de la empresa y no el precio de la acción? Porque el aumento de valor llega íntegro a los accionistas, como muestra el caso de Divided Airlines. Sin deuda, la empresa vale \$\,500 y tiene 100 acciones de \$\,5. Al anunciar que emitirá \$\,200 de deuda para recomprar acciones, el mercado incorpora de inmediato el escudo de \$\,70 al balance (activos de 500 más 70): cada acción pasa a valer \$\,5,70, y el aumento ocurre el día del anuncio, no el del canje. Luego se emite la deuda y se recompran $200/5{,}70 = 35{,}09$ acciones; quedan 64,91 acciones que valen $370/64{,}91 = \$\,5{,}70$. Quienes venden sus acciones cobran lo mismo que valen las que se conservan, y los accionistas captan todo el escudo. El escudo es un activo aunque no sea tangible: tiene valor porque reduce los impuestos futuros.

\begin{figure}[H]
  \centering
  \includegraphics[width=0.5\linewidth,height=0.42\textheight,keepaspectratio]{U3-65}
\end{figure}
\begin{figure}[H]
  \centering
  \includegraphics[width=0.5\linewidth,height=0.42\textheight,keepaspectratio]{U3-66}
  \caption{Balances a valor de mercado de Divided Airlines al anunciarse la emisión de deuda y después del canje de deuda por acciones.}
\end{figure}

\begin{ejemplo}[Las empresas A y B (alícuota del 40\,\%)]
Ambas tienen un EBIT de 20 y $k_u = 20\,\%$. A no tiene deuda: paga 8 de impuestos, gana 12 y vale $V_U = 12/0{,}20 = 60$. B debe 50 al 10\,\%: paga 5 de intereses y 6 de impuestos, gana 9, y su ahorro fiscal es $50\times 0{,}10\times 0{,}40 = 2$ por año, que vale $50\times 0{,}40 = 20$. Entonces $V_L = 60 + 20 = 80$, $E = 30$, $k_e = 0{,}20 + (0{,}20-0{,}10)\,\frac{50}{30}\,0{,}60 = 30\,\%$ (o $9/30$) y $WACC = 0{,}30\times\frac{30}{80} + 0{,}10\times 0{,}60\times\frac{50}{80} = 15\,\% = 12/80$. El Estado paga, en valor presente, el 40\,\% de la deuda.
\end{ejemplo}

\begin{formulaBox}[WACC en función del apalancamiento (deuda perpetua)]
\[
WACC = k_u\,(1 - t\,L), \quad L = \frac{D}{V}
\qquad\qquad
V = \frac{EBIT(1-t)}{WACC} = \frac{EBIT(1-t)}{k_u} + D\,t
\]
$k_u$ es el rendimiento exigido a la empresa sin deuda, $t$ la alícuota y $L$ la proporción de deuda en el valor de mercado: el WACC baja linealmente a medida que crece $L$.
\end{formulaBox}

\begin{figure}[H]
  \centering
  \includegraphics[width=0.5\linewidth,height=0.42\textheight,keepaspectratio]{screenshot006}
  \caption{Con impuesto de sociedades el valor de la empresa crece linealmente con la deuda, el costo del capital propio sube, el de la deuda se mantiene y el WACC desciende.}
\end{figure}

\begin{definicion}{Proposiciones de MM con impuestos corporativos: síntesis}
\textbf{Supuestos}: las empresas pagan la tasa $t_c$ sobre las utilidades después de intereses; no hay costos de transacción; individuos y empresas se endeudan a la misma tasa. \textbf{Proposición I}: como los intereses son deducibles y los dividendos no, el apalancamiento reduce los impuestos y aumenta el valor ($V_L = V_U + t_cB$). \textbf{Proposición II}: el costo del capital propio crece con el apalancamiento porque crece su riesgo, aunque más despacio que sin impuestos.
\end{definicion}

\begin{figure}[H]
  \centering
  \includegraphics[width=0.5\linewidth,height=0.42\textheight,keepaspectratio]{U3-67}
\end{figure}
\begin{figure}[H]
  \centering
  \includegraphics[width=0.5\linewidth,height=0.42\textheight,keepaspectratio]{U3-68}
  \caption{Resultados de las proposiciones de MM con impuestos corporativos.}
\end{figure}

Los cálculos anteriores tienden a \textbf{sobrestimar} el escudo por tres razones: no es realista suponer una deuda fija y perpetua, porque la capacidad de endeudamiento cambia con las utilidades y el valor; muchas empresas enfrentan tasas marginales menores que la legal; y el escudo solo sirve si hay utilidades futuras que proteger, algo que nadie puede garantizar. Además, según la teoría todas las empresas deberían endeudarse al máximo, lo que contradice la realidad; los límites al uso de la deuda explican la diferencia \vease{sec:u3-limites}.

\subsection{El CAPM combinado con las proposiciones de MM: betas apalancadas}\label{sec:u3-hamada}

En 1972 Robert Hamada combinó la Proposición II con impuestos con el CAPM, modelo según el cual el rendimiento exigido a una acción es la tasa libre de riesgo más su beta por la prima de mercado \vease{sec:u5-capm}. El resultado separa tres componentes del rendimiento que exige el accionista: el valor tiempo del dinero ($r_f$), un premio por el \textbf{riesgo de negocio} ($\beta_u(r_m-r_f)$) y un premio por el \textbf{riesgo financiero}, proporcional a la deuda neta del escudo fiscal, $D(1-t)$.

\begin{definicion}{Beta del activo (desapalancado) y beta del capital (apalancado)}
El \textbf{beta del activo} o desapalancado ($\beta_u$) es el beta que tendría la acción si la empresa se financiara solo con capital: depende únicamente del riesgo de negocio. El \textbf{beta del capital} o apalancado ($\beta_e$) agrega el riesgo financiero que provoca la deuda.
\end{definicion}

\begin{formulaBox}[Betas y apalancamiento (Hamada)]
\[
\begin{gathered}
\beta_A = \beta_D\,\frac{D}{V} + \beta_E\,\frac{E}{V}
\qquad\qquad
\text{sin impuestos: } \beta_E = \beta_A\Bigl(1+\frac{D}{E}\Bigr)
\\[4pt]
\text{con impuestos y } \beta_D = 0: \quad
\beta_e = \beta_u\Bigl[1+\frac{D(1-t)}{E}\Bigr]
\qquad\Longleftrightarrow\qquad
\beta_u = \frac{\beta_e}{1+\frac{D(1-t)}{E}}
\\[4pt]
k_e = r_f + \beta_u\,(r_m - r_f) + \beta_u\,(r_m - r_f)\,\frac{D}{E}\,(1-t)
\end{gathered}
\]
\end{formulaBox}

En la práctica, el consultor \emph{desapalanca} el beta observado de una empresa con deuda, para llevar imaginariamente su endeudamiento a cero, y lo \emph{reapalanca} con la estructura de la empresa que valúa.

\begin{ejemplo}[Desapalancar y reapalancar un beta]
El beta de una acción es 1,66; la deuda y el capital valen 100 y 140 ($D/E = 0{,}714$) y $t = 40\,\%$. Beta del activo: $\beta_u = 1{,}66/[1 + 0{,}714\times 0{,}60] = 1{,}16$. Para una relación $D/E = 0{,}80$: $\beta_e = 1{,}16\,[1 + 0{,}80\times 0{,}60] = 1{,}72$.
\end{ejemplo}

%% =====================================================================
%%  UNIDAD 3 (parte B)
%% =====================================================================
\section{Estructura de capital: límites al uso de la deuda}\label{sec:u3-limites}

\subsection{Costos de las quiebras financieras}

La deuda ahorra impuestos, pero también presiona a la empresa, porque intereses y principal son obligaciones legales; los dividendos, en cambio, no lo son. Si las obligaciones no se cumplen, la empresa entra en dificultades financieras, y la peor de ellas es la \emph{quiebra}, en la que la propiedad de los activos pasa legalmente de los accionistas a los acreedores. Los costos de esa reorganización tienden a cancelar las ventajas de la deuda.

\begin{ejemplo}[Day Corporation (inversores neutrales al riesgo, tasa del 10\,\%)]
Los acreedores de Day, a quienes se prometen \$\,60, están dispuestos a pagar solo \$\,50 por la deuda, porque anticipan que en una recesión no cobrarán todo: la valúan «con los ojos bien abiertos», con un rendimiento prometido de $60/50 - 1 = 20\,\%$. Sin costos de quiebra, Day vale lo mismo que una empresa idéntica sin riesgo de quiebra (\$\,68,16). Pero si en la recesión el flujo de \$\,50 se reduce a \$\,35 por honorarios de abogados y costos judiciales (\$\,15), el valor de la empresa cae a \$\,61,36 y los acreedores pagan solo \$\,43,18, con un rendimiento prometido cercano al 39\,\%: es un bono chatarra. Si Day emitiera esa deuda para pagar un dividendo, el dividendo sería de \$\,43,18 en lugar de \$\,50.
\end{ejemplo}

\begin{figure}[H]
  \centering
  \includegraphics[width=0.5\linewidth,height=0.42\textheight,keepaspectratio]{U3-1}
\end{figure}
\begin{figure}[H]
  \centering
  \includegraphics[width=0.5\linewidth,height=0.42\textheight,keepaspectratio]{U3-2}
\end{figure}
\begin{figure}[H]
  \centering
  \includegraphics[width=0.5\linewidth,height=0.42\textheight,keepaspectratio]{U3-3}
  \caption{Valor de la deuda y rendimiento prometido sin costos de quiebra y con ellos, caso Day Corporation.}
\end{figure}

\begin{definicion}{Riesgo de quiebra y costos de quiebra}
La posibilidad de quiebra reduce el valor de la empresa, pero no por el riesgo en sí, sino por los \textbf{costos} asociados a la quiebra, que consumen una porción del pastel. La quiebra no causa la caída del valor: es su resultado; lo que destruye valor es el costo de usar el mecanismo legal por el cual los acreedores toman el control. Como los acreedores pagan un precio justo cuando son realistas sobre la probabilidad y el costo de la quiebra, son los \textbf{accionistas} quienes soportan esos costos futuros.
\end{definicion}

\subsection{Descripción de los costos de las dificultades financieras}

Como el mismo efecto existe aunque se evite la quiebra legal, conviene hablar de \emph{costos de las dificultades financieras}. Hay dificultades cuando los compromisos con los acreedores no se cumplen o se cumplen con esfuerzo; su costo esperado depende de su probabilidad y de su magnitud. Se retroalimentan: una caída del resultado operativo genera dificultades, estas obligan a decidir bajo presión, y esas decisiones deterioran aún más el resultado.

\paragraph{Costos directos: legales y administrativos de la liquidación o reorganización.} Honorarios de abogados, contadores, peritos y costas judiciales, que se pagan antes que a los acreedores y solo si la quiebra ocurre. Son grandes en monto —la quiebra del condado de Orange en California costó más de 2,4 millones de dólares por mes; las de Enron y WorldCom, unos 1.000 y 600 millones—, pero pequeños en proporción al valor: los estudios los ubican entre el 1 y el 3\,\% del valor de mercado, con estimaciones de hasta el 2--10\,\% de los activos. Dependen también de la composición del activo: una empresa con muchos intangibles vale menos en la liquidación, porque marcas y patentes valen mientras se usan. Por economías de escala, pesan relativamente más en las empresas pequeñas.

\paragraph{Costos indirectos: el deterioro de la capacidad para hacer negocios.} Son los más importantes, difíciles de medir, y aparecen ante la sola posibilidad de la cesación de pagos:
\begin{itemize}
\item \textbf{pérdida de ventas}: nadie compra a quien quizá no preste servicio, garantía ni repuestos (clientes de Chrysler en los años setenta, los apostadores que evitaban un casino insolvente). Sin ventas no hay resultado operativo ni con qué pagar intereses o dividendos: es el costo mayor;
\item \textbf{subinversión}: para cumplir, la empresa elige proyectos de recupero rápido y menor rentabilidad, y posterga los mejores;
\item \textbf{restricción crediticia}: el deterioro de los indicadores que miran los bancos (endeudamiento, cobertura de intereses, rentabilidad) se traduce en tasas más altas, plazos cortos, más garantías o la pérdida del crédito;
\item \textbf{conflicto con proveedores}, que se niegan a operar o endurecen las condiciones, y \textbf{con empleados}, cuya productividad cae por temor a perder el empleo;
\item \textbf{pérdida del ahorro fiscal}, porque sin utilidades no hay impuesto que proteger.
\end{itemize}
Las estimaciones de los costos totales, directos más indirectos, van del 8--10\,\% del valor operativo para empresas hoy sanas a más del 20\,\% del valor para las que ya están en problemas (Andrade y Kaplan: entre 10 y 20\,\%). El litigio entre Texaco y Pennzoil es ilustrativo: el valor de Texaco cayó 4.100 millones y el de Pennzoil subió solo 682, una pérdida neta de unos 3.500 millones que los honorarios legales no explican; la causa probable son los costos indirectos (proveedores que exigían pago al contado, embarques cancelados, bancos que restringieron operaciones).

\paragraph{Costos de agencia.} Con deuda surgen conflictos entre accionistas y acreedores, que se agravan en las dificultades y llevan a los accionistas a estrategias egoístas que reducen el valor de toda la empresa.

\begin{definicion}{Estrategias egoístas de los accionistas}
\begin{enumerate}
\item \textbf{Incentivo a asumir grandes riesgos} (traslado del riesgo): cerca de la quiebra, los accionistas «juegan con el dinero de otros» y prefieren proyectos riesgosos aun con VPN menor; expropian valor a los acreedores (errores de comisión).
\item \textbf{Incentivo a subinvertir} (negarse a aportar capital): rechazan proyectos con VPN positivo cuando gran parte de su valor iría a los acreedores (errores de omisión).
\item \textbf{«Ordeñar» la empresa} (cobrar y huir): pagar dividendos extraordinarios en plena dificultad, dejando menos para los acreedores.
\end{enumerate}
A ellas se suman \textbf{ganar tiempo} —demorar la ejecución con cambios contables, falsas expectativas o recortes de mantenimiento e I+D— y el \textbf{cebo y cambio} —emitir poca deuda segura y reemplazarla de golpe por deuda riesgosa—.
\end{definicion}

\begin{ejemplo}[Las dos primeras estrategias]
\textbf{Riesgo.} Una empresa que debe \$\,100 elige entre dos proyectos con recesión y auge equiprobables. El de bajo riesgo da un valor de 100 o 200 (acciones 0 o 100; bonos 100 en ambos casos): valor esperado 150, acciones 50. El de alto riesgo da 50 o 240 (acciones 0 o 140; bonos 50 o 100): valor esperado 145, pero acciones 70. Sin deuda se elegiría el primero; con deuda, los accionistas eligen el segundo, porque ganan en el auge y la pérdida de la recesión la soportan los acreedores.

\textbf{Subinversión.} Otra empresa debe pagar \$\,4.000 y tendrá flujos de 5.000 en auge y 2.400 en recesión (acciones: 1.000 y 0; valor esperado 500). Un proyecto de \$\,1.000 aportado por los accionistas agrega 1.700 en ambos estados y evita la quiebra: los flujos pasan a 6.700 y 4.100, y las acciones, a 2.700 y 100 (valor esperado 1.400). Aunque el VPN es positivo, la participación de los accionistas sube 900 por una inversión de 1.000, porque en la recesión casi todo el beneficio va a los acreedores: el proyecto se rechaza.
\end{ejemplo}

Estas distorsiones existen solo si hay una probabilidad real de dificultades: no afectan a una empresa diversificada y de bajo riesgo, cuya deuda es casi libre de riesgo, ni a las empresas reguladas, pero sí a las pequeñas de industrias riesgosas. En la empresa sin deuda se aceptan siempre los proyectos de VPN positivo; con deuda, la política de inversión puede distorsionarse. Al final, los costos los pagan los accionistas: los acreedores racionales anticipan esas conductas y se protegen exigiendo tasas más altas, de modo que a estas empresas el crédito les resulta difícil y caro, y terminan con bajos niveles de deuda.

\subsection{¿Pueden reducirse los costos de la deuda?}

Los métodos disponibles pueden reducir los costos de la deuda, pero no eliminarlos.

\begin{definicion}{Cláusulas de protección (\emph{protective covenants})}
Son acuerdos incorporados al contrato de préstamo o a la escritura de emisión, que los accionistas aceptan para obtener tasas más bajas; violarlos puede provocar el incumplimiento. Las \textbf{negativas} limitan o prohíben acciones: pagar dividendos por encima de cierto monto, gravar activos a favor de otros prestamistas, fusionarse, vender o arrendar los activos principales sin autorización, emitir más deuda de largo plazo. Las \textbf{positivas} obligan a hacer algo: mantener un capital de trabajo mínimo, entregar estados financieros periódicos.
\end{definicion}

Cada cláusula responde a una conducta posible de los accionistas: las restricciones sobre los estados financieros (capital de trabajo, cobertura de intereses y patrimonio mínimos) desalientan las inversiones de alto riesgo cerca de la crisis; las restricciones a la disposición de activos (límites a dividendos y a la venta de activos, garantías e hipotecas) impiden transferirlos a los accionistas y subinvertir; las restricciones al intercambio de activos impiden aumentar el riesgo; y las restricciones a la dilución (límites al arrendamiento y a nueva deuda de igual o mayor prioridad) protegen los derechos existentes. En una muestra clásica de emisiones (Smith y Warner), el 91\,\% restringía la deuda adicional, el 39\,\% las fusiones, el 36\,\% la venta de activos y el 23\,\% los dividendos.

Frente al conflicto, los accionistas tienen tres opciones: no emitir deuda, carísimo por la pérdida de la ventaja fiscal; emitirla sin cláusulas, a cambio de tasas altas; o incluir cláusulas claras que protejan a los acreedores a bajo costo para los accionistas. Por eso las cláusulas, aunque reduzcan la flexibilidad, aumentan el valor de la empresa y suelen ser la solución más barata; los accionistas deberían favorecer toda cláusula razonable. Otras dos vías son la \textbf{recompra de deuda} antes de la quiebra y la \textbf{consolidación de la deuda}: si uno o pocos prestamistas concentran la deuda, o si los acreedores también son accionistas —práctica habitual de los grandes bancos japoneses—, se reducen los costos de negociar en una crisis, lo que permite razones de deuda más altas.

\subsection{Integración de los efectos fiscales y de los costos de las dificultades financieras}

Según MM con impuestos, todas las empresas deberían endeudarse al máximo, lo que no ocurre. La explicación es que los costos de las dificultades financieras reducen el valor de la empresa apalancada.

\begin{definicion}{Teoría estática del equilibrio (\emph{trade-off})}
La estructura de capital es un equilibrio entre los beneficios fiscales de la deuda y los costos de las dificultades financieras. Con poca deuda, la probabilidad de crisis es baja y predominan los beneficios; a medida que se agrega deuda, el valor presente de esos costos crece a tasa creciente. El \textbf{óptimo} $B^*$ está donde el valor presente del ahorro fiscal de un peso más de deuda iguala el aumento del valor presente de los costos esperados de dificultades: allí el valor de la empresa es máximo y el WACC, mínimo. Cada empresa tiene así un nivel objetivo de deuda (su \emph{capacidad de endeudamiento}), aunque no exista una fórmula exacta para calcularlo.
\end{definicion}

\begin{formulaBox}[Valor de la empresa con impuestos, dificultades y agencia]
\[
V_L = V_U + t_c\,D - VP(\text{costos de dificultades financieras}) - VP(\text{costos de agencia})
\]
\end{formulaBox}

\begin{figure}[H]
  \centering
  \includegraphics[width=0.5\linewidth,height=0.42\textheight,keepaspectratio]{U3-4}
\end{figure}
\begin{figure}[H]
  \centering
  \includegraphics[width=0.5\linewidth,height=0.42\textheight,keepaspectratio]{U3-5}
  \caption{Cantidad óptima de deuda y valor de la empresa: el valor sigue una curva en forma de $\cap$ y el WACC alcanza su mínimo en $B^*$ (figura 17.1).}
\end{figure}

\begin{figure}[H]
  \centering
  \includegraphics[width=0.9\linewidth,height=0.42\textheight,keepaspectratio]{screenshot008}
  \caption{La estructura de capital óptima: el valor presente del ahorro fiscal frente al valor presente de los costos de las dificultades financieras.}
\end{figure}

La teoría admite objetivos distintos para cada empresa, evita predicciones extremas, justifica razones moderadas y explica muchas diferencias entre sectores; no explica, en cambio, por qué algunas de las empresas más exitosas prosperan con muy poca deuda.

\paragraph{Nuevamente el pastel.} Los flujos de la empresa tienen cuatro reclamantes: accionistas ($S$), acreedores ($B$), gobierno ($G$) y, en las dificultades, abogados y otros ($L$). A ellos se suman los pleitos potenciales ($LS$): cualquier persona tiene un derecho económico eventual sobre una gran empresa, porque podría demandarla.

\begin{figure}[H]
  \centering
  \includegraphics[width=0.5\linewidth,height=0.42\textheight,keepaspectratio]{U3-6}
\end{figure}
\begin{figure}[H]
  \centering
  \includegraphics[width=0.5\linewidth,height=0.42\textheight,keepaspectratio]{U3-7}
\end{figure}
\begin{figure}[H]
  \centering
  \includegraphics[width=0.5\linewidth,height=0.42\textheight,keepaspectratio]{U3-8}
\end{figure}
\begin{figure}[H]
  \centering
  \includegraphics[width=0.5\linewidth,height=0.42\textheight,keepaspectratio]{U3-9}
  \caption{El modelo del pastel con factores del mundo real: $V_T = S + B + G + L$ (figura 17.2).}
\end{figure}

\begin{definicion}{Derechos negociables y no negociables}
Los derechos de accionistas y acreedores son \textbf{negociables} ($V_M$): se compran y venden en los mercados, y quienes los adquieren pagan a la empresa por ellos. Los del fisco y de los litigantes son \textbf{no negociables} ($V_N$): nadie paga a la empresa por el privilegio de cobrarle impuestos u honorarios.
\end{definicion}

\begin{formulaBox}[Valor total de la empresa]
\[
V_T = S + B + G + L = V_M + V_N
\]
La estructura de capital no cambia $V_T$, que depende de los flujos de la empresa, pero sí su reparto: todo aumento de $V_M$ implica una disminución igual de $V_N$. Los administradores racionales eligen la estructura que maximiza $V_M$, es decir, que minimiza impuestos, costos de quiebra y demás derechos no negociables.
\end{formulaBox}

\begin{figure}[H]
  \centering
  \includegraphics[width=0.5\linewidth,height=0.42\textheight,keepaspectratio]{U3-10}
\end{figure}
\begin{figure}[H]
  \centering
  \includegraphics[width=0.5\linewidth,height=0.42\textheight,keepaspectratio]{U3-11}
\end{figure}
\begin{figure}[H]
  \centering
  \includegraphics[width=0.5\linewidth,height=0.42\textheight,keepaspectratio]{U3-12}
  \caption{Derechos negociables y no negociables: la estructura de capital solo reparte el valor total.}
\end{figure}

\paragraph{Cómo estimar los costos de insolvencia.} En la teoría de MM la empresa siempre gana y aprovecha todo el escudo; en la realidad las ganancias no son permanentes y un endeudamiento alto puede reducir el propio resultado operativo. Por eso conviene una razón de deuda que capture la ventaja fiscal, deje un margen ante coyunturas adversas y se apoye en un indicador de cobertura que tenga en cuenta la variabilidad y la tendencia del resultado operativo.

\begin{formulaBox}[Índice normal de solvencia]
\[
\begin{gathered}
\text{Resultado estándar} = \frac{1}{n}\sum_{j=1}^{n}\frac{\text{resultado}_j}{\text{activo operativo promedio}_j}\times\text{Activo operativo}_n
\\[4pt]
INS = \frac{\text{Resultado estándar}}{\text{Intereses}}
\end{gathered}
\]
El resultado estándar es el rendimiento promedio del activo multiplicado por el activo actual.
\end{formulaBox}

\begin{definicion}{Modelo de estructura de capital óptima (ECO)}
Intercambia los beneficios del endeudamiento con la calificación (\emph{rating}) de la deuda, ligada a la cobertura normalizada de intereses. Supone que el ahorro fiscal depende de la capacidad de los activos de generar renta, por lo que se descuenta a $k_u$; que la deuda tiene riesgo y su tasa $k_d$ depende del rating; que el mayor riesgo de insolvencia se reparte entre accionistas y acreedores, por lo que $k_e$ se calcula como $k_d$ más una prima aproximadamente fija; y que los costos de insolvencia modifican el riesgo de los activos ($\beta_u$) y su rendimiento exigido ($k_u$). Con él, el valor presente de las dificultades financieras es $V_{ECO} - V_{MM}$. Un método afín, el \emph{WACC approach}, tabula el valor de la empresa para distintos niveles de deuda recalculando el beta y el rendimiento exigido en cada uno.
\end{definicion}

\subsection{Emisión de señales}

En el equilibrio estático, el subsidio marginal de la deuda supera su costo marginal con poca deuda y ocurre lo contrario con mucha. Una empresa que espera utilidades bajas se endeuda poco, porque una pequeña deducción le alcanza y más deuda elevaría los costos de quiebra; una más exitosa puede endeudarse más con poco aumento del riesgo. Las empresas racionales, por lo tanto, se endeudan más cuando esperan que sus utilidades crezcan, y los inversores interpretan la deuda como una \textbf{señal} del valor de la empresa.

Un gerente podría endeudarse de más para engañar al mercado y subir hoy el precio de la acción; pero cuando el engaño se descubre, el precio cae por debajo del que habría tenido, porque la deuda quedó por encima del óptimo. Esos costos impiden que las empresas menos valiosas imiten a las más valiosas, de modo que la deuda sigue siendo una señal creíble. La evidencia lo confirma: el precio de las acciones sube con fuerza al anunciarse una oferta de canje que aumenta el apalancamiento y cae cuando el canje lo reduce.

\begin{figure}[H]
  \centering
  \includegraphics[width=0.5\linewidth,height=0.42\textheight,keepaspectratio]{U3-13}
  \caption{Rendimientos de las acciones al anunciarse ofertas de canje que aumentan o reducen el apalancamiento (figura 17.3).}
\end{figure}

\subsection{Evasión de obligaciones, prerrogativas y malas inversiones: el costo de agencia del capital}

Adam Smith ya lo había advertido: quien administra dinero ajeno no lo cuida con la vigilancia con que cuida el propio. Un individuo trabaja más si es dueño que si es empleado, y más cuanto mayor es su participación. Por eso, cuando un propietario-gerente emite capital nuevo y diluye su participación, tiende a trabajar menos, tomar más prerrogativas (\emph{perquisites}) y emprender proyectos de VPN negativo —los sueldos gerenciales suelen crecer con el tamaño—; la pérdida por malos proyectos es mucho mayor que la causada por la holgazanería o los gastos excesivos. El problema pesa en una empresa chica que hace una gran emisión, no en una gran corporación cuyo gerente ya tiene una porción ínfima. Lo soporta el propietario, porque los nuevos inversores, si invierten con los ojos abiertos, pagan menos por las acciones; puede atenuarlo aceptando vigilancia. Las \textbf{compras apalancadas} (LBO), en las que los gerentes compran la empresa y concentran la propiedad, reducen mucho estos costos.

\begin{definicion}{Modelo estático ampliado}
Al sustituir capital por deuda, el valor de la empresa cambia por tres efectos: más escudo fiscal, menos costos de agencia del capital y más costos de dificultades financieras (incluidos los de agencia de la deuda). Por eso la razón óptima de deuda es mayor con costos de agencia del capital que sin ellos, aunque nunca llega al 100\,\% por el peso de los costos de las dificultades.
\end{definicion}

\begin{definicion}{Hipótesis del flujo de efectivo libre (Jensen, 1986)}
Un gerente con poca participación tiene el \emph{motivo} para derrochar; el flujo de efectivo libre le da la \emph{oportunidad}. Por eso cabe esperar más actividades derrochadoras —como malas adquisiciones— en empresas que generan grandes flujos. Dividendos y servicio de la deuda reducen ese flujo; la deuda lo hace con más fuerza, porque dejar de pagarla lleva a la quiebra, mientras que recortar un dividendo no. Pasar de capital a deuda aumenta, entonces, el valor de la empresa.
\end{definicion}

\subsection{Teoría de selección jerárquica}

El equilibrio estático omite una consideración: la oportunidad. Un gerente que cree que su acción, que cotiza a \$\,50, vale \$\,60, no emitirá acciones, porque regalaría valor a los nuevos accionistas; emitirá bonos, cuyo valor depende sobre todo de una tasa pública. Si la acción cotiza a \$\,70, en cambio, querrá emitir, porque ganaría \$\,10 para sus accionistas actuales; ganar eso eliminando costos de agencia o de quiebra llevaría años.

\begin{definicion}{Información asimétrica}
Los gerentes saben más que los inversores sobre las perspectivas de la empresa, porque trabajan en ella todos los días, y actúan en interés de los accionistas actuales. Si ambos supuestos se cumplen, emitirán acciones cuando las crean sobrevaluadas y deuda cuando las crean subvaluadas.
\end{definicion}

Pero los inversores lo saben: interpretan una emisión de acciones como señal de sobrevaluación y una de deuda como señal de subvaluación, y el precio cae al anunciarse una emisión de acciones. Es un juego de póquer en el que solo las empresas más sobrevaluadas tendrían incentivo para emitir acciones; una moderadamente sobrevaluada sería confundida con ellas. Resultado extremo: nadie emitiría acciones, tan extremo como la indiferencia de MM sin impuestos o el endeudamiento total de MM con impuestos. En la práctica, el gerente también pondera impuestos, costos de dificultades y de agencia: emite deuda hasta que la crisis se vuelve una posibilidad real y recién entonces considera las acciones. Por eso muchas empresas mantienen una \emph{reserva de endeudamiento}, una razón de deuda inferior a la óptima para endeudarse cuando aparezca una buena oportunidad.

\begin{definicion}{Teoría de la selección jerárquica (\emph{pecking order})}
Las inversiones se financian primero con fondos internos, luego con deuda y solo al final con nuevas acciones. Sus dos reglas:
\begin{enumerate}
\item \textbf{usar financiamiento interno}, porque aun la deuda riesgosa despierta el escepticismo de los inversores, que se evita si no se recurre a ellos;
\item \textbf{emitir primero los títulos más seguros}: deuda directa, luego híbridos como los convertibles y, por último, acciones.
\end{enumerate}
Además, las empresas adaptan sus metas de pago de dividendos a sus oportunidades de inversión; como los dividendos son rígidos y la rentabilidad fluctúa, cuando el flujo interno sobra cancelan deuda o compran valores, y cuando falta usan primero su efectivo.
\end{definicion}

Sus implicaciones difieren de las del equilibrio:
\begin{itemize}
\item \textbf{No hay un apalancamiento meta.} En el equilibrio, el óptimo está donde el beneficio marginal de la deuda iguala su costo marginal; en la jerarquía, el apalancamiento resulta de las necesidades de fondos y de la casualidad de los proyectos disponibles, porque hay dos clases de capital, el interno y el externo.
\item \textbf{Las empresas rentables se endeudan menos}, porque generan fondos propios; el equilibrio predice lo contrario: más flujo significa más capacidad de endeudarse y más escudo que aprovechar.
\item \textbf{A las empresas les gusta la holgura financiera}: acumulan efectivo, valores negociables, activos vendibles y acceso rápido al crédito para no depender del mercado cuando surja un proyecto. Pero esa holgura tiene un límite, porque el exceso de efectivo agrava los problemas de agencia.
\end{itemize}

\subsection{El crecimiento y la razón de deuda a capital}

Algunos críticos señalan que, como los costos de quiebra parecen mucho menores que el subsidio fiscal, el equilibrio implicaría razones de deuda cercanas al 100\,\%. La jerarquía es más coherente con la realidad, y hay otro argumento: según Berens y Cuny, el crecimiento exige financiamiento con capital aun con costos de quiebra bajos.

\begin{ejemplo}[Crecimiento y deuda óptima (tasa del 10\,\%)]
Una empresa sin crecimiento con una UAII perpetua de \$\,100 puede emitir \$\,1.000 de deuda: los intereses absorben toda la utilidad, no paga impuestos y su razón deuda/valor óptima es del 100\,\%, porque más deuda provocaría incumplimiento. Si la UAII crece al 5\,\% anual, para seguir sin pagar impuestos la deuda debe crecer también al 5\,\% (1.000, 1.050, 1.102,50…), y la nueva deuda puede pagarse como dividendo. El valor de la empresa es $100/(0{,}10-0{,}05) = \$\,2.000$; el capital vale \$\,1.000, igual al valor de los dividendos crecientes $50/(0{,}10-0{,}05)$, y la razón deuda/valor es del 50\,\%.
\end{ejemplo}

\begin{figure}[H]
  \centering
  \includegraphics[width=0.5\linewidth,height=0.42\textheight,keepaspectratio]{U3-14}
\end{figure}
\begin{figure}[H]
  \centering
  \includegraphics[width=0.5\linewidth,height=0.42\textheight,keepaspectratio]{U3-16}
\end{figure}
\begin{figure}[H]
  \centering
  \includegraphics[width=0.5\linewidth,height=0.42\textheight,keepaspectratio]{U3-15}
  \caption{Valor de la empresa con crecimiento, valor del capital y valuación equivalente como $V_U$ más el valor presente del escudo fiscal.}
\end{figure}

La clave es que, para no pagar impuestos, alcanza con que el interés de cada año iguale la utilidad de ese año: el crecimiento aumenta el valor de la empresa, pero no la deuda necesaria para proteger la utilidad actual, de modo que aumenta el valor del capital y la razón deuda/valor que protege toda la utilidad queda por debajo del 100\,\%. La conclusión vale también con inflación: el financiamiento 100\,\% con deuda es subóptimo, y las empresas de alto crecimiento deberían tener razones de deuda menores.

\subsection{Impuestos personales}

Las empresas no son las únicas que pagan impuestos: los inversores tributan por dividendos e intereses. De cada peso de utilidad antes de impuestos, el accionista de una empresa sin deuda recibe $(1-t_c)(1-t_S)$; si la empresa se financia con deuda, el peso entero se paga como interés y el acreedor recibe $1-t_B$.

\begin{figure}[H]
  \centering
  \includegraphics[width=0.5\linewidth,height=0.42\textheight,keepaspectratio]{U3-17}
\end{figure}
\begin{figure}[H]
  \centering
  \includegraphics[width=0.5\linewidth,height=0.42\textheight,keepaspectratio]{U3-18}
\end{figure}
\begin{figure}[H]
  \centering
  \includegraphics[width=0.5\linewidth,height=0.42\textheight,keepaspectratio]{U3-19}
\end{figure}
\begin{figure}[H]
  \centering
  \includegraphics[width=0.5\linewidth,height=0.42\textheight,keepaspectratio]{U3-20}
  \caption{Impuestos corporativos y personales sobre un peso de utilidad pagado a accionistas o a acreedores.}
\end{figure}

La empresa debe elegir la estructura que ponga más efectivo en manos de sus inversores, es decir, la que minimice el valor actual de \emph{todos} los impuestos, corporativos y personales. Tres casos lo aclaran:
\begin{enumerate}
\item si $t_S = t_B$, los acreedores reciben más que los accionistas y conviene la deuda: la renta del capital tributa dos veces y la de la deuda una sola (la ausencia de impuestos personales es un caso particular);
\item la empresa es indiferente cuando ambos flujos se igualan;
\item en el mundo real, con $t_c = t_B = 35\,\%$ y $t_S = 15\,\%$ (Estados Unidos, 2005), el accionista recibe $0{,}65\times 0{,}85 \approx 0{,}55$ y el acreedor 0,65. Como las recompras tributan solo sobre la ganancia de capital, la tasa efectiva sobre las distribuciones es menor, por ejemplo 10\,\%; aun así el accionista recibe 0,59 y la deuda conserva su ventaja. En los años setenta, con intereses gravados hasta el 70\,\% y ganancias de capital a tasas mucho menores, las acciones probablemente tenían ventaja.
\end{enumerate}

\begin{formulaBox}[Impuestos personales y equilibrio de Miller]
\[
\begin{gathered}
\text{Indiferencia (equilibrio de Miller):}\quad (1-t_c)(1-t_S) = 1-t_B
\\[4pt]
\text{Ventaja relativa de la deuda} = \frac{1-t_B}{(1-t_S)(1-t_c)}
\\[4pt]
V_U = \frac{EBIT\,(1-t_c)(1-t_S)}{k_u}
\qquad\qquad
V_L = V_U + \left[1 - \frac{(1-t_c)(1-t_S)}{1-t_B}\right] B
\end{gathered}
\]
$t_B$: impuesto personal sobre los intereses; $t_S$: impuesto personal efectivo sobre las distribuciones a los accionistas. Si $t_S = t_B$, el término entre corchetes se reduce a $t_c$ y vuelve la Proposición I con impuestos.
\end{formulaBox}

% Figura opcional (el material la menciona como ilegible en el OCR):
% se incluye solo si su archivo está en la carpeta.
\IfFileExists{ross process-13.png}{%
\begin{figure}[H]
  \centering
  \includegraphics[width=0.5\linewidth,height=0.42\textheight,keepaspectratio]{ross process-13}
  \caption{Modelo de Miller: valor de la empresa apalancada con impuestos corporativos y personales.}
\end{figure}}{}

\begin{ejemplo}[Las empresas A y B con impuestos personales]
Con $t_c = 40\,\%$, un impuesto personal del 20\,\% sobre la renta de las acciones y del 25\,\% sobre los intereses: $V_U = 20\times 0{,}60\times 0{,}80/0{,}20 = 48$ y $V_L = 48 + [1 - 0{,}48/0{,}75]\times 50 = 48 + 0{,}36\times 50 = 66$. La ganancia por endeudarse (18) es menor que la de MM con impuestos corporativos (20).
\end{ejemplo}

\begin{figure}[H]
  \centering
  \includegraphics[width=0.5\linewidth,height=0.42\textheight,keepaspectratio]{screenshot007}
  \caption{El equilibrio de Miller: el valor de la empresa según los impuestos corporativos y personales.}
\end{figure}

Si la renta del capital tributa menos que los intereses, la ventaja fiscal de la deuda es menor que sin impuestos personales, y el monto óptimo de deuda también. Si el impuesto personal sobre los intereses superara la alícuota corporativa (con la renta del capital exenta), endeudarse destruiría valor, porque la empresa tendría que «sobornar» a los acreedores con tasas más altas. Hay dos casos especiales: cuando deuda y capital tributan igual a nivel personal, los impuestos personales salen del cálculo; y cuando los impuestos corporativos y personales se compensan, la política de endeudamiento es irrelevante, lo que exige una alícuota corporativa menor que la personal y una tasa baja sobre la renta del capital.

El \textbf{modelo de Miller} (1977) aplica esta idea con un impuesto personal progresivo: en conjunto, el sector corporativo emite la deuda justa para que la mantengan los inversores con tasas iguales o menores que la corporativa, mientras los de tasas más altas tienen acciones; para cada empresa la decisión es indiferente por la competencia, y el rendimiento de los bonos supera al de acciones de riesgo comparable. Dos críticas lo limitan: con las tasas reales, la desventaja personal de la deuda reduce pero no elimina el incentivo corporativo (Graham), por lo que el modelo implicaría financiarse solo con deuda; y los intereses solo se deducen hasta el monto de las ganancias, de modo que el beneficio esperado de la deuda es menor, las empresas ofrecerían menos deuda y la primera unidad de deuda vale más que la última.

\subsection{Cómo establecen las empresas su estructura de capital}

A diferencia del VPN, la teoría de la estructura de capital no ofrece una fórmula exacta, por lo que se recurre a regularidades empíricas:
\begin{enumerate}
\item \textbf{La mayoría de las empresas tiene razones de deuda bajas} y usa menos deuda que capital; no se endeudan hasta agotar el escudo (los impuestos corporativos estadounidenses superaron los 460.000 millones de dólares en 2007), lo que revela límites como los costos de las dificultades.
\item \textbf{Muchas empresas no usan deuda}: en un centenar de empresas de la Bolsa de Nueva York sin deuda de largo plazo, los gerentes tenían alta participación accionaria, había más presencia familiar y más efectivo; gerentes poco diversificados rehúyen el riesgo adicional del apalancamiento.
\item \textbf{Las estructuras difieren entre industrias} de manera persistente: son muy bajas en sectores de alto crecimiento con grandes oportunidades de inversión (farmacéutico, electrónico) y altas en los de activos tangibles (bienes raíces).
\item \textbf{La mayoría fija razones meta}: solo el 19\,\% de 392 directores financieros encuestados no las usaba, aunque su rigidez varía y no se reajustan ante cada cambio de precio de la acción.
\end{enumerate}

\begin{figure}[H]
  \centering
  \includegraphics[width=0.5\linewidth,height=0.42\textheight,keepaspectratio]{U3-21}
  \caption{Razones de deuda a valor total de empresas de distintos países (figura 17.4).}
\end{figure}

\begin{figure}[H]
  \centering
  \includegraphics[width=0.5\linewidth,height=0.42\textheight,keepaspectratio]{U3-22}
  \caption{Uso de razones de deuda a capital fijadas como meta según una encuesta a directores financieros (figura 17.5).}
\end{figure}

\begin{figure}[H]
  \centering
  \includegraphics[width=0.5\linewidth,height=0.42\textheight,keepaspectratio]{U3-23}
  \caption{Razones de estructura de capital de empresas no financieras por industria (tabla 17.3).}
\end{figure}

Tres factores orientan la razón meta:
\begin{itemize}
\item \textbf{Impuestos}: como los intereses solo se deducen hasta el monto de las utilidades, las empresas muy rentables tienden a metas más altas; endeudarse crea valor si la alícuota corporativa supera la de los acreedores.
\item \textbf{Tipo de activos}: los tangibles conservan valor de reventa; la I+D lo pierde casi todo en una crisis; por eso las empresas con activos tangibles pueden tener metas de deuda más altas.
\item \textbf{Incertidumbre del resultado operativo}: las empresas con resultados inciertos (como los laboratorios) pueden caer en dificultades aun sin deuda y deben financiarse sobre todo con capital; las reguladas de servicios públicos, de resultados estables, usan mucha deuda.
\end{itemize}

Los estudios agregan cuatro determinantes, coherentes con el equilibrio y con la jerarquía: las empresas \textbf{grandes} se endeudan más; las que tienen más \textbf{activos tangibles}, también; las más \textbf{rentables}, menos; y las de mayor relación \textbf{valor de mercado/valor contable}, menos. Los anuncios que aumentan el apalancamiento elevan el precio de la acción y los que lo reducen lo deprimen, lo que es compatible tanto con la señalización del equilibrio (el gerente se endeuda si cree que bajó la probabilidad de quiebra) como con la jerarquía (convierte deuda en acciones cuando las cree sobrevaluadas). Como estas guías son imprecisas, muchas empresas se apoyan en los promedios de su industria: al menos evitan alejarse de las prácticas aceptadas, y las empresas existentes son, después de todo, las sobrevivientes.

%% =====================================================================
%%  UNIDAD 3 (parte C)
%% =====================================================================
\section{El costo de capital}\label{sec:u3-ccpp}

\subsection{Costo de capital promedio ponderado}

El costo de capital es la tasa mínima que la empresa debe ganar sobre sus inversiones para aumentar su valor. Como se financia con varias fuentes, la tasa relevante combina sus costos.

\begin{definicion}{Costo de capital promedio ponderado (CCPP)}
Es el costo futuro promedio esperado de los fondos de largo plazo de la empresa, $k_a$. Se calcula ponderando el costo de cada tipo específico de capital (deuda de largo plazo, acciones preferidas, ganancias retenidas o acciones comunes nuevas) por su proporción en la estructura de capital. Es la generalización del WACC \vease{sec:u3-conceptos}.
\end{definicion}

\begin{formulaBox}[Costo de capital promedio ponderado]
\[
k_a = w_i\,k_i + w_p\,k_p + w_s\,k_{r\ \text{o}\ n}
\qquad\qquad
w_i + w_p + w_s = 1
\]
$w_i$, $w_p$ y $w_s$ son las proporciones de deuda de largo plazo, acciones preferidas y capital común; $k_i$ es el costo de la deuda después de impuestos, $k_p$ el de las preferidas, $k_r$ el de las ganancias retenidas y $k_n$ el de las acciones comunes nuevas.
\end{formulaBox}

Tres aspectos del cálculo: conviene expresar las ponderaciones en forma decimal y los costos en porcentaje; las ponderaciones deben ser positivas y sumar 1, de modo que se incluyan todas las fuentes; y el peso del capital común se multiplica por $k_r$ o por $k_n$, según se financie con ganancias retenidas o con nuevas acciones. Los costos de cada componente se desarrollan en el capítulo~\ref{cap:u5} \vease{sec:u5-costocap}.

\begin{ejemplo}[Duchess Corporation]
Los costos de sus fuentes son: deuda 5,6\,\%, preferidas 10,6\,\%, ganancias retenidas 13,0\,\% y acciones nuevas 14,0\,\%; sus ponderaciones, 40\,\%, 10\,\% y 50\,\%. Como espera contar con \$\,300.000 de ganancias retenidas, usa $k_r$:
\begin{center}
\begin{tabular}{@{}lccc@{}}
\toprule
Fuente & Ponderación & Costo & Costo ponderado\\
\midrule
Deuda de largo plazo & 0,40 & 5,6\,\% & 2,2\,\%\\
Acciones preferidas & 0,10 & 10,6\,\% & 1,1\,\%\\
Capital común & 0,50 & 13,0\,\% & 6,5\,\%\\
\midrule
Total & 1,00 & & \textbf{CCPP = 9,8\,\%}\\
\bottomrule
\end{tabular}
\end{center}
Si el riesgo no cambia, debería aceptar los proyectos que rindan más del 9,8\,\%.
\end{ejemplo}

\begin{figure}[H]
  \centering
  \includegraphics[width=0.5\linewidth,height=0.42\textheight,keepaspectratio]{U3-69}
\end{figure}
\begin{figure}[H]
  \centering
  \includegraphics[width=0.5\linewidth,height=0.42\textheight,keepaspectratio]{U3-70}
  \caption{Cálculo del costo de capital promedio ponderado, caso Duchess Corporation.}
\end{figure}

\subsection{Esquemas de ponderación}

\begin{definicion}{Ponderaciones del costo de capital}
\begin{itemize}
\item \textbf{De valor en libros}: miden la proporción de cada fuente con valores contables.
\item \textbf{De valor de mercado}: la miden con valores de mercado.
\item \textbf{Históricas}: se basan en las proporciones reales, pasadas o presentes, de la estructura.
\item \textbf{Objetivo}: se basan en las proporciones deseadas, es decir, en la estructura «óptima» que la empresa quiere alcanzar.
\end{itemize}
\end{definicion}

Las ponderaciones de mercado son preferibles: se acercan a lo que realmente se obtendría por la venta de los títulos, son coherentes con costos calculados a precios de mercado y con flujos de inversión medidos en términos de mercado. Dada la naturaleza aproximada del cálculo, la elección quizá no sea decisiva, pero en teoría el esquema preferido es el de \textbf{proporciones objetivo a valor de mercado}.

\begin{ejemplo}[El promedio ponderado en las finanzas personales]
Chuck Solis debe \$\,26.000 al 9,6\,\%, \$\,9.000 al 10,6\,\% y \$\,45.000 al 7,4\,\%, y le ofrecen \$\,80.000 al 9,2\,\% para cancelar todo. El costo promedio de su deuda actual es $0{,}325\times 9{,}6\,\% + 0{,}1125\times 10{,}6\,\% + 0{,}5625\times 7{,}4\,\% = 8{,}47\,\% \approx 8{,}5\,\%$, menor que el 9,2\,\% ofrecido: le conviene no cambiar nada.
\end{ejemplo}

\begin{figure}[H]
  \centering
  \includegraphics[width=0.5\linewidth,height=0.42\textheight,keepaspectratio]{U3-71}
  \caption{Resumen del capítulo sobre el costo de capital.}
\end{figure}

\paragraph{Enfoque en el valor.} El costo de capital es la tasa principal de las decisiones de inversión, y su exactitud determina su calidad: subestimarlo hace atractivos proyectos malos, y sobreestimarlo hace rechazar proyectos buenos. Estimarlo bien aumenta la probabilidad de que las decisiones de largo plazo sean coherentes con la meta de maximizar la riqueza de los dueños. El costo relevante para nuevas inversiones es el costo marginal de los fondos que se necesita obtener \vease{sec:u5-costocap}.

\subsection{El WACC después de impuestos y su uso práctico}

Con impuestos, las decisiones de inversión y de financiamiento ya no pueden separarse como sugería MM. El WACC se calcula con las características actuales de la empresa, pero se usa para descontar flujos futuros; eso supone dos cosas: que el riesgo de negocio del proyecto es igual al del resto de los activos y se mantiene durante su vida, y que el proyecto conserva la misma proporción de deuda sobre valor que la empresa, también constante.

\begin{formulaBox}[WACC después de impuestos y costo de oportunidad del capital]
\[
\begin{gathered}
WACC = r_D\,(1-T_c)\,\frac{D}{V} + r_E\,\frac{E}{V}
\qquad\qquad
r = r_D\,\frac{D}{V} + r_E\,\frac{E}{V}
\qquad\qquad
r_E = r + (r - r_D)\,\frac{D}{E}
\end{gathered}
\]
$r$ es el costo de oportunidad del capital de la empresa (el WACC sin impuestos, que no depende de la estructura); con más fuentes, como las preferidas, se agrega un término por cada una, ponderado por su valor de mercado.
\end{formulaBox}

Reglas prácticas: la deuda de corto plazo suele omitirse, lo que es teóricamente incorrecto pero poco grave si es transitoria o está compensada con efectivo y valores negociables; los demás pasivos corrientes se restan de los activos corrientes y la diferencia figura como capital de trabajo neto en el activo; la suma del financiamiento de largo plazo es la \emph{capitalización total}. Al aumentar la deuda el costo del capital propio sube, pero el WACC baja, y no porque se use deuda «barata» en lugar de capital «caro», sino por el escudo fiscal de los intereses. Si cambia la estructura, el WACC se recalcula en tres pasos: se desapalanca para obtener el costo de oportunidad $r$; se estiman el nuevo costo de la deuda y el nuevo costo del capital propio con la nueva razón $D/E$; y se vuelve a ponderar.

\paragraph{Valuación de la empresa y valor de horizonte.} Para valuar la empresa entera se la trata como un gran proyecto. Si se descuenta al WACC, los flujos deben calcularse como si la empresa se financiara solo con capital propio: se usa el \emph{free cash flow}, que difiere de la utilidad neta en tres aspectos: se calcula antes de intereses, suma las depreciaciones y otros gastos no monetarios, y resta los gastos de capital y la inversión en capital de trabajo, que no son gastos contables \vease{sec:u1-fcff}. Como las empresas son prácticamente inmortales, se proyecta un horizonte de mediano plazo y se agrega un valor terminal, que suele concentrar la mayor parte del valor y debe estimarse con cuidado.

\begin{formulaBox}[Valuación por FCF descontado al WACC]
\[
VP = \frac{FCF_1}{1+WACC} + \dots + \frac{FCF_H}{(1+WACC)^H} + \frac{VP_H}{(1+WACC)^H}
\qquad\qquad
VP_H = \frac{FCF_{H+1}}{WACC - g}
\]
El resultado valúa los activos y las operaciones; para obtener el valor del capital propio se resta la deuda.
\end{formulaBox}

El punto clave es el \textbf{reequilibrio}: el WACC calculado con la estructura actual supone que la empresa reajustará su financiamiento para mantener la misma razón de deuda sobre valor de mercado.

%% ---------------------------------------------------------------------
\section{Valoración y presupuesto de capital para la empresa apalancada}\label{sec:u3-metodos}

Cuando un proyecto o una empresa se financia en parte con deuda, su valor puede calcularse de varias maneras, que difieren en el flujo que se descuenta y en la tasa que se usa.

\subsection{El enfoque del valor presente ajustado (APV)}

\begin{definicion}{Valor presente ajustado (APV)}
Es el valor del proyecto como si se financiara solo con capital propio (caso base) más el valor presente de los efectos secundarios del financiamiento: $APV = VPN + VPNF$. Los efectos son cuatro: el subsidio fiscal de la deuda (positivo, el más importante), los costos de emitir títulos (negativos), los costos de las dificultades financieras (negativos) y los subsidios al financiamiento otorgados por un proveedor o por el gobierno (positivos).
\end{definicion}

\begin{formulaBox}[Valor presente ajustado]
\[
APV = \sum_{t=1}^{\infty}\frac{UCF_t}{(1+r_0)^t} + \text{efectos adicionales de la deuda} - \text{inversión inicial}
\]
$UCF_t$: flujo del proyecto para los accionistas de una empresa no apalancada; $r_0$: costo de capital del proyecto sin deuda.
\end{formulaBox}

El APV es útil cuando la deuda está atada al valor contable o debe cancelarse según un cronograma fijo, funciona bien en las compras apalancadas (LBO) y rinde más cuando los efectos del financiamiento son numerosos e importantes, como en las inversiones internacionales.

\subsection{El enfoque del flujo a capital (FTE) y el método WACC}

\begin{formulaBox}[Flujo a capital y método WACC]
\[
\begin{gathered}
FTE = \sum_{t=1}^{\infty}\frac{LCF_t}{(1+r_S)^t} - (\text{inversión inicial} - \text{monto prestado})
\\[4pt]
VPN_{WACC} = \sum_{t=1}^{\infty}\frac{UCF_t}{(1+r_{WACC})^t} - \text{inversión inicial}
\end{gathered}
\]
$LCF_t$: flujo del proyecto para los accionistas de la empresa apalancada, neto de intereses; $r_S$: costo del capital propio apalancado.
\end{formulaBox}

\subsection{Comparación de los enfoques APV, FTE y WACC}

El APV y el WACC ponen el flujo no apalancado en el numerador y ajustan por el beneficio fiscal de la deuda: el APV lo suma como un término separado; el WACC lo incorpora de modo más sutil, reduciendo la tasa por debajo de $r_0$. El FTE valúa otra entidad: resta de la inversión solo el aporte propio y trabaja con flujos netos de intereses.

\begin{definicion}{Regla de elección del método}
Usar \textbf{WACC} o \textbf{FTE} cuando se mantiene durante la vida del proyecto una razón meta de deuda sobre valor; usar \textbf{APV} cuando lo conocido es el \emph{nivel} de deuda en cada período. En el mundo real las empresas suelen fijar una razón de deuda sobre valor de mercado —los bancos prestan a los desarrolladores inmobiliarios un porcentaje del valor del proyecto, no de la inversión—, y el WACC es, por lejos, el método más usado. Si la empresa está apalancada, $r_S$ se obtiene con el beta apalancado en la línea del mercado de títulos \vease{sec:u3-hamada}.
\end{definicion}

\subsection{Los cuatro métodos de descuento de flujos y su equivalencia}

Con los flujos definidos en el capítulo~\ref{cap:u2} \vease{sec:u2-flujos}, hay cuatro caminos para valuar la empresa, y todos dan el mismo resultado si cada flujo se descuenta a la tasa que corresponde a su riesgo:
\begin{itemize}
\item \textbf{Equity cash flow}: el ECF descontado a $k_e$ da el valor de las acciones, $E$; los intereses descontados a $k_d$ dan el de la deuda, $D$; y $V = E + D$.
\item \textbf{Free cash flow}: el FCF descontado al WACC después de impuestos.
\item \textbf{Capital cash flow}: el CCF, que ya incluye el ahorro fiscal, descontado al WACC antes de impuestos ($WACC_{ai}$), para no computar dos veces el escudo.
\item \textbf{APV}: el FCF descontado a $k_u$ más el valor presente del ahorro fiscal.
\end{itemize}

\begin{ejemplo}[Equivalencia de los cuatro métodos (perpetuidades; $t = 40\,\%$)]
Datos: $r_f = 5\,\%$, prima de mercado 6\,\%, $\beta_e = 1{,}66$, $k_d = 5\,\%$, intereses 5, ECF 21 y FCF 24.
\begin{itemize}
\item $k_e = 0{,}05 + 0{,}06\times 1{,}66 = 15\,\%$; $E = 21/0{,}15 = 140$; $D = 5/0{,}05 = 100$; $V = 240$.
\item $WACC = 0{,}05\times 0{,}60\times\frac{100}{240} + 0{,}15\times\frac{140}{240} = 10\,\%$, y $V = 24/0{,}10 = 240$.
\item $CCF = ECF + \text{intereses} = 26 = FCF + t\times\text{intereses}$; $WACC_{ai} = 0{,}05\times\frac{100}{240} + 0{,}15\times\frac{140}{240} = 10{,}83\,\%$, y $V = 26/0{,}1083 = 240$.
\item $\beta_u = 1{,}66/(1 + 100\times 0{,}60/140) = 1{,}16$; $k_u = 0{,}05 + 0{,}06\times 1{,}16 = 12\,\%$; $V_U = 24/0{,}12 = 200$; escudo $D\,t = 40$; $V = 240$.
\end{itemize}
El costo del capital propio que surge del CAPM coincide con el de la Proposición II de MM con impuestos: $k_e = 0{,}12 + (0{,}12 - 0{,}05)\times\frac{100\times 0{,}60}{140} = 15\,\%$.
\end{ejemplo}

\begin{formulaBox}[Equivalencia de métodos]
\begin{center}
\begin{tabular}{@{}llc@{}}
\toprule
Flujo & Tasa & Valor de la empresa\\
\midrule
ECF e intereses & $k_e$, $k_d$ & $V = \dfrac{ECF}{k_e} + \dfrac{\text{Intereses}}{k_d}$\\[8pt]
FCF & $WACC = k_d(1-t)\dfrac{D}{V} + k_e\dfrac{E}{V}$ & $V = \dfrac{FCF}{WACC}$\\[8pt]
CCF & $WACC_{ai} = k_d\dfrac{D}{V} + k_e\dfrac{E}{V}$ & $V = \dfrac{CCF}{WACC_{ai}}$\\[8pt]
FCF y ahorro fiscal & $k_u$, $k_d$ & $V = \dfrac{FCF}{k_u} + D\,t$\\
\bottomrule
\end{tabular}
\end{center}
\end{formulaBox}

Los resultados con las tasas de la Proposición II con impuestos coinciden con los del CAPM si los betas se ajustan por el escudo $(1-t)$ y por el endeudamiento. Ninguno de estos cálculos incluye, sin embargo, una compensación por los costos de insolvencia de un endeudamiento elevado \vease{sec:u3-limites}.

%% ---------------------------------------------------------------------
\section{Análisis de riesgos, opciones reales y presupuesto de capital}\label{sec:u3-escenarios}

El VPN de un proyecto depende de pronósticos que pueden fallar. Varias herramientas exploran esa variabilidad y la gravedad de un error de pronóstico.

\begin{definicion}{Análisis de sensibilidad, de escenarios y de punto de equilibrio}
El \textbf{análisis de sensibilidad} recalcula el VPN modificando una sola variable por vez y manteniendo las demás en su valor esperado. El \textbf{análisis de escenarios} considera el movimiento conjunto de varias variables ante situaciones probables (una guerra, un salto del precio del petróleo, una ola de accidentes). El \textbf{análisis del punto de equilibrio} determina las ventas necesarias para no perder. Además, el descuento de flujos suele esconder \textbf{opciones reales}: tanto la de expandirse como la de abandonar tienen valor.
\end{definicion}

\subsection{Análisis de escenarios}

Solar Electronics Corporation (SEC) evalúa un proyecto de motores de avión impulsados por energía solar. Unos pocos accidentes aéreos reducirían los vuelos y la demanda de motores y, aunque no involucraran aviones solares, podrían volver al público reacio a toda tecnología innovadora: el mercado caería al 70\,\% de lo esperado y la participación de SEC a dos tercios. En ese escenario el proyecto daría un flujo anual negativo de \$\,156 millones durante cinco años y, con una inversión de \$\,1.500 millones y una tasa del 15\,\%, el VPN sería $-1.500 - 156\times A^{5}_{0{,}15} = -\$\,2.023$ millones. Una serie de escenarios así ilumina los problemas del proyecto mejor que el análisis de sensibilidad estándar.

\begin{figure}[H]
  \centering
  \includegraphics[width=0.5\linewidth,height=0.42\textheight,keepaspectratio]{U3-72}
  \caption{Pronóstico de flujos de efectivo de SEC en el escenario de un accidente de aviación (tabla 7.4).}
\end{figure}

\subsection{Análisis del punto de equilibrio}

Complementa al análisis de sensibilidad porque muestra cuán graves serían los errores de pronóstico. Se calcula en términos de utilidad contable y de valor presente, y solo con las estimaciones esperadas de cada variable.

\paragraph{Utilidad contable.} El precio de cada motor es \$\,2 millones y su costo variable \$\,1 millón: cada motor aporta \$\,1 millón de \emph{margen de contribución} antes de impuestos. Los costos fijos son \$\,1.791 millones y la depreciación, \$\,300 millones: \$\,2.091 millones por año cualquiera sea la venta. La utilidad neta pasa de $-1.380$ millones con cero unidades a $-720$ con 1.000, 600 con 3.000 y 5.220 con 10.000.

\begin{formulaBox}[Punto de equilibrio contable]
\[
Q_{\text{contable}} = \frac{\text{Costos fijos} + \text{Depreciación}}{\text{Precio de venta} - \text{Costo variable}} = \frac{2.091}{1} = 2.091 \text{ motores}
\]
Los impuestos no lo alteran: con utilidad antes de impuestos nula tampoco hay impuesto.
\end{formulaBox}

\begin{figure}[H]
  \centering
  \includegraphics[width=0.5\linewidth,height=0.42\textheight,keepaspectratio]{U3-73}
  \caption{Punto de equilibrio usando cifras contables: las rectas de ingresos y costos se cruzan en 2.091 motores (figura 7.1).}
\end{figure}

\begin{figure}[H]
  \centering
  \includegraphics[width=0.5\linewidth,height=0.42\textheight,keepaspectratio]{U3-74}
\end{figure}
\begin{figure}[H]
  \centering
  \includegraphics[width=0.5\linewidth,height=0.42\textheight,keepaspectratio]{U3-75}
  \caption{Ingresos, costos y utilidad del proyecto de SEC según distintos supuestos de ventas.}
\end{figure}

\paragraph{Valor presente.} Interesa más el valor presente que la utilidad. Con una tasa del 15\,\%, el VPN del proyecto va de $-5.120$ millones con cero ventas a $-2.908$ con 1.000, 1.517 con 3.000 y 17.004 con 10.000. La inversión inicial de \$\,1.500 millones equivale a un \textbf{costo anual equivalente} (CAE) de cinco años: $1.500/A^{5}_{0{,}15} = \$\,447{,}5$ millones, mayor que la depreciación anual de 300 porque supone que los fondos podrían haberse invertido al 15\,\%.

\begin{formulaBox}[Punto de equilibrio del valor presente (financiero)]
\[
\begin{aligned}
Q_{VPN} &= \frac{CAE + \text{Costos fijos}\,(1-t_c) - \text{Depreciación}\times t_c}{(\text{Precio} - \text{Costo variable})(1-t_c)}\\[4pt]
&= \frac{447{,}5 + 1.791\times 0{,}66 - 300\times 0{,}34}{1\times 0{,}66} = \frac{1.528}{0{,}66} = 2.315
\end{aligned}
\]
La depreciación entra con signo negativo porque su escudo fiscal compensa costos.
\end{formulaBox}

\begin{figure}[H]
  \centering
  \includegraphics[width=0.5\linewidth,height=0.42\textheight,keepaspectratio]{U3-76}
\end{figure}
\begin{figure}[H]
  \centering
  \includegraphics[width=0.5\linewidth,height=0.42\textheight,keepaspectratio]{U3-77}
\end{figure}
\begin{figure}[H]
  \centering
  \includegraphics[width=0.5\linewidth,height=0.42\textheight,keepaspectratio]{U3-78}
\end{figure}
\begin{figure}[H]
  \centering
  \includegraphics[width=0.5\linewidth,height=0.42\textheight,keepaspectratio]{U3-79}
  \caption{Punto de equilibrio usando el valor presente neto: costo anual equivalente y cruce de los valores presentes de ingresos y costos en 2.315 motores (figura 7.2).}
\end{figure}

¿Por qué difieren ambos puntos? Con 2.091 motores, SEC cubre la depreciación de \$\,300 millones y los demás costos, pero no el costo de oportunidad de los \$\,1.500 millones invertidos: el costo anual verdadero de recuperar la inversión es de \$\,447,5 millones. La depreciación subestima ese costo, y por eso las empresas que alcanzan el equilibrio contable en realidad pierden dinero: pierden el costo de oportunidad de la inversión.

\begin{figure}[H]
  \centering
  \includegraphics[width=0.5\linewidth,height=0.42\textheight,keepaspectratio]{U3-80}
\end{figure}
\begin{figure}[H]
  \centering
  \includegraphics[width=0.5\linewidth,height=0.42\textheight,keepaspectratio]{U3-81}
\end{figure}
\begin{figure}[H]
  \centering
  \includegraphics[width=0.5\linewidth,height=0.42\textheight,keepaspectratio]{U3-82}
\end{figure}
\begin{figure}[H]
  \centering
  \includegraphics[width=0.5\linewidth,height=0.42\textheight,keepaspectratio]{U3-83}
  \caption{Análisis complementarios del riesgo del proyecto de SEC.}
\end{figure}

%% =====================================================================
%%  UNIDAD 3 (parte D)
%% =====================================================================
\section{Punto de equilibrio}\label{sec:u3-pe}

El análisis del punto de equilibrio supone funciones lineales de ingresos y costos, aunque puede extenderse a funciones no lineales y aun con discontinuidades a medida que crece la producción. Su propósito principal es descubrir cómo cambian los costos y los ingresos de la organización con el volumen de producción.

\subsection{Determinación del punto de equilibrio}

\begin{definicion}{Punto de equilibrio y margen de contribución}
El \textbf{punto de equilibrio} es el nivel de producción o de ventas que hace nula la utilidad total, es decir, que iguala los ingresos totales con los costos totales (fijos más variables). Requiere estimar con cuidado los costos fijos, los variables y los ingresos. La diferencia entre el precio y el costo variable unitario es el \textbf{margen} o \textbf{contribución marginal unitaria}; expresada como proporción del precio es la \textbf{contribución marginal porcentual}.
\end{definicion}

\begin{formulaBox}[Punto de equilibrio en unidades y en valor]
\[
\begin{gathered}
CT = CF + CVU\cdot Q
\qquad\qquad
IT = P\cdot Q
\qquad\qquad
PE:\ IT = CT \;\Longrightarrow\; Q\,(P - CVU) = CF
\\[4pt]
Q_E = \frac{CF}{P - CVU}
\qquad\qquad
IT_E = \frac{CF}{1 - \dfrac{CVU}{P}} = \frac{CF}{1 - \dfrac{CVT}{V}}
\end{gathered}
\]
Con varios productos, el equilibrio en pesos se obtiene con el margen de contribución sobre ventas.
\end{formulaBox}

Los textos difieren en qué costos fijos incluir. En el enfoque de este apartado solo se consideran los costos \emph{erogables}, de modo que depreciaciones y amortizaciones quedan excluidas; en el \textbf{punto de equilibrio económico} (PEE) de Dumrauf, en cambio, la depreciación integra los costos fijos, se calcula antes de intereses e impuestos (que no lo afectan, porque en el equilibrio no hay ganancia gravable), y el equilibrio sin depreciación se llama \emph{de efectivo} \vease{sec:u3-pef}.

% Figura opcional (el material la menciona como ilegible en el OCR):
% se incluye solo si su archivo está en la carpeta.
\IfFileExists{U3-84.png}{%
\begin{figure}[H]
  \centering
  \includegraphics[width=0.5\linewidth,height=0.42\textheight,keepaspectratio]{U3-84}
  \caption{Representación gráfica del punto de equilibrio con funciones lineales de ingreso y costo total.}
\end{figure}}{}

\begin{ejemplo}[Ejercicio de aplicación 5.4: punto de equilibrio]
Una planta de 168.000 unidades mensuales vende a \$\,25 por unidad; su costo unitario total es de \$\,18 trabajando al 80\,\% de la capacidad, y su costo fijo anual, de \$\,20.000.000. La capacidad anual es de 2.016.000 unidades (1.612.800 al 80\,\%); el costo fijo unitario a ese nivel es $20.000.000/1.612.800 = \$\,12{,}40$ y el variable unitario, $18 - 12{,}40 = \$\,5{,}60$.
\begin{enumerate}
\item $Q_E = \dfrac{20.000.000}{25 - 5{,}60} = 1.030.928$ unidades por año.
\item Al 90\,\% de la capacidad (1.814.400 unidades): $R = (25 - 5{,}60)\times 1.814.400 - 20.000.000 = \$\,15.199.360$ de ganancia anual.
\end{enumerate}
\end{ejemplo}

% Figura opcional (el material la menciona como ilegible en el OCR):
% se incluye solo si su archivo está en la carpeta.
\IfFileExists{U3-85.png}{%
\begin{figure}[H]
  \centering
  \includegraphics[width=0.5\linewidth,height=0.42\textheight,keepaspectratio]{U3-85}
  \caption{Punto de equilibrio de una planta que trabaja al 80\,\% de su capacidad: resolución del ejercicio de aplicación 5.4.}
\end{figure}}{}

La herramienta es útil para decisiones sobre la estructura operativa que aumentan costos fijos o variables, para evaluar una baja de precios que busque más volumen o el lanzamiento de un producto.

\begin{formulaBox}[Margen de seguridad y equilibrio con utilidad deseada]
\[
\text{Margen de seguridad} = \frac{\text{Ventas} - \text{Ventas de equilibrio}}{\text{Ventas}}
\qquad
Q = \frac{CF + \text{ut. deseada}}{P - CVU}
\]
El \textbf{margen de seguridad} mide cuánto pueden caer las ventas antes de entrar en pérdida. Para obtener una utilidad deseada, esta se suma a los costos fijos.
\end{formulaBox}

\begin{ejemplo}[Margen de seguridad y utilidad deseada]
Con costos fijos de 20, precio 1 y costo variable unitario 0,50, el equilibrio es de 40 unidades ($20/0{,}50$) o \$\,40 de ventas ($20/[1 - 50/100]$). Si se venden 100, el margen de seguridad es $(100-40)/100 = 60\,\%$. Para ganar 15 hacen falta $(20+15)/0{,}50 = 70$ unidades, o \$\,70.
\end{ejemplo}

\begin{figure}[H]
  \centering
  \includegraphics[width=0.5\linewidth,height=0.42\textheight,keepaspectratio]{screenshot001}
\end{figure}
\begin{figure}[H]
  \centering
  \includegraphics[width=0.5\linewidth,height=0.42\textheight,keepaspectratio]{screenshot003}
  \caption{Punto de equilibrio económico, margen de seguridad y equilibrio con utilidad deseada.}
\end{figure}

El equilibrio sube si aumentan los costos fijos o los variables y baja si sube el precio; un cambio en la cantidad vendida no lo modifica, solo amplía la ganancia o la pérdida. El modelo admite tres matices: los costos fijos suelen crecer en escalera por intervalos de actividad; los economistas objetan la linealidad, porque los costos variables unitarios y los precios cambian con el volumen, lo que genera una \emph{zona de ganancia} y no un punto único; y por eso el cálculo vale dentro de un \emph{rango relevante} de actividad, el habitual de la empresa.

\paragraph{Punto de equilibrio y estrategia de precios.} La contribución marginal se relaciona con la rotación: los productos de menor salida suelen tener márgenes más altos (los artículos de lujo más que la ropa, y la ropa rota menos que los alimentos). Dos locales del mismo centro comercial con costos fijos de \$\,3.000 alcanzan el equilibrio de modo muy distinto: una relojería con contribución del 60\,\% necesita vender \$\,5.000, diez relojes de \$\,500; una tienda de ropa con contribución del 20\,\% necesita \$\,15.000, 300 prendas de \$\,50.

\subsection{Punto de equilibrio de efectivo y punto de equilibrio financiero}\label{sec:u3-pef}

\begin{definicion}{Puntos de equilibrio de efectivo y financiero}
El \textbf{punto de equilibrio de efectivo} es el nivel de ventas que anula el flujo de caja, suponiendo ventas al contado y costos erogables: excluye la depreciación, que no es salida de caja, y por eso requiere menos ventas que el PEE (en el PEE la empresa ya genera un flujo igual a su depreciación). El \textbf{punto de equilibrio financiero} agrega el pago de intereses y la devolución del capital o, mejor, exige que el VAN sea nulo, es decir, que los flujos descontados al costo de capital recuperen la inversión. Es el verdadero umbral: por debajo de él la empresa destruye valor aunque contablemente parezca en equilibrio o en ganancia.
\end{definicion}

\begin{formulaBox}[Equilibrio de efectivo y financiero]
\[
\begin{gathered}
FCF = (P - CVU)\,Q - CF + Dep = (P - CVU)\,Q - (CF - Dep)
\\[4pt]
Q_{\text{efectivo}} = \frac{CF - Dep}{P - CVU}
\qquad\qquad
Q_{\text{financiero}} = \frac{(CF - Dep) + \text{Intereses} + \text{Amortización}}{P - CVU}
\\[4pt]
\text{en valor presente:}\quad
\text{Inversión} = FCF\times\frac{(1+WACC)^n - 1}{(1+WACC)^n\,WACC}
\end{gathered}
\]
Sin impuestos; $CF$ son los costos fijos totales, que incluyen la depreciación ($Dep$).
\end{formulaBox}

\begin{ejemplo}[Los tres puntos de equilibrio]
Con los datos anteriores, los costos fijos de 20 incluyen 5 de depreciación. El equilibrio de efectivo es $15/0{,}50 = 30$ unidades, menor que el PEE de 40. Si además hay que pagar 75 de intereses y amortizaciones, el equilibrio financiero es $90/0{,}50 = 180$. Con el criterio del valor presente, una inversión de 500 que debe recuperarse en diez años al 10\,\% exige un flujo anual de 81,37; como $81{,}37 = 0{,}50\,Q - 15$, el equilibrio es $Q = 96{,}37/0{,}50 = 192{,}74$ unidades.
\end{ejemplo}

\begin{figure}[H]
  \centering
  \includegraphics[width=0.5\linewidth,height=0.42\textheight,keepaspectratio]{screenshot004}
\end{figure}
\begin{figure}[H]
  \centering
  \includegraphics[width=0.5\linewidth,height=0.42\textheight,keepaspectratio]{screenshot005}
  \caption{Punto de equilibrio de efectivo y punto de equilibrio financiero.}
\end{figure}

\subsection{Determinación del punto de cierre}

\begin{definicion}{Punto de cierre}
Es el nivel de operaciones en el que los ingresos totales igualan el costo fijo total. Aun en la zona de pérdidas conviene mantener abierta la planta mientras los ingresos cubran el costo fijo y, al menos, una parte de los costos variables; cuando no alcanzan a cubrir el costo fijo total, conviene cerrarla. Puede representarse gráficamente, como el punto de equilibrio.
\end{definicion}

\subsection{Selección entre alternativas tecnológicas}

Dervitsiotis plantea que el punto de equilibrio cambia con la técnica de producción. Si los ingresos son iguales con cualquier tecnología pero sus costos varían con el volumen —cada una combina costos fijos y variables distintos—, las rectas de costo total se cruzan en ciertos volúmenes $Q_1$ y $Q_2$. Por debajo de $Q_1$ conviene la alternativa 1, entre $Q_1$ y $Q_2$ la 2 y por encima de $Q_2$ la 3: el criterio es minimizar el costo total con un precio fijo. En el ejercicio de aplicación 5.5 la recomendación es comprar el modelo B, que minimiza el costo total para el volumen de producción esperado. La representación gráfica muestra a la vez todas las opciones de selección del problema.

% Figura opcional (el material la menciona como ilegible en el OCR):
% se incluye solo si su archivo está en la carpeta.
\IfFileExists{U3-87.png}{%
\begin{figure}[H]
  \centering
  \includegraphics[width=0.5\linewidth,height=0.42\textheight,keepaspectratio]{U3-87}
\end{figure}
\begin{figure}[H]
  \centering
  \includegraphics[width=0.8\linewidth,height=0.42\textheight,keepaspectratio]{U3-88}
  \caption{Selección entre alternativas tecnológicas según el volumen de producción.}
\end{figure}}{}

% Figura opcional (el material la menciona como ilegible en el OCR):
% se incluye solo si su archivo está en la carpeta.
\IfFileExists{U3-89.png}{%
\begin{figure}[H]
  \centering
  \includegraphics[width=0.5\linewidth,height=0.42\textheight,keepaspectratio]{U3-89}
\end{figure}
\begin{figure}[H]
  \centering
  \includegraphics[width=0.5\linewidth,height=0.42\textheight,keepaspectratio]{U3-90}
\end{figure}
\begin{figure}[H]
  \centering
  \includegraphics[width=0.5\linewidth,height=0.42\textheight,keepaspectratio]{U3-91}
  \caption{Ejercicio de aplicación 5.5: enunciado, resolución y solución gráfica.}
\end{figure}}{}

\subsection{Leverage operativo, financiero y combinado}\label{sec:u3-leverage}

Los costos fijos amplifican los cambios: los operativos, los del resultado operativo ante cambios en las ventas; los financieros (intereses), los del resultado neto ante cambios en el resultado operativo.

\begin{definicion}{Grados de leverage}
El \textbf{leverage operativo} (LOP) mide la variación porcentual del EBIT ante una variación porcentual de las ventas; siempre es mayor que 1 si hay costos fijos, tiende a infinito cerca del punto de equilibrio y disminuye a medida que las ventas se alejan de él. Cuanto mayores los costos fijos, mayor el \emph{riesgo de negocio}, es decir, la variabilidad del resultado operativo. El \textbf{leverage financiero} (LF) mide la variación porcentual del resultado neto ante una variación porcentual del EBIT: si los activos rinden más que el costo de la deuda la palanca es positiva y mejora el ROE, y si rinden menos es negativa; cuanto mayor la carga de intereses, mayor el \emph{riesgo financiero}, la variabilidad de la utilidad neta. El \textbf{leverage combinado} (LC) es el producto de ambos: en los buenos momentos la palanca financiera refuerza a la operativa; en los malos, ambas juegan en contra.
\end{definicion}

\begin{formulaBox}[Leverage operativo, financiero y combinado]
\[
\begin{gathered}
LOP = \frac{\Delta EBIT/EBIT}{\Delta Q/Q} = \frac{Q\,(P - CVU)}{Q\,(P - CVU) - CF} = \frac{\text{Contribución marginal}}{EBIT}
\\[4pt]
LF = \frac{\Delta RN/RN}{\Delta EBIT/EBIT} = \frac{Q\,(P - CVU) - CF}{Q\,(P - CVU) - CF - \text{Intereses}} = \frac{EBIT}{EBIT - \text{Intereses}}
\\[4pt]
LC = LOP\times LF = \frac{\text{Ventas} - \text{Costo variable total}}{EBIT - \text{Intereses}}
\\[4pt]
\%\Delta EBIT = LOP\times\%\Delta\text{Ventas}
\qquad
\%\Delta RN = LF\times\%\Delta EBIT = LC\times\%\Delta\text{Ventas}
\end{gathered}
\]
\end{formulaBox}

El análisis del punto de equilibrio y del leverage tiene una advertencia general: suele hacerse sobre el beneficio contable cuando sería más apropiado hacerlo sobre el VAN, y trabaja solo con las estimaciones esperadas de cada variable. El efecto del leverage operativo sobre el beta se analiza en el capítulo~\ref{cap:u5} \vease{sec:u5-detbeta}.

%% ---------------------------------------------------------------------
\begin{control}
\pregunta{1. ¿Cuál es el tema principal?}
Si la forma de financiar la empresa —la mezcla de deuda y capital propio— cambia su valor y su costo de capital, y cómo se mide el riesgo operativo y financiero que asume.

\pregunta{2. ¿Dónde y cómo se desarrolla?}
Parte de los instrumentos de financiamiento y del modelo del pastel; demuestra con MM que sin impuestos la estructura es irrelevante; incorpora los impuestos (escudo fiscal), los costos de las dificultades financieras, los conflictos de agencia, la información asimétrica y los impuestos personales; define el costo de capital promedio ponderado y los métodos de valuación con deuda; y termina con el análisis de escenarios, los puntos de equilibrio y el leverage.

\pregunta{3. ¿Cuáles son los temas secundarios?}
Apalancamiento casero y análisis EBIT-EPS; visión tradicional; Proposición III; betas apalancadas; estrategias egoístas y cláusulas de protección; derechos negociables y no negociables; modelo ECO; señales; flujo de efectivo libre; selección jerárquica; crecimiento; equilibrio de Miller; esquemas de ponderación; APV, FTE, WACC y CCF; punto de cierre y alternativas tecnológicas.

\pregunta{4. ¿Las ideas principales y secundarias están relacionadas entre sí?}
Sí. Cada imperfección que se agrega a MM —impuestos, quiebras, agencia, información, impuestos personales— explica por qué las empresas no son indiferentes al endeudamiento ni se endeudan al máximo. El resultado es un costo de capital que depende de la estructura y que se usa, con el flujo adecuado, para valuar; y el riesgo que ese costo remunera crece con los costos fijos operativos y financieros que miden los puntos de equilibrio y el leverage.

\pregunta{5. ¿Cuál es la conclusión o la idea general?}
La deuda crea valor por el escudo fiscal y por disciplinar a los gerentes, pero lo destruye por los costos de las dificultades financieras y de agencia; la estructura óptima equilibra ambos efectos, y en la práctica las empresas siguen además una jerarquía de fuentes. El WACC, bien calculado y aplicado al flujo correcto, traduce esa decisión en una tasa para valuar, y un proyecto solo crea valor si supera su equilibrio financiero, no el contable.
\end{control}

%% =====================================================================
%%  UNIDAD 4 (parte A): política de dividendos y de distribución
%% =====================================================================
\chapter{Política de dividendos y valuación de bonos y acciones}\label{cap:u4}

\begin{ficha}
\fichabase{4 — Política de dividendos y valuación de bonos y acciones}
\fila{Texto base}{S. Ross, R. Westerfield y J. Jaffe, \emph{Finanzas corporativas}, McGraw-Hill: «Política de dividendos: ¿por qué importa?» (esquema de repaso del capítulo).}
\fila{}{R. Brealey, S. Myers y F. Allen, \emph{Principios de finanzas corporativas}, McGraw-Hill: «Política de distribución» (resumen del capítulo).}
\fila{}{G.\,L. Dumrauf, \emph{Finanzas corporativas: un enfoque latinoamericano}, 3.ª ed., Alfaomega: cap.~6 «Valuación de bonos y acciones», incluido «Determinantes del costo de oportunidad del dinero».}
\fila{Complementaria}{Brealey, Myers y Allen, ob. cit.: «El valor de las acciones ordinarias» (resumen). Ross, Westerfield y Jaffe, ob. cit.: resumen «Cómo valorar bonos y acciones» y preguntas de repaso sobre tipos de dividendos, fechas de pago, dividendos caseros y modelo de crecimiento de dividendos.}
\fila{}{Apuntes de cátedra: modelo de descuento de dividendos, crecimiento, PVGO y política de dividendos.}
\fila{Fórmulas previas}{Perpetuidades y anualidades, simples y crecientes \vease{sec:u1-va}; razón precio-utilidad \vease{sec:u2-ratios}; FCF y ECF \vease{sec:u2-flujos}; acciones preferidas y acciones en cartera \vease{sec:u3-flp}.}
\end{ficha}


%% ---------------------------------------------------------------------
\section{Política de dividendos: ¿por qué importa?}\label{sec:u4-dividendos}

Las sociedades distribuyen efectivo a sus accionistas de dos maneras: pagan dividendos o recompran parte de las acciones en circulación. Cuando retienen utilidades lo hacen por dos razones: para invertir en crecimiento o para «sobrevivir» a tiempos difíciles. La recompra fue rara hasta los años ochenta; desde entonces es cada vez más común. La \emph{decisión de distribución} \vease{sec:u1-alcance} fija, a la vez, lo que reciben los accionistas y los fondos que la empresa conserva para invertir.

\begin{definicion}{Índice de pago (\emph{payout}) y tasa de retención}
El \textbf{índice de pago} $d$ es la proporción de las utilidades que se distribuye como dividendos: dividendos sobre utilidad neta o, por acción, DPA/BPA. La \textbf{tasa de retención} $b = 1-d$ es la proporción que se reinvierte. La política de dividendos de una empresa suele describirse por su índice de pago.
\end{definicion}

\subsection{Tipos de dividendos y mecánica del pago}

\begin{definicion}{Dividendo, distribución y dividendo en acciones}
Un \textbf{dividendo} es un pago directo de la sociedad a sus accionistas con cargo a las ganancias corrientes o acumuladas; si el pago proviene de otras fuentes se lo llama \textbf{distribución}. El tipo más común es el \textbf{dividendo en efectivo}: las empresas que cotizan suelen pagarlo cuatro veces al año. El \textbf{dividendo en acciones} se expresa como proporción —con un 2\,\%, el accionista recibe una acción nueva por cada 50 que posee— y no es un dividendo verdadero, porque no sale efectivo de la empresa: equivale a un desdoblamiento (\emph{split}).
\end{definicion}

Por eso suele preferirse el dividendo en efectivo: el dividendo en acciones solo reparte el mismo patrimonio entre más acciones, de modo que el precio de cada una baja en proporción (dilución) y el accionista no recibe dinero.

\begin{definicion}{Fechas del pago de un dividendo en efectivo}
\begin{itemize}
\item \textbf{Fecha de declaración}: el directorio aprueba pagar el dividendo a los accionistas registrados en una fecha determinada.
\item \textbf{Fecha ex dividendo}: dos días hábiles antes de la fecha de registro. Hasta ese día la acción cotiza \emph{con dividendo} (\emph{cum dividend}); desde entonces el comprador ya no lo cobra, por lo que el precio debería caer.
\item \textbf{Fecha de registro}: la sociedad confecciona la lista de las personas que considera accionistas.
\item \textbf{Fecha de pago}: se envían los pagos.
\end{itemize}
\end{definicion}

\begin{figure}[H]
  \centering
  \includegraphics[width=0.5\linewidth,height=0.42\textheight,keepaspectratio]{bma process-32}
  \caption{Dividendos y recompras: cómo distribuyen efectivo las empresas (Brealey, Myers y Allen).}
\end{figure}

\subsection{La proposición de indiferencia de MM}

Modigliani y Miller (MM) analizan la política de dividendos con tres supuestos: \textbf{mercados perfectos} (sin impuestos ni comisiones, y ningún participante puede mover el precio con sus operaciones); \textbf{expectativas homogéneas} (todos comparten las mismas creencias sobre inversiones, ganancias y dividendos futuros), y una \textbf{política de inversión fijada de antemano}, que no cambia por la política de dividendos.

\begin{definicion}{Proposición de indiferencia de los dividendos}
Con esos supuestos, el inversor es indiferente a la política de dividendos y el valor de la empresa depende solo de su política de inversión, incluidas sus oportunidades de crecimiento. Aumentar el dividendo emitiendo acciones nuevas, o reducirlo recomprando acciones, no ayuda ni perjudica al accionista.
\end{definicion}

El argumento es el siguiente. Fijados el programa de inversiones y el endeudamiento, el dinero extra para un dividendo mayor solo puede salir de una emisión de acciones. Para que el valor de mercado total no cambie, tiene que haber una \emph{transferencia de valor} de los accionistas antiguos a los nuevos: estos reciben las acciones emitidas, cada acción vale menos que antes del anuncio y los antiguos sufren una pérdida de capital que compensa exactamente el dividendo extra que cobran. Además, cada accionista puede deshacer por su cuenta la política de la empresa, reinvirtiendo lo que le sobra o vendiendo acciones si quiere más efectivo: son los \emph{dividendos caseros}.

\begin{ejemplo}[Dividendos caseros]
Una inversora prefiere cobrar \$\,10 por acción al final de cada uno de los dos próximos años, y la empresa anuncia otra política: \$\,11 y \$\,8,90. No tiene por qué inquietarse: reinvierte al 10\,\% el peso que no necesita en la primera fecha y en la segunda obtiene \$\,1,10 más, con lo que vuelve a cobrar $8{,}90 + 1{,}10 = 10$. Con una tasa del 10\,\%, ambas corrientes valen lo mismo:
\[
P_0 = \frac{10}{1{,}10} + \frac{10}{1{,}10^{2}} = \frac{11}{1{,}10} + \frac{8{,}90}{1{,}10^{2}} = 17{,}35
\]
\end{ejemplo}

El tercer supuesto deja un corolario práctico: nunca debe resignarse un proyecto con VPN positivo para aumentar un dividendo o para pagarlo por primera vez.

\subsection{Impuestos, costos de emisión y dividendos}

Los dividendos en efectivo tributan como renta ordinaria cuando se cobran. Las ganancias de capital suelen gravarse a tasas algo menores y, además, el impuesto se difiere hasta que se vende la acción. Para el accionista individual, entonces, la tasa efectiva sobre los dividendos es mayor que la tasa sobre las ganancias de capital. De allí surgen tres conclusiones:
\begin{itemize}
\item Si la empresa no tiene efectivo suficiente, hay consenso en que, en un mundo con impuestos personales, no debe emitir acciones para pagar un dividendo: el accionista paga el impuesto y la empresa, los costos de emisión.
\item Si le sobra efectivo después de aceptar todos los proyectos con VPN positivo, las alternativas al dividendo tienen costos: aceptar proyectos adicionales de VPN negativo deja a la empresa expuesta a tomas de control, compras apalancadas y luchas por poderes (\emph{proxy fights}); adquirir otras empresas rara vez es rentable para la compradora, aun con un propósito comercial válido; comprar activos financieros conviene o no según las tasas personales y corporativas, y una empresa inversora excluye de su impuesto el 70\,\% de los dividendos que cobra; la cuarta alternativa, la recompra, se trata a continuación.
\item Los directivos tienen un incentivo para buscar usos alternativos de los fondos y reducir los dividendos; aun así, los impuestos personales no alcanzan para que las empresas eliminen todos los dividendos.
\end{itemize}

\subsection{Recompra de acciones}

La recompra se convirtió en una forma importante de distribuir ganancias y complementa la política de dividendos cuando el ahorro de impuestos pesa. En un mercado perfecto la empresa es indiferente entre pagar un dividendo y recomprar; con la legislación fiscal vigente, en cambio, el accionista suele preferir la recompra. Por eso el fisco (el IRS, en Estados Unidos) puede penalizar a la sociedad cuyo único motivo para recomprar sea eludir el impuesto sobre los dividendos. Muchas empresas recompran, además, porque creen que sus propias acciones son su mejor inversión, sobre todo cuando consideran que el precio está transitoriamente deprimido.

\begin{definicion}{Formas de recomprar acciones}
Las acciones recompradas quedan en la cartera de la empresa (acciones de tesorería) y pueden revenderse si necesita fondos. Se recompran de cuatro maneras:
\begin{itemize}
\item anunciando un plan de compras en el \textbf{mercado abierto} (la forma más común);
\item con una \textbf{oferta pública de adquisición} (\emph{tender offer}) de un número determinado de acciones a un precio fijo;
\item con una \textbf{subasta holandesa}, en la que la empresa fija una serie de precios a los que está dispuesta a recomprar;
\item por \textbf{negociación directa} con uno o más accionistas importantes.
\end{itemize}
\end{definicion}

La reacción inmediata del mercado al anuncio de una recompra suele ser favorable, e Ikenberry, Lakonishok y Vermaelen (1995) muestran que, a largo plazo, las acciones de las empresas que recompran rinden significativamente más que las de empresas comparables que no lo hacen.

\subsection{Rendimiento esperado, dividendos e impuestos personales}

Si los dividendos tributan más, el rendimiento esperado antes de impuestos de un título con alto rendimiento por dividendos debe superar al de un título idéntico con rendimiento bajo, para compensar el impuesto. La implicancia es que un inversor exento debería concentrarse en títulos de alto rendimiento por dividendos. La evidencia está dividida: Brennan (1970) y Litzenberger y Ramaswamy (1979) encuentran una asociación positiva entre los rendimientos esperados antes de impuestos y los rendimientos por dividendos; Black y Scholes (1974) y Miller y Scholes (1982) no encuentran ninguna relación, y Fama y French (1992) sostienen que la relación aparece solo porque el rendimiento por dividendos está correlacionado con la razón valor contable/valor de mercado (B/M).

\subsection{Factores del mundo real que favorecen dividendos altos}

\begin{itemize}
\item \textbf{Deseo de ingresos corrientes}. Ciertos inversores necesitan efectivo regular, y el dividendo les evita los costos de transacción y las molestias de vender acciones periódicamente.
\item \textbf{Resolución de la incertidumbre}. Según Gordon, un dividendo alto beneficia al accionista porque le despeja dudas: cobra hoy lo que, retenido, sería un resultado futuro incierto.
\item \textbf{Arbitraje fiscal}. Hay estrategias que permiten cobrar dividendos sin soportar todo su costo impositivo (Miller y Scholes, 1978).
\item \textbf{Costos de agencia}. Entre accionistas y bonistas, un dividendo alto es una forma de retirar valor de la empresa antes que los acreedores, por lo que estos suelen restringir los dividendos en los contratos de préstamo. Entre accionistas y directivos, a los gerentes les resulta más fácil perseguir objetivos propios cuando la empresa tiene mucho flujo de fondos libre \vease{sec:u3-limites}; el dividendo les quita ese efectivo y funciona como disciplina, de modo que puede reducir los costos de agencia.
\end{itemize}

\subsection{¿Una resolución de los factores del mundo real?}

Hay dos posiciones opuestas. Por los impuestos, el precio de la acción podría relacionarse \emph{negativamente} con el dividendo corriente, manteniendo constantes las ganancias futuras; por el deseo de ingresos corrientes y factores afines, podría relacionarse \emph{positivamente}. Está comprobado que el precio sube cuando se aumenta el dividendo y baja cuando se lo reduce u omite, lo que parece apoyar la segunda posición. Pero la observación es compatible también con la primera: como a las empresas no les gusta recortar dividendos, solo los aumentan cuando esperan que las ganancias y el flujo de fondos futuros crezcan lo suficiente para no tener que volver atrás. El aumento es, entonces, una señal de la dirección al mercado.

\begin{definicion}{Contenido informativo de los dividendos}
Es la suba del precio que sigue a la señal de un aumento del dividendo: el dividendo da a los inversores una pista sobre la rentabilidad «real» de la empresa, sobre todo en países donde la información contable es poco confiable por el secreto o por estructuras societarias de varios niveles. Varios modelos de política de dividendos incorporan este incentivo de la dirección a comunicar información mediante los dividendos.
\end{definicion}

La recompra informa algo distinto, porque suele ser un hecho aislado y no un compromiso de largo plazo de distribuir más efectivo. Las empresas recompran cuando acumularon un exceso de efectivo o cuando quieren sustituir capital propio por deuda; la recompra también puede señalar que los directivos confían en el futuro, es decir, que consideran subvaluada la acción.

\begin{definicion}{Efecto clientela}
Inversores de distintas categorías impositivas prefieren empresas con distintas políticas de dividendos, hasta que la oferta y la demanda de acciones de alto dividendo quedan en equilibrio. Por eso una empresa solo puede aumentar su precio cambiando su política si existe una clientela insatisfecha. La evidencia muestra que las acciones con mayor rendimiento por dividendos están en manos de individuos de categorías impositivas bajas.
\end{definicion}

\subsection{Lo que sabemos y lo que no sabemos sobre la política de dividendos}

Tres hechos están bien establecidos. Los dividendos son sustanciales: el índice de pago promedio fue del 43\,\% entre 1963 y 1998. Cada vez menos empresas pagan dividendos: Fama y French (2001) muestran que el porcentaje de las que lo hacen cayó en los últimos años. Y las empresas \textbf{suavizan} los dividendos: según Lintner (1954), fijan índices de pago objetivo de largo plazo y, como saben que solo una parte de cualquier cambio en las ganancias será permanente y necesitan tiempo para juzgarlo, los cambios de dividendos se rezagan varios períodos respecto de los cambios en las ganancias. De ese comportamiento se desprenden tres rasgos de la decisión de dividendos: los directivos son reacios a hacer cambios que luego deban revertir, suavizan los dividendos para no tener que reducirlos y se fijan más en los cambios del dividendo que en su nivel absoluto.

\subsection{Dividendos en acciones y desdoblamientos}

En un mercado perfecto, ni el dividendo en acciones ni el desdoblamiento tienen valor para la empresa. Se los explica por el \textbf{rango de cotización popular}: cuando el precio lo supera, muchos inversores no pueden comprar un \emph{lote redondo} de 100 acciones, y comprar \emph{lotes impares} (menos de 100) encarece las comisiones; desdoblando, la empresa devuelve el precio al rango. Sin embargo, Copeland observa que después de un desdoblamiento suben dos costos de transacción medidos en porcentaje: las comisiones de corretaje y el diferencial entre precio comprador y vendedor.

\begin{definicion}{Contrasplit (\emph{reverse split})}
Es el canje de varias acciones por una nueva, por ejemplo una por cada tres. Se justifica por tres razones: los costos de transacción para el accionista suelen bajar; la liquidez y la negociabilidad mejoran cuando el precio vuelve al rango de cotización popular, y las acciones de precio muy bajo no se consideran «respetables», por lo que los inversores tienden a subestimar las ganancias, el crecimiento y la estabilidad de esas empresas.
\end{definicion}

%% ---------------------------------------------------------------------
\section{Política de distribución}\label{sec:u4-distribucion}

\subsection{La controversia sobre la distribución}

Sobre la política de distribución hay tres posiciones. A la \textbf{derecha}, los conservadores sostienen que los inversores prefieren dividendos más altos. A la \textbf{izquierda}, otro grupo afirma que los dividendos altos reducen el valor porque tributan más que las ganancias de capital. En el \textbf{centro}, los moderados sostienen que la política de distribución no cambia nada, lo que es cierto en mercados de capitales perfectos según la proposición de MM \vease{sec:u4-dividendos}.

\subsubsection{La derecha}

Responde a MM con las imperfecciones del mundo real que harían preferibles los índices de pago altos. A medida que la empresa madura, dispone de efectivo para distribuir; los dividendos regulares ahorran a muchos accionistas costos de transacción y molestias considerables. Su mejor argumento es el de agencia: repartir los fondos impide que los directivos los usen mal o los derrochen.

\subsubsection{Los impuestos y la izquierda radical}

Su credo es simple: mientras los dividendos tributen más que las ganancias de capital, la empresa debe pagar el menor dividendo en efectivo posible y retener el efectivo disponible o usarlo para recomprar acciones, transformando así los dividendos en ganancias de capital. Hay evidencia de que los impuestos influyeron en la elección de acciones de los inversores, aunque como una consideración secundaria.

\subsubsection{Los moderados}

Sostienen que la política de dividendos no afecta el valor porque la oferta de dividendos se ajusta libremente a la demanda: si las empresas pudieran aumentar su precio cambiando el reparto, ya lo habrían hecho. Las empresas bien administradas pagan dividendos para señalar su valor, aunque eso les cueste más impuestos. Como el valor depende de la política de inversión, y la distribución se decide después de invertir y financiarse, la distribución debería cambiar con el \textbf{ciclo de vida} de la empresa: la empresa joven, con muchas oportunidades de inversión, retiene y reinvierte; la madura tiene menos oportunidades y, como se teme el derroche por sobreinversión, los inversores presionan para que empiece a distribuir; y a medida que envejece debe distribuir cada vez más.

\begin{center}
\small
\begin{tabular}{@{}p{0.17\linewidth}p{0.36\linewidth}p{0.39\linewidth}@{}}
\toprule
Posición & Tesis & Fundamento principal\\
\midrule
Derecha & Dividendos altos aumentan el valor. & Ingresos corrientes, menores costos de transacción y disciplina de los directivos (agencia).\\
\addlinespace
Izquierda & Dividendos altos reducen el valor. & Mayor carga impositiva que las ganancias de capital: retener o recomprar.\\
\addlinespace
Moderados & La política es irrelevante. & MM: el valor depende de la inversión; la oferta de dividendos se adapta a la demanda; el reparto sigue el ciclo de vida.\\
\bottomrule
\end{tabular}
\end{center}

%% =====================================================================
%%  UNIDAD 4 (parte B): valuación de bonos y acciones
%% =====================================================================
\section{Valuación de bonos y acciones}\label{sec:u4-valuacion}

Las técnicas del valor tiempo del dinero \vease{sec:u1-va} se aplican ahora a la fijación de precios de los dos títulos principales que emiten las empresas: las obligaciones y las acciones. El principio es uno solo: el precio de un activo debe reflejar el valor presente de su flujo de fondos futuro, descontado a una tasa que represente el rendimiento de las inversiones de riesgo similar. Con él se puede valuar una obligación y una acción, y calcular el rendimiento asociado a cada inversión.

\subsection{Valor de los activos en un mercado eficiente: la parábola revisada}

Quien compra un activo quiere ganar al menos lo que rinden las inversiones de riesgo similar, de modo que su precio justo es el valor actual de lo que paga descontado a ese rendimiento: para una renta perpetua, la renta dividida por la tasa. Un precio más alto implica un rendimiento menor, que no compensa el riesgo cuando con activos similares se gana más. En un mercado eficiente ese precio no se sostiene: al tenedor le convendría vender el activo y comprar otro de riesgo similar más barato, y el arbitraje devolvería el precio a su valor justo y el rendimiento al exigido por el mercado.

\begin{ejemplo}[Una renta perpetua de \$\,10]
Un activo paga \$\,10 por año a perpetuidad y las inversiones de riesgo similar rinden el 20\,\%: su precio justo es $10/0{,}20 = \$\,50$. Si alguien pagara \$\,60, su rendimiento sería $10/60 = 16{,}66\,\%$; el tenedor ganaría vendiendo el activo y comprando otro de riesgo similar a \$\,50, y el arbitraje devolvería el precio a \$\,50 y el rendimiento al 20\,\%.
\end{ejemplo}

%% ---------------------------------------------------------------------
\subsection{Valuación y rendimiento de la inversión en obligaciones}\label{sec:u4-bonos}

Cuando una empresa necesita una suma grande para un proyecto, reunirla con un solo prestamista es difícil. Por eso recurre al mercado de capitales y divide el crédito en porciones iguales, representadas por títulos negociables que toman los ahorristas.

\begin{definicion}{Bono u obligación}
Título negociable que representa una porción alícuota de un préstamo tomado por el emisor en el mercado de capitales. A diferencia del préstamo tradicional, suele tener plazos largos; las condiciones las fija el deudor, que se obliga a pagar periódicamente un interés a la tasa enunciada en el título y a rescatarlo en la fecha establecida; tiene una pluralidad de acreedores que el emisor no individualiza, porque los títulos se negocian en el mercado; tiene un \textbf{valor nominal} o facial y puede emitirse a la par, bajo la par o sobre la par según el precio de colocación (\textbf{precio de emisión}); y combina una tasa de interés nominal pactada con una tasa de rendimiento implícita en su precio. Las condiciones de intereses, capital, fechas, garantías y restricciones —que a veces limitan el endeudamiento o los dividendos— figuran en el \textbf{prospecto de emisión}.
\end{definicion}

Quien compra un bono adquiere el derecho a cobrar una corriente de intereses y capital. Los bonos pueden tener tasa fija o flotante, no pagar intereses (cupón cero), amortizar en cuotas o en un solo pago, tener período de gracia, incluir opciones y tener o no garantía. Tres tipos básicos sirven de referencia: los \textbf{bonos de descuento puro} o cupón cero, que pagan solo el nominal al vencimiento; los \textbf{bonos con cupón nivelado}, que pagan un cupón fijo y el nominal al final, y los \textbf{consols}, perpetuidades sin vencimiento \vease{sec:u1-va}.

\subsubsection{Conceptos fundamentales en la valuación de bonos}

Los \textbf{intereses} se pagan a intervalos periódicos (semestrales, trimestrales, mensuales) con dos modalidades. Con \emph{tasa fija}, calculada sobre el valor nominal —o sobre el residual, cuando el bono ya devuelve capital—, el inversor obtiene un flujo previsible y gana o pierde según bajen o suban las tasas; al emisor le ocurre lo contrario: si las tasas bajan querrá emitir nueva deuda más barata y, si suben, se felicitará por haberse financiado a tasa fija. Con \emph{tasa flotante}, en cambio, se fija una tasa de referencia —muy usada es la LIBOR, \emph{London Interbank Offered Rate}, tasa interbancaria del mercado de Londres— y en cada período se devenga el interés según su valor al comienzo del período, de modo que los flujos quedan «indexados» a tasas de mercado.

El \textbf{capital} o principal se amortiza en un pago único al vencimiento —el bono \emph{bullet}, cuyo flujo coincide con el sistema americano de amortización— o en pagos periódicos, que pueden coincidir o no con los de intereses.

\begin{definicion}{Valor residual, valor técnico y paridad}
El \textbf{valor residual} (VR) es el capital aún no amortizado; sobre él se calculan los intereses, que bajan con cada pago de capital. En un bono \emph{bullet} coincide siempre con el valor nominal. El \textbf{valor técnico} (VT) suma al residual los intereses devengados del cupón en curso (cupón corrido): es lo que el bono «debería valer» técnicamente. La \textbf{paridad} expresa el precio de mercado $V$ como porcentaje del valor técnico.
\end{definicion}

\begin{formulaBox}[Valor residual, valor técnico y paridad]
\[
\begin{gathered}
\text{VR} = \text{Valor nominal} - \text{Total amortizado}
\qquad\qquad
\text{VT} = \text{VR} + \text{Interés corrido}
\\[4pt]
\text{Paridad} = \frac{V}{\text{VT}}\times 100
\end{gathered}
\]
\end{formulaBox}

\subsubsection{Valuación de un bono con pago del principal al vencimiento}

Si el bono tiene precio de mercado, ¿para qué valuarlo? Valuar es calcular su \textbf{valor intrínseco}: el valor normativo que debería tener si su flujo se descuenta al rendimiento de un activo de riesgo similar. Es el valor presente de los cupones $C$ y del principal $P$:

\begin{formulaBox}[Precio de un bono \emph{bullet}]
\[
B = \frac{C}{1+kd} + \frac{C}{(1+kd)^{2}} + \dots + \frac{C+P}{(1+kd)^{n}}
  = C\,\frac{(1+kd)^{n}-1}{(1+kd)^{n}\,kd} + \frac{P}{(1+kd)^{n}}
\]
$kd$ es el rendimiento que exige el comprador; el primer término del segundo miembro es el valor actual de la renta de cupones.
\end{formulaBox}

\begin{ejemplo}[Santa Emilia S.\,A.]
El 1 de diciembre de 2001 emitió una obligación por 100 millones con cupón del 10\,\% anual y amortización del capital al final del quinto año, cuando los inversores exigían el 11\,\% para riesgos similares:
\[
B = \frac{10}{1{,}11} + \frac{10}{1{,}11^{2}} + \frac{10}{1{,}11^{3}} + \frac{10}{1{,}11^{4}} + \frac{110}{1{,}11^{5}}
  = 10\times\frac{1{,}11^{5}-1}{1{,}11^{5}\times 0{,}11} + \frac{100}{1{,}11^{5}} = 96{,}3
\]
\end{ejemplo}

\sinejemplos{Así, el 1 de diciembre de 2001 Santa Emilia S.\,A. emitió una obligación a cinco años, con cupón del 10\,\% anual y amortización del capital al vencimiento, cuando los inversores exigían el 11\,\% para riesgos similares: descontados a esa tasa, sus pagos valen \$\,96,3 por cada \$\,100 de nominal. }El 11\,\% es la tasa implícita o de arbitraje que iguala los pagos futuros con el precio de \$\,96,3. Que el precio surja de la oferta y la demanda no quiere decir que no haya un modelo detrás: el precio refleja las decisiones de inversores racionales en un mercado eficiente. Si el 11\,\% fuera un rendimiento demasiado alto para ese riesgo, el arbitraje empujaría el precio hacia arriba y el rendimiento hacia abajo.

\subsubsection{¿Por qué la tasa de rendimiento esperada difiere de la tasa del cupón?}

Al emitir, los cupones suelen fijarse en el nivel de las tasas de mercado. Santa Emilia intentó colocar su obligación con un cupón algo menor (10\,\%) y el mercado respondió con un descuento: suscribió el título al 96,3\,\% del nominal y obtuvo así el 11\,\%.

\begin{definicion}{Cotización a la par, bajo la par y sobre la par}
El bono cotiza \textbf{a la par} cuando su precio es igual al valor residual: su TIR coincide con la tasa de contrato del cupón $j_{(m)}$. Como el cupón es fijo, si el precio está por debajo del valor residual el rendimiento supera la tasa del cupón (\textbf{bajo la par}); si está por encima, es menor (\textbf{sobre la par}).
\end{definicion}

\begin{center}
\small
\begin{tabular}{@{}lccc@{}}
\toprule
 & A la par & Bajo la par & Sobre la par\\
\midrule
TIR frente al cupón & TIR $= j_{(m)}$ & TIR $> j_{(m)}$ & TIR $< j_{(m)}$\\
TIR frente al rendimiento corriente & TIR $=$ RC & TIR $>$ RC & TIR $<$ RC\\
\bottomrule
\end{tabular}
\end{center}

\begin{ejemplo}[Bono perpetuo con cupón de \$\,10]
Si las tasas de mercado para riesgos similares son del 10\,\%, el bono vale su nominal, \$\,100. Si suben al 11\,\%, nadie paga 100 para ganar el 10\,\%: aumenta la oferta de bonos y el precio cae hasta que los \$\,10 fijos rinden el 11\,\%, es decir, $10/0{,}11 = 90{,}91$. Si bajan al 9\,\%, el bono se vuelve atractivo, sube su demanda y su precio llega a $10/0{,}09 = 111{,}11$.
\end{ejemplo}

\subsubsection{Medidas de rendimiento de la inversión en bonos}

El bono ofrece dos ganancias en «metálico»: los intereses (\emph{coupon yield}) y las ganancias de capital (\emph{capital gain yield}), que dependen de vender el título más caro de lo que se pagó; si el precio baja, hay pérdida de capital. Como el cupón es una suma fija, funciona para el inversor como una «prima de seguro» —la imagen es de Keynes— que lo cubre hasta ese importe contra una baja de precio: el inversor de Santa Emilia está cubierto hasta \$\,10.

\begin{definicion}{Rendimiento al vencimiento (\emph{yield to maturity})}
Es la TIR del bono: la tasa que iguala el valor presente de su flujo con su precio actual. Representa lo que ganaría el inversor si mantuviera el bono hasta el vencimiento, el emisor no cayera en cesación de pagos (\emph{default}) y los cupones se reinvirtieran a esa misma tasa —para lo cual habría que comprar otro bono igual a una paridad que mantenga la TIR—; no considera, por lo tanto, las oportunidades reales de reinversión.
\end{definicion}

Las calculadoras financieras y las planillas la calculan directamente, pero la intuición es simple: se prueba una tasa y se corrige según el precio obtenido. Al 10\,\% el bono de Santa Emilia valdría \$\,100 y al 15\,\%, \$\,83,2; como cotiza a \$\,96,3, uniendo ambos puntos y buscando la tasa que corresponde a ese precio se obtiene una aproximación de la TIR (11\,\%). El procedimiento exacto, la interpolación lineal reiterada, es una técnica de presupuesto de capital.

\begin{figure}[H]
  \centering
  \includegraphics[width=0.5\linewidth,height=0.42\textheight,keepaspectratio]{Finanzas20corporativas203ed20Dumrauf-38}
  \caption{Aproximación gráfica a la TIR del bono: precio en función de la tasa de descuento.}
\end{figure}

\begin{definicion}{Rendimiento corriente y rendimiento total}
El \textbf{rendimiento corriente} (\emph{current yield}) relaciona el cupón del período con el precio de mercado y solo mide la ganancia por intereses. La TIR, en cambio, usa todos los cupones, la ganancia de capital y la reinversión a la propia TIR; por eso supera al rendimiento corriente cuando el bono cotiza bajo la par, es menor sobre la par y reacciona más ante los cambios de precio. El \textbf{rendimiento total} es la suma del rendimiento corriente y de la ganancia (o pérdida) de capital del período.
\end{definicion}

\begin{formulaBox}[Medidas de rendimiento de un bono]
\[
\text{RC} = \frac{C}{B_0}
\qquad
\text{Ganancia de capital} = \frac{B_1 - B_0}{B_0}
\qquad
\text{Rendimiento total} = \frac{C}{B_0} + \frac{B_1 - B_0}{B_0}
\]
\end{formulaBox}

Si el rendimiento exigido no cambia, el rendimiento total del período es igual a él. En un bono comprado con descuento, el rendimiento corriente es menor que el exigido y lo completa una ganancia de capital; en uno comprado con prima, es mayor y lo compensa una pérdida de capital.

\begin{ejemplo}[Rendimiento total con descuento y con prima]
\textbf{Con descuento.} Si después de un año los inversores siguen exigiendo el 11\,\%, al 1 de diciembre de 2002 la obligación de Santa Emilia, con un cupón menos, vale
\[
B_1 = \frac{10}{1{,}11} + \frac{10}{1{,}11^{2}} + \frac{10}{1{,}11^{3}} + \frac{110}{1{,}11^{4}} = 96{,}89
\]
El rendimiento corriente es $10/96{,}3 = 10{,}38\,\%$, la ganancia de capital $(96{,}89 - 96{,}3)/96{,}3 = 0{,}62\,\%$ y el total, $10{,}38\,\% + 0{,}62\,\% = 11\,\%$.

\textbf{Con prima.} Si la emisión se hubiera hecho cuando se exigía el 8\,\%, con el mismo cupón:
\[
\begin{aligned}
B_0 &= 10\times\frac{1{,}08^{5}-1}{1{,}08^{5}\times 0{,}08} + \frac{100}{1{,}08^{5}} = 107{,}99\\
B_1 &= 10\times\frac{1{,}08^{4}-1}{1{,}08^{4}\times 0{,}08} + \frac{100}{1{,}08^{4}} = 106{,}62
\end{aligned}
\]
El rendimiento corriente sería $10/107{,}99 = 9{,}26\,\%$ y la pérdida de capital, $(106{,}62 - 107{,}99)/107{,}99 = -1{,}26\,\%$: el total vuelve a ser el 8\,\%.
\end{ejemplo}

\paragraph{Consideraciones impositivas.} ¿Convenía comprar la obligación de Santa Emilia bajo la par, con cupón del 10\,\%, u otra con un cupón más realista del 11\,\%? Para un inversor gravado puede ser preferible la primera: tributa sobre \$\,10 de intereses y puede postergar el impuesto sobre la ganancia de capital hasta el vencimiento, y un impuesto más lejano tiene menor valor presente. Quien compra un bono con prima paga impuestos sobre intereses más altos, aunque puede usar la prima para compensarlos. Como los inversores computan el rendimiento después de impuestos, en un mercado eficiente los precios se ajustan hasta ofrecer rendimientos idénticos después de impuestos.

\paragraph{Evolución del precio.} A medida que se acerca la amortización final, el precio converge al valor par: el del bono con prima baja y el del bono con descuento sube, y la prima o el descuento tienden a cero, porque al faltar menos tiempo el descuento compuesto pesa menos en el valor presente de los flujos que restan. Por la misma razón, cuanto más tiempo falta para el vencimiento, más cambia el precio ante un cambio del rendimiento exigido. (Algunos bonos soberanos sudamericanos, emitidos muy por debajo de la par, subieron de precio al acercarse el vencimiento también porque disminuía el riesgo de impago.) La tabla lo muestra para la obligación de Santa Emilia según el rendimiento exigido.

\begin{center}
\small
\begin{tabular}{@{}cccc@{}}
\toprule
Año & Precio con $kd = 11\,\%$ & Precio con $kd = 10\,\%$ & Precio con $kd = 8\,\%$\\
\midrule
0 & 96,30 & 100 & 107,99\\
1 & 96,90 & 100 & 106,62\\
2 & 97,56 & 100 & 105,15\\
3 & 98,29 & 100 & 103,57\\
4 & 99,10 & 100 & 101,85\\
5 & 100,00 & 100 & 100,00\\
\bottomrule
\end{tabular}
\end{center}

\begin{figure}[H]
  \centering
  \includegraphics[width=0.5\linewidth,height=0.42\textheight,keepaspectratio]{Finanzas20corporativas203ed20Dumrauf-43}
  \caption{Evolución del precio hasta el vencimiento para un bono comprado con descuento o con prima.}
\end{figure}

\subsubsection{Detalles en la construcción del flujo de fondos del bono}

Construir el flujo es el paso previo a la valuación (\emph{pricing}). El prospecto especifica cómo se cuentan los días de cada período y del año, y los mercados no siguen todos la misma regla, de modo que una convención mal aplicada da otro precio.

\begin{definicion}{Convenciones para contar los días}
\begin{itemize}
\item \textbf{Actual/365}: año de 365 días y número exacto de días entre fechas de pago consecutivas.
\item \textbf{Actual/Actual}: días exactos del período y del año (366 en un año bisiesto).
\item \textbf{Actual/360}: días exactos del período sobre un año de 360 días.
\item \textbf{30/360}: todos los meses se cuentan de 30 días y el año, de 360; el semestre tiene 180 días y el cupón semestral es la mitad de la tasa anual por el valor residual.
\end{itemize}
\end{definicion}

\begin{ejemplo}[Convención Actual/365]
El 15 de diciembre de 2000 la compañía Apalancamiento Financiero emitió una obligación a cinco años, cupón del 10\,\% anual e intereses semestrales los 15 de junio y 15 de diciembre. Entre el 15/12/2000 y el 15/06/2001 hay 182 días exactos, de modo que el primer cupón es $100\times 0{,}10/365\times 182 = 4{,}99$ (el texto cuenta 183 días y obtiene 5,013). Con la convención 30/360 sería $100\times 0{,}10/2 = 5$.
\end{ejemplo}

\paragraph{Bonos comprados en períodos intermedios de renta.} En la práctica los bonos se negocian entre fechas de pago. Se necesitan dos datos: $p$, los días entre el momento de valuación y el pago del próximo cupón, y $d$, los días del período de renta, que informa el prospecto. Los intereses se calculan con la convención del prospecto, pero para descontar los días se cuentan con exactitud: se usa una TIR anual ajustada por exponente según los días que faltan para cada cobro.

\begin{formulaBox}[Cupón y descuento entre fechas de pago]
\[
C = \text{VR}\times\frac{j_{(m)}}{360/180} \quad(\text{30/360, pago semestral})
\qquad\qquad
VP(C_k) = \frac{C_k}{(1+\text{TIR})^{\,p_k/365}}
\]
$p_k$ son los días exactos entre la compra y el cobro $k$; para cada flujo siguiente se suman al exponente los días que faltan.
\end{formulaBox}

\begin{ejemplo}[Estancias del Pilar]
El 11 de marzo de 2001 se compra a \$\,95 una obligación \emph{bullet} con vencimiento el 15 de diciembre de 2002, cupón del 10\,\% anual y pagos semestrales; con 30/360 cada cupón es $100\times 0{,}10/360\times 180 = 5$.
\begin{center}
\begin{tabular}{@{}lccccc@{}}
\toprule
Fecha & 11/03/01 & 15/06/01 & 15/12/01 & 15/06/02 & 15/12/02\\
\midrule
Días para el cobro & & 96 & 279 & 461 & 644\\
Flujo de caja & $-95$ & 5 & 5 & 5 & 105\\
\bottomrule
\end{tabular}
\end{center}
\[
95 = \frac{5}{(1+\text{TIR})^{96/365}} + \frac{5}{(1+\text{TIR})^{279/365}} + \frac{5}{(1+\text{TIR})^{461/365}} + \frac{105}{(1+\text{TIR})^{644/365}}
\quad\Longrightarrow\quad \text{TIR} \approx 15{,}4\,\%
\]
\end{ejemplo}

\paragraph{Costo efectivo de una obligación negociable.} El rendimiento del inversor puede diferir del costo del emisor porque este, además del descuento de mercado, paga los costos de la colocación. En diciembre de 1993 Lácteos S.\,A. emitió obligaciones por U\$S 30 millones, a tres años, con interés del 10\,\% anual pagadero por semestre y convención Actual/365; el mercado las compró bajo la par, con un descuento de 830.000.

\begin{center}
\small
\begin{tabular}{@{}lrrrr@{}}
\toprule
Fecha de pago & Colocación neta & Intereses & Comisiones y gastos & Flujo del emisor\\
\midrule
03/12/93 & 29.170,0 & & $-$826,9 & 28.343,1\\
01/06/94 & & $-$1.500,0 & $-$1,5 & $-$1.501,5\\
05/12/94 & & $-$1.541,7 & $-$26,0 & $-$1.567,7\\
01/06/95 & & $-$1.500,0 & $-$1,5 & $-$1.501,5\\
03/12/95 & & $-$1.541,7 & $-$26,0 & $-$1.567,7\\
03/06/96 & & $-$1.508,3 & $-$1,5 & $-$1.509,8\\
03/12/96 & & $-$1.541,7 & $-$1,5 & $-$31.543,2\\
\bottomrule
\end{tabular}
\end{center}

En miles; la colocación neta es el nominal menos el descuento, y el último flujo incluye la amortización de 30.000. Los gastos iniciales suman 826,9: comisiones de estructuración (0,10\,\%) y de agencia (0,80\,\%), \emph{underwriters} (entre 1,25 y 1,5\,\% de lo colocado por cada banco), estudio jurídico, auditores, dos calificadoras de riesgo, impresión del prospecto y de las láminas, publicidad y tasa de la CNV; luego se pagan la comisión del agente de pago y los honorarios anuales de las calificadoras. Hasta el descuento, el costo del emisor coincide con el rendimiento del inversor, de alrededor del 11,6\,\%; los costos de emisión —casi el 3\,\% del nominal— elevan el costo efectivo al \textbf{12,85\,\%} anual (TIR no periódica del flujo del emisor, función TIR.NO.PER de Excel).

\subsubsection{La estructura temporal de la tasa de interés}

Para entender por qué bonos distintos se venden a precios distintos hay que relacionar la TIR con el plazo.

\begin{definicion}{Estructura temporal de la tasa de interés (curva de rendimientos)}
Es la relación entre el rendimiento de los bonos y su plazo o vida media; busca medir el precio puro del tiempo, es decir, el rendimiento que exige el mercado por un pago que se hará en un año determinado. Su forma más común es \textbf{ascendente} (los plazos largos rinden más), pero puede ser \textbf{horizontal} (\emph{flat}), \textbf{invertida} o \textbf{con joroba}. La teoría que explica las tasas de los distintos plazos es la teoría de la estructura temporal.
\end{definicion}

Analistas y administradores de carteras siguen la curva en busca de oportunidades de arbitraje. Un procedimiento práctico para estimarla es una regresión logarítmica de los rendimientos de los bonos soberanos sobre el plazo en años o sobre la \emph{duration} modificada, una medida de la vida ponderada del bono que tiene en cuenta el peso de cada flujo y el momento en que se cobra.

\begin{figure}[H]
  \centering
  \includegraphics[width=0.5\linewidth,height=0.42\textheight,keepaspectratio]{Finanzas20corporativas203ed20Dumrauf-48}
  \caption{Curva de rendimientos estimada por regresión logarítmica de la TIR de bonos soberanos sobre el plazo.}
\end{figure}

\begin{definicion}{Tasa contado (\emph{spot}) y tasa futura (\emph{forward})}
La \textbf{tasa contado} de un plazo es la tasa implícita en un bono cupón cero de ese plazo (un bono que no paga intereses y se vende con descuento). La \textbf{tasa futura} de un período es la que debe regir para ese período, según las tasas contado, para que no haya arbitraje entre invertir de una vez a todo el plazo o renovar una inversión más corta.
\end{definicion}

\begin{formulaBox}[Tasa futura implícita]
\[
(1+i_1)\,(1+f_2) = (1+i_2)^{2}
\qquad\Longrightarrow\qquad
f_2 = \frac{(1+i_2)^{2}}{1+i_1} - 1
\]
$i_1$ e $i_2$ son las tasas contado anuales a uno y dos años; $f_2$, la tasa futura del segundo año.
\end{formulaBox}

\begin{ejemplo}[Tasas contado y futura]
Un cupón cero a un año que paga 100 cuesta 90,9: $90{,}9 = 100/(1+i_1)$, de modo que $i_1 = 10\,\%$. Uno a dos años cuesta 81,96: $81{,}96 = 100/1{,}22$, una tasa de 22\,\% en dos años, equivalente a 10,45\,\% anual. La tasa futura del segundo año es $1{,}22/1{,}10 - 1 = 10{,}9\,\%$. Si la tasa del segundo año fuera mayor, todos preferirían comprar el bono a un año y renovarlo; si fuera menor, el de dos años. Como en un mercado eficiente los inversores tienen ambas alternativas, los precios no deben permitir el arbitraje.
\end{ejemplo}

No hay bonos cupón cero para todos los plazos, pero las tasas contado faltantes pueden despejarse con bonos con cupón, mediante la técnica de \emph{bootstrapping}: los primeros flujos de un bono con cupón se descuentan a las tasas contado ya conocidas, de modo que en la ecuación de su precio queda una sola incógnita, la tasa contado de su último plazo.

\begin{ejemplo}[\emph{Bootstrapping}]
Se conocen el cupón cero a un año (precio 90,9), el cupón cero a dos años (81,9) y un bono \emph{bullet} a tres años con cupón del 10\,\% que cotiza a 90:
\[
90 = \frac{10}{1{,}10} + \frac{10}{1{,}22} + \frac{110}{(1+i_3)^{3}}
\quad\Longrightarrow\quad
i_3 \approx 14{,}8\,\% \ \text{anual (el texto informa 14,76\,\%)}
\]
\end{ejemplo}

La conclusión de fondo es que un bono no debería valuarse con una única TIR, sino descontando cada flujo con la tasa contado de su plazo; de lo contrario habría posibilidad de arbitraje.

\subsubsection{Riesgos asociados a la inversión en bonos}

Aunque el bono promete un rendimiento, su inversión tiene riesgos.

\begin{definicion}{Riesgos de la inversión en bonos}
\begin{itemize}
\item \textbf{Tasa de interés}: si las tasas suben, el precio baja hasta que el bono vuelve a rendir lo que ofrece el mercado. El efecto es máximo en el bono cupón cero y poco importa a quien mantiene el bono hasta el vencimiento, siempre que el emisor no incumpla, porque el precio converge a la par.
\item \textbf{Reinversión}: si las tasas bajan, los flujos cobrados se reinvierten a tasas menores y se reduce el capital futuro; afecta sobre todo a inversores institucionales (fondos de pensión, aseguradoras). En el \emph{bullet} se limita a los cupones, es mayor en los bonos que amortizan capital y no existe en el cupón cero.
\item \textbf{Rescate anticipado} (\emph{call}): el emisor puede rescatar el bono a un precio preestablecido y lo hará si las tasas bajan, reemplazándolo por deuda más barata; el inversor tendría que reinvertir todo el capital a una tasa menor y por eso descuenta ese riesgo en el precio.
\item \textbf{Venta anticipada} (\emph{put}): da al tenedor el derecho a vender el bono al emisor al valor par en fechas fijadas, pagadero en efectivo u otros activos.
\item \textbf{Conversión}: el tenedor puede canjear el bono por acciones a un precio fijado; el bono convertible vale como obligación más una opción de compra sobre las acciones.
\item \textbf{Inflación}: la sufren todos los bonos, más en países de inflación aguda; los inversores exigen más rendimiento (menores precios). Se atenúa emitiendo en moneda dura o con cláusulas de indexación.
\item \textbf{Devaluación}: los bonos de emisores de países que devalúan pagan rendimientos más altos.
\item \textbf{\emph{Default}}: posibilidad de que el emisor no pague intereses o capital; lo evalúan las agencias calificadoras de crédito.
\item \textbf{Liquidez}: dificultad para vender un título sin un mercado amplio, frecuente en obligaciones de empresas sudamericanas.
\end{itemize}
\end{definicion}

\begin{definicion}{\emph{Yield-to-call} y \emph{yield-to-worst}}
Si el bono es rescatable, además de la TIR se calcula la \textbf{\emph{yield-to-call}}, la TIR que se obtendría si el rescate ocurriera en la fecha establecida; importa sobre todo cuando las tasas bajan y el cupón supera las tasas de mercado. Con varias fechas de rescate, la peor de esas TIR es la \textbf{\emph{yield-to-worst}}. Para el emisor, la decisión de rescatar considera además los ahorros impositivos y los gastos de emisión.
\end{definicion}

\begin{ejemplo}[Rescate anticipado de Santa Emilia]
Si la obligación pudiera rescatarse desde el inicio del segundo año, no se la rescataría mientras las tasas sigan en 11\,\%, pero sí si bajaran al 8\,\%. Si el precio de 96,3 anticipa ese rescate, la TIR de un flujo de dos años es la \emph{yield-to-call}:
\[
96{,}3 = \frac{10}{1+kd} + \frac{110}{(1+kd)^{2}}
\qquad\Longrightarrow\qquad kd = 12{,}19\,\%
\]
\end{ejemplo}

En los convertibles, el emisor suele reservarse el rescate con una agenda de precios decrecientes (\emph{call provision}); así puede forzar la conversión cuando el valor de conversión supera el de rescate, porque el tenedor prefiere acciones a recibir menos dinero.

%% ---------------------------------------------------------------------
\subsection{Valuación y rendimiento de la inversión en acciones}\label{sec:u4-acciones}

\paragraph{Cómo se negocian y cómo se valoran.} Las acciones nuevas se venden en el mercado primario y la mayor parte de la negociación ocurre en las bolsas, que son mercados secundarios \vease{sec:u1-alcance}. El inversor puede dar una \emph{orden de mercado}, que se ejecuta al mejor precio disponible, o una \emph{orden limitada}, que fija un precio tope. Las sociedades procuran que sus acciones coticen por encima del valor contable (activos menos pasivos), señal de que sus proyectos valen más de lo que cuestan.

Una acción da dos flujos: los dividendos y el precio de venta. Pero el precio de venta depende, a su vez, de los dividendos posteriores a la venta; por eso el valor de una acción para el inversor es el valor presente de todos los dividendos futuros esperados. Solo los dividendos —no las ganancias— llegan al accionista, y solo ellos deben descontarse. Muchas empresas en crecimiento no pagan dividendos durante años para reinvertir en proyectos que generarán efectivo más adelante; su valor descansa en los dividendos futuros.

\begin{definicion}{Tasa de capitalización de mercado (costo del capital propio)}
Es la tasa $ke$ (o $r$) a la que se descuentan los dividendos: el rendimiento que esperan los accionistas, compuesto por el dividendo más la apreciación esperada del precio. Es un costo de oportunidad, porque todos los títulos de una misma clase de riesgo se cotizan para ofrecer el mismo rendimiento esperado. Aquí se la toma como dato; su estimación se desarrolla en el capítulo~\ref{cap:u5} \vease{sec:u5-capm}.
\end{definicion}

\begin{formulaBox}[Modelo de descuento de dividendos]
\[
\begin{gathered}
P_0 = \frac{D_1 + P_1}{1+ke}
\qquad\qquad
ke = \frac{D_1 + P_1 - P_0}{P_0}
\\[4pt]
P_0 = \frac{D_1}{1+ke} + \frac{D_2}{(1+ke)^{2}} + \dots + \frac{D_n}{(1+ke)^{n}} + \frac{P_n}{(1+ke)^{n}}
\qquad\xrightarrow{\;n\to\infty\;}\qquad
P_0 = \sum_{t=1}^{\infty}\frac{D_t}{(1+ke)^{t}}
\end{gathered}
\]
\end{formulaBox}

\subsubsection{Valuación sobre la base de dividendos esperados sin crecimiento}

Si la empresa no retiene utilidades, no reinvierte y, por lo tanto, no crece ni cambia sus necesidades de capital de trabajo; si además la depreciación alcanza para renovar los bienes de uso, el dividendo es igual a la utilidad neta y forma una perpetuidad igual al dividendo del primer año: $P_0 = D_1/ke$.

\begin{ejemplo}[Calypso S.\,A.]
No crece y paga un dividendo anual de \$\,10 a perpetuidad; sus accionistas exigen el 20\,\%: $P_0 = 10/0{,}20 = 50$. Descontando el dividendo y el precio del año próximo se llega al mismo valor: $P_0 = (10 + 50)/1{,}20 = 50$.
\end{ejemplo}

\paragraph{Las ganancias de capital no son importantes en la perpetuidad.} Reemplazando cada precio futuro por el dividendo y el precio del período siguiente, el precio de hoy queda como una corriente de dividendos más un precio terminal muy lejano. A medida que el horizonte se alarga, el valor actual de los dividendos acumulados gana participación y el del precio futuro tiende a cero: el valor de la acción está dado, en lo fundamental, por el valor actual de sus dividendos. La tabla lo muestra para \conejemplos{Calypso}\sinejemplos{una acción que paga \$\,10 por año a perpetuidad, con $ke = 20\,\%$ y, por lo tanto, un precio de \$\,50}.

\begin{center}
\small
\begin{tabular}{@{}cccc@{}}
\toprule
Horizonte (años) & VA de los dividendos acumulados & VA del precio futuro & Total\\
\midrule
1 & 8,33 & 41,67 & 50\\
2 & 15,28 & 34,72 & 50\\
3 & 21,06 & 28,94 & 50\\
4 & 25,89 & 24,11 & 50\\
10 & 41,92 & 8,08 & 50\\
50 & 49,99 & 0,01 & 50\\
\bottomrule
\end{tabular}
\end{center}

\begin{figure}[H]
  \centering
  \includegraphics[width=0.5\linewidth,height=0.42\textheight,keepaspectratio]{Finanzas20corporativas203ed20Dumrauf-62}
  \caption{Composición relativa del valor de la acción según el horizonte: dividendos acumulados y precio futuro.}
\end{figure}

\subsubsection{Rendimientos de la inversión en acciones}

Como los bonos, las acciones rinden dividendos y ganancias de capital.

\begin{definicion}{Rendimiento por dividendos (\emph{dividend yield}) y ganancia de capital}
El \textbf{rendimiento por dividendos} (DY) relaciona el dividendo con el precio y es el equivalente del rendimiento corriente de un bono; en una acción sin crecimiento coincide exactamente con $ke$. La \textbf{ganancia de capital} relaciona el precio del año próximo con el actual. Su suma es el \textbf{rendimiento total}.
\end{definicion}

\begin{formulaBox}[Rendimiento de una acción]
\[
DY = \frac{D}{P_0}
\qquad
G = \frac{P_1 - P_0}{P_0}
\qquad
\text{Rendimiento total} = DY + G
\]
\end{formulaBox}

\begin{ejemplo}[Calypso S.\,A., un año después]
Si la acción sube a \$\,55: $DY = 10/50 = 20\,\%$, $G = (55-50)/50 = 10\,\%$ y el rendimiento total del año es del 30\,\%.
\end{ejemplo}

\subsubsection{Valuación de acciones con crecimiento constante}

Una empresa con muchas posibilidades de expansión reinvierte parte de sus utilidades y reparte el resto. Si sus ventas, costos y activos crecen a una tasa $g$, es razonable que los dividendos crezcan a la misma tasa, como en Telefonía Celular S.\,A., que crece al 5\,\% anual y reparte la mitad de su utilidad.

\begin{center}
\small
\begin{tabular}{@{}lrrrrr@{}}
\toprule
Telefonía Celular S.\,A. & Año 0 & Año 1 & Año 2 & Año 3 & Año 4\\
\midrule
Ventas & 100 & 105,0 & 110,3 & 115,8 & 121,6\\
Costos & 50 & 52,5 & 55,1 & 57,9 & 60,8\\
Impuestos & 20 & 21,0 & 22,1 & 23,2 & 24,3\\
Utilidad neta & 30 & 31,5 & 33,1 & 34,7 & 36,5\\
Dividendos (50\,\%) & 15 & 15,8 & 16,5 & 17,4 & 18,2\\
\bottomrule
\end{tabular}
\end{center}

La corriente de dividendos es una perpetuidad creciente \vease{sec:u1-va}, y su valor es la fórmula de Gordon.

\begin{formulaBox}[Modelo de Gordon (crecimiento constante)]
\[
P_0 = \frac{D_1}{1+ke} + \frac{D_1(1+g)}{(1+ke)^{2}} + \frac{D_1(1+g)^{2}}{(1+ke)^{3}} + \dots
    = \frac{D_1}{ke-g} = \frac{D_0\,(1+g)}{ke-g}
\qquad\qquad
ke = \frac{D_1}{P_0} + g
\]
Válido solo si $ke > g$. Con crecimiento estable, el costo del capital propio es el rendimiento por dividendos más la tasa de crecimiento.
\end{formulaBox}

\begin{ejemplo}[Telefonía Celular S.\,A.]
Con un dividendo próximo de \$\,10, $ke = 20\,\%$ y $g = 5\,\%$: $P_0 = 10/(0{,}20-0{,}05) = 66{,}66$, y a la inversa $ke = 10/66{,}66 + 0{,}05 = 20\,\%$. Si el dividendo no creciera valdría \$\,50: la diferencia de 16,66 es lo que el mercado paga por el crecimiento de los dividendos. (Con los dividendos de la tabla, $D_1 = 15{,}75$, el precio sería $15{,}75/0{,}15 = 105$.)
\end{ejemplo}

\paragraph{Cuestionamientos al modelo de Gordon.} Se le hacen tres críticas, y todas tienen respuesta.
\begin{enumerate}
\item \emph{Exige $ke > g$.} No es un problema real: a la larga ninguna empresa crece más rápido que la economía en la que opera, de modo que el crecimiento de largo plazo ronda, en términos reales, el 3 o 4\,\% (el del PBI, siempre de un dígito), mientras que $ke$ es siempre mayor. Las empresas que crecen fuerte al principio se valúan con un modelo por fases.
\item \emph{Los dividendos no crecen a una tasa estable.} Es cierto, pero $g$ es una media geométrica: la tasa compuesta a la que se espera que crezcan en promedio, aunque haya variaciones anuales.
\item \emph{Hay empresas que pagan dividendos muy irregularmente.} Para las que no tienen una política definida —situación frecuente en países sudamericanos— el modelo no es adecuado; sí puede serlo para empresas maduras con una política establecida.
\end{enumerate}

\paragraph{Estimación de la tasa de crecimiento.} La forma más sencilla multiplica la rentabilidad del capital propio por la tasa de retención: si el ROE es del 10\,\% y se retiene la mitad de las utilidades, esa mitad vuelve a ganar el 10\,\% y la empresa crece al $0{,}10\times 0{,}50 = 5\,\%$. La otra vía es la tasa histórica: el promedio geométrico, que solo usa el primer y el último dato, o un ajuste lineal o exponencial por mínimos cuadrados, que usa todas las observaciones.

\begin{formulaBox}[Tasa de crecimiento de los dividendos]
\[
\begin{gathered}
g = \text{ROE}\times b
  = \left(1 - \frac{DIV_1}{EPS_1}\right)\times\frac{EPS_1}{\text{Valor contable del patrimonio por acción}}
\\[4pt]
g_{\text{histórica}} = \left(\frac{D_n}{D_1}\right)^{1/(n-1)} - 1
\end{gathered}
\]
$b$ es la tasa de retención; el ROE se interpreta como el rendimiento de las utilidades retenidas.
\end{formulaBox}

\begin{ejemplo}[Tasa histórica de crecimiento]
\begin{center}
\begin{tabular}{@{}lccccc@{}}
\toprule
Año & 1 & 2 & 3 & 4 & 5\\
\midrule
Dividendo (\$) & 0,80 & 0,84 & 0,90 & 0,92 & 0,96\\
Variación anual & & 5,0\,\% & 7,0\,\% & 2,0\,\% & 4,3\,\%\\
\bottomrule
\end{tabular}
\end{center}
La tasa compuesta resuelve $0{,}80\,(1+g)^{4} = 0{,}96$: $g = (0{,}96/0{,}80)^{1/4} - 1 = 4{,}66\,\%$. Con el ajuste exponencial (función ESTIMACION.LOGARITMICA de Excel) el texto obtiene alrededor de 4,8\,\%, algo más que la media geométrica (con la serie exacta de la tabla ambos métodos dan $\approx 4{,}7\,\%$).
\end{ejemplo}

\begin{figure}[H]
  \centering
  \includegraphics[width=0.5\linewidth,height=0.42\textheight,keepaspectratio]{Finanzas20corporativas203ed20Dumrauf-71}
  \caption{Ajuste exponencial («mejor ajuste») de la serie de dividendos.}
\end{figure}

\subsubsection{Precio, utilidades y oportunidades de crecimiento}

Los inversores distinguen las \textbf{acciones de crecimiento}, cuya ganancia esperada proviene del aumento futuro de las utilidades, de las \textbf{acciones de renta}, que se compran por su dividendo regular.

\begin{definicion}{Valor presente de las oportunidades de crecimiento (PVGO)}
Es la parte del precio que excede el valor capitalizado de las utilidades bajo una política de crecimiento nulo, es decir, de una empresa que repartiera toda su utilidad sin reinvertir ($EPS_1/r$). Mide lo que el mercado paga hoy por los rendimientos superiores de inversiones todavía no realizadas; por eso las acciones de crecimiento cotizan con razones precio-utilidad altas.
\end{definicion}

\begin{formulaBox}[Precio, utilidades y PVGO]
\[
\begin{gathered}
P_0 = \frac{EPS_1}{r} + PVGO
\qquad\qquad
\frac{P_0}{EPS_1} = \frac{1}{r}\left[1 + \frac{PVGO}{EPS_1/r}\right]
\qquad\qquad
\frac{EPS_1}{P_0} = r\left(1 - \frac{PVGO}{P_0}\right)
\\[4pt]
P_0 = \frac{EPS_1}{r} + \frac{NPV_1}{r-g}
\end{gathered}
\]
La última es el modelo NPVGO: si la empresa invierte cada año en proyectos cuyo VPN por acción es $NPV_1$ en el primer año y crece a la tasa $g$, la PVGO es una perpetuidad creciente.
\end{formulaBox}

La PVGO solo es positiva si lo reinvertido rinde más que el costo del capital propio: crecer reinvirtiendo a una tasa igual a $r$ no agrega valor, y hacerlo a una tasa menor lo destruye. Por eso comparar el precio con crecimiento, $D_1/(ke-g)$, con el que tendría el mismo dividendo sin crecer, $D_1/ke$, mide el valor del crecimiento del dividendo, mientras que la PVGO compara el precio con $EPS_1/r$; el análisis del PER que sigue lo muestra con números.

\subsubsection{Valuación de acciones con fases de crecimiento variable}

Un buen negocio suele crecer mucho al principio y luego moderarse hacia el ritmo de la economía (PBI). La acción se valúa entonces en dos etapas: una fase de $t$ períodos de crecimiento $g$, cuyo valor es el de una anualidad creciente \vease{sec:u1-va}, y una perpetuidad que empieza en $t+1$ con crecimiento $g'$, actualizada $t$ períodos.

\begin{formulaBox}[Modelo de dos fases]
\[
P_0 = D_1\,\frac{1 - \left(\dfrac{1+g}{1+ke}\right)^{t}}{ke-g}
    + \frac{1}{(1+ke)^{t}}\,\frac{D_{t+1}}{ke-g'}
\qquad\qquad
D_{t+1} = D_1\,(1+g)^{t-1}\,(1+g')
\]
Con crecimiento diferencial sin una pauta simple, se descuenta cada dividendo por separado y se suman.
\end{formulaBox}

\begin{ejemplo}[Telefonía Celular S.\,A. en dos fases]
Crece al 5\,\% durante diez años y luego, por la competencia, los sustitutos y la saturación del mercado, al 2\,\%; $D_1 = 10$ y $ke = 20\,\%$. El dividendo del año 10 es $10\times 1{,}05^{9} = 15{,}51$ y el del año 11, $15{,}51\times 1{,}02 = 15{,}82$:
\[
P_0 = 10\times\frac{1-(1{,}05/1{,}20)^{10}}{0{,}20-0{,}05} + \frac{1}{1{,}20^{10}}\times\frac{15{,}82}{0{,}20-0{,}02}
    = 49{,}13 + 14{,}20 = 63{,}33
\]
\end{ejemplo}

\begin{figure}[H]
  \centering
  \includegraphics[width=0.5\linewidth,height=0.42\textheight,keepaspectratio]{Finanzas20corporativas203ed20Dumrauf-77}
  \caption{Crecimiento de los dividendos en dos fases: 5\,\% durante diez años y 2\,\% después.}
\end{figure}

\subsubsection{Valuación de las acciones por descuento del \emph{Equity Cash Flow}}

El modelo de dividendos usa una definición estricta del flujo del accionista: supone que solo recibe dividendos. El \emph{Equity Cash Flow} \vease{sec:u2-flujos} es más amplio: es el flujo residual después de inversiones, intereses y amortizaciones de deuda, es decir, lo que la empresa \emph{puede} pagar como dividendos. Descontarlo es una alternativa al modelo de dividendos, más explícita, porque cada partida surge de integrar resultados, inversión y financiamiento en un modelo proyectado. Como los dividendos pueden ser mayores o menores que el ECF, las dos estimaciones suelen diferir. Cuatro factores los separan: una política de dividendos estable, los requerimientos de inversión, los impuestos personales —que pueden hacer preferir ganancias de capital— y la señalización.

\begin{itemize}
\item \textbf{Dividendos iguales al ECF}: cuando la empresa reparte todo el flujo disponible después de inversiones, intereses y amortizaciones, o cuando reinvierte el ECF en proyectos con VAN nulo (títulos bien valuados o adquisiciones que no crean valor), porque esa reinversión no aumenta el flujo futuro del accionista.
\item \textbf{Dividendos menores que el ECF}: cuando reinvierte para crecer. Si reinvierte en títulos de corto plazo con VAN nulo o en adquisiciones con VAN nulo o negativo, crece en tamaño pero no en valor, y la menor razón de endeudamiento que resulta también puede restar valor.
\item \textbf{Dividendos mayores que el ECF}: la empresa debe emitir acciones o deuda para pagarlos, quizá como señal; el capital externo tiene costos de emisión, la razón de endeudamiento puede subir y el exceso puede generar racionamiento de capital si después no consigue financiar una expansión.
\end{itemize}

\subsubsection{Valuación de acciones preferidas}

La acción preferida \vease{sec:u3-flp} se valúa como el valor presente de sus dividendos preferidos, una perpetuidad descontada a $kp$.

\begin{formulaBox}[Acción preferida]
\[
P_p = \frac{D_{p}}{1+kp} + \frac{D_{p}}{(1+kp)^{2}} + \dots = \frac{D_p}{kp}
\qquad\qquad\text{por ejemplo: } \frac{10}{0{,}10} = 100
\]
\end{formulaBox}

\subsubsection{Valuación de acciones a partir del \emph{price earning}}

El \textbf{PER} (\emph{price earning ratio}), la razón precio-utilidad de §\,\ref{sec:u2-ratios}, es el múltiplo más usado como medida rápida de cómo valora el mercado una acción: indica cuántas veces cabe el beneficio por acción (BPA) en el precio, o cuántos pesos paga el inversor por cada peso de ganancia.

\begin{formulaBox}[PER y modelo de Gordon]
\[
\begin{gathered}
PER = \frac{P}{BPA} = \frac{P\times\text{N.º de acciones}}{\text{Utilidad neta}}
\qquad\qquad
P = BPA\times PER_{\text{sector}}
\\[4pt]
PER = \frac{P}{BPA} = \frac{DPA}{(ke-g)\,BPA} = \frac{\text{Payout}}{ke-g}
\end{gathered}
\]
\end{formulaBox}

\begin{ejemplo}[Omega S.\,A.]
Ganó \$\,3.000.000 con 10.000 acciones en circulación (BPA = \$\,300) y su acción cotiza a \$\,2.400: $PER = 2.400/300 = 8$. Los inversores pagan 8 pesos por cada peso de ganancia.
\end{ejemplo}

Según la fórmula, un PER alto (bajo) puede indicar que los inversores esperan un crecimiento alto (bajo) de los dividendos, que la acción tiene poco riesgo y por eso $ke$ es baja, o que se espera un índice de reparto alto. El primero y el tercero interactúan: repartir más significa crecer menos. Desde otra perspectiva, el PER depende de tres factores: el valor por acción de las oportunidades de crecimiento (más oportunidades, PER más alto), el riesgo (más riesgo, PER más bajo) y el método contable (los contadores conservadores, al informar utilidades menores, elevan el PER).

Para estimar el valor de una acción se multiplica su BPA por el PER de empresas comparables o del sector. Se usa mucho porque da una medida normalizada para comparar precios —cuánto paga el mercado por cada peso de beneficio— y porque es más sencillo que el modelo de Gordon y sirve para acciones que no pagan dividendos. A veces se lo lee como una medida de la calidad y el aprecio de los beneficios de la firma, pero también se presta a usos impropios:

\begin{center}
\small
\begin{tabular}{@{}p{0.45\linewidth}p{0.45\linewidth}@{}}
\toprule
Ventajas & Desventajas\\
\midrule
Simple de calcular y disponible para las empresas que cotizan; facilita la comparación entre acciones. & Puede creerse que elimina la necesidad de hacer hipótesis sobre riesgo, crecimiento y dividendos.\\
\addlinespace
Relaciona el precio con los beneficios actuales. & Puede creerse que refleja mejor los humores y percepciones del mercado.\\
\addlinespace
Sustituye otras características de la firma: riesgo, rendimiento y crecimiento. & Depende de las convenciones contables con que se calcula el beneficio.\\
\bottomrule
\end{tabular}
\end{center}

\paragraph{Factores que afectan al PER: ROE, crecimiento y rentabilidad exigida.} Influyen el ROE, el crecimiento —ligado a la política de dividendos—, la rentabilidad exigida y las tasas de interés. En el cuadro siguiente, seis empresas financiadas solo con acciones tienen un patrimonio de 100 millones y $ke = 10\,\%$; se suponen precios iguales al valor descontado de los dividendos.

\begin{center}
\small
\begin{tabular}{@{}lrrrrrr@{}}
\toprule
Empresa & A & B & C & D & E & F\\
\midrule
Utilidad neta (ROE) & 10 (10\,\%) & 12 (12\,\%) & 14 (14\,\%) & 16 (16\,\%) & 18 (18\,\%) & 8 (8\,\%)\\
\addlinespace
\multicolumn{7}{@{}l}{\emph{Sin retención: reparto del 100\,\%, $g = 0$}}\\
Valor de mercado & 100 & 120 & 140 & 160 & 180 & 80\\
PER & 10 & 10 & 10 & 10 & 10 & 10\\
\addlinespace
\multicolumn{7}{@{}l}{\emph{Con retención del 50\,\%: $g = \text{ROE}\times 0{,}50$}}\\
Dividendos & 5 & 6 & 7 & 8 & 9 & 4\\
Crecimiento $g$ & 5\,\% & 6\,\% & 7\,\% & 8\,\% & 9\,\% & 4\,\%\\
Valor de mercado & 100 & 150 & 233 & 400 & 900 & 67\\
PER & 10,0 & 12,5 & 16,7 & 25,0 & 50,0 & 8,3\\
\bottomrule
\end{tabular}
\end{center}

El valor de mercado se obtiene descontando los dividendos con $ke$ (con Gordon, por ejemplo, B vale $6/(0{,}10-0{,}06) = 150$) y el PER, dividiéndolo por la utilidad neta. En términos de PVGO ($P - \text{utilidad}/ke$), la retención crea 30 en B, 93 en C, 240 en D y 720 en E, nada en A y destruye 13 en F. Las conclusiones son cuatro:
\begin{enumerate}
\item El crecimiento esperado $g$ es el de los beneficios y los dividendos, y depende de la tasa de retención. No basta un ROE alto para tener un PER alto si no hay crecimiento, ni basta crecer si se invierte por debajo del costo de los recursos.
\item Si la empresa reparte todo y no crece, el PER es constante (aquí, $1/ke = 10$), aunque el ROE supere el costo del capital.
\item Si retiene y crece, el PER aumenta con el ROE, tanto más cuanto mayor es el crecimiento; pero baja cuando el ROE es inferior al costo de los recursos (F, 8,3), porque esa empresa destruye valor al reinvertir.
\item Cuanto mayor es la rentabilidad exigida a las acciones, menor es el PER.
\end{enumerate}

\begin{formulaBox}[Regla del PER]
\[
g = 0 \;\Longrightarrow\; PER = \frac{1}{ke}\ \text{constante}
\qquad\qquad
PER\uparrow \iff g > 0 \ \text{y}\ \text{ROE} > ke
\]
Si se cumple la condición, el PER sube tanto más cuanto mayor es el crecimiento.
\end{formulaBox}

\paragraph{Valoración de la empresa por flujo de efectivo descontado.} En lugar de los dividendos por acción puede descontarse el flujo de efectivo libre de toda la empresa con un modelo de dos etapas: se proyectan y descuentan los flujos libres hasta un horizonte y se suma el valor descontado de un valor de horizonte \vease{sec:u3-ccpp}.

%% ---------------------------------------------------------------------
\section{Determinantes del costo de oportunidad del dinero}\label{sec:u4-costooport}

La tasa de interés es el precio que se paga por «alquilar» dinero durante un tiempo. Quienes ahorran lo prestan, a cambio de un interés, a quienes consumen más de lo que ganan: unos cambian consumo de hoy por más consumo futuro y otros aumentan su poder de compra actual pagando con consumo futuro. Además del interés «puro», que remunera el valor tiempo del dinero, el costo de oportunidad depende de cuatro factores.

\begin{definicion}{Componentes del costo de oportunidad del dinero}
\begin{itemize}
\item \textbf{Riesgo}: a mayor riesgo, mayor rendimiento exigido. Incluye el \emph{riesgo de default} \vease{sec:u4-bonos}, factor crítico en las tasas de las obligaciones porque depende de la capacidad de pago del emisor (como referencia sin riesgo de impago se toman los títulos del Tesoro de EE.\,UU.), el \emph{riesgo de variabilidad del precio} (aun sin riesgo de impago, cuando cambian las tasas cambian los precios de bonos y acciones) y la diferencia de \emph{derechos}: la deuda obliga a pagar bajo pena de consecuencias legales, mientras que al accionista no se le promete dividendo alguno, por lo que su costo de oportunidad suele ser mayor que el del obligacionista.
\item \textbf{Inflación esperada}: los inversores se fijan en el poder adquisitivo, y los rendimientos incorporan la inflación esperada \vease{sec:u1-va}. Si los bancos no la convalidaran, los ahorristas retirarían su dinero para comprar bienes, que suben más; y las empresas que piden prestado pueden pagar tasas nominales más altas porque también suben sus precios de venta.
\item \textbf{Impuestos}: los inversores calculan rendimientos después de impuestos. Para la empresa, como los intereses se deducen del impuesto a las ganancias, el costo de la deuda es $kd\,(1-t)$; para el prestamista, si suben los impuestos, debe cobrar una tasa mayor antes de impuestos para mantener su rendimiento neto.
\item \textbf{Vencimiento}: según la estructura temporal \vease{sec:u4-bonos}, normalmente ascendente, los pagos más lejanos se descuentan a tasas más altas, porque los títulos de vencimiento largo son más sensibles a los cambios de tasas y ese mayor riesgo de precio exige mayor rendimiento.
\end{itemize}
\end{definicion}

%% ---------------------------------------------------------------------
\begin{control}
\pregunta{1. ¿Cuál es el tema principal?}
Cómo afecta al valor la decisión de distribuir efectivo a los accionistas y cómo se valúan los títulos que emite la empresa —bonos y acciones— como valor presente de sus flujos.

\pregunta{2. ¿Dónde y cómo se desarrolla?}
Parte de los tipos y la mecánica de los dividendos y de la proposición de indiferencia de MM; incorpora impuestos, recompras, señales, clientelas y la evidencia sobre cómo deciden las empresas, y ordena el debate en tres posiciones. Luego aplica el descuento de flujos a las obligaciones (precio, rendimientos, convenciones, estructura temporal y riesgos) y a las acciones (modelos de dividendos sin crecimiento, de Gordon y por fases, ECF, preferidas y PER), y cierra con los determinantes de la tasa de descuento.

\pregunta{3. ¿Cuáles son los temas secundarios?}
Fechas del dividendo; dividendos caseros; formas de recompra; contenido informativo y efecto clientela; suavizamiento (Lintner); desdoblamientos y contrasplits; ciclo de vida; valor residual, técnico y paridad; rendimiento corriente y total; costo efectivo de una obligación; tasas contado, futuras y \emph{bootstrapping}; \emph{yield-to-call}; estimación de $g$; PVGO; ventajas y límites del PER.

\pregunta{4. ¿Las ideas principales y secundarias están relacionadas entre sí?}
Sí. La política de dividendos solo importa por las imperfecciones del mercado —impuestos, costos de emisión, agencia, información—, porque en mercados perfectos el valor depende de la inversión. Ese mismo principio sostiene la valuación: un bono o una acción valen el flujo que prometen descontado al costo de oportunidad, y ese costo refleja riesgo, inflación, impuestos y plazo. El crecimiento de los dividendos y el PER vinculan ambas partes: retener solo crea valor si la reinversión rinde más que $ke$.

\pregunta{5. ¿Cuál es la conclusión o la idea general?}
En mercados perfectos la forma de repartir no cambia el valor; en el mundo real, impuestos, agencia y señales explican por qué las empresas suavizan dividendos, recompran acciones y ajustan el reparto a su ciclo de vida. El precio de cualquier título es el valor presente de su flujo al rendimiento de inversiones de riesgo similar, y el crecimiento solo agrega valor cuando lo reinvertido supera el costo de oportunidad.
\end{control}

%% =====================================================================
%%  UNIDAD 5 (parte A): VPN, criterios de inversión y evaluación de proyectos
%% =====================================================================
\chapter{Decisiones de inversión, riesgo y costo del capital}\label{cap:u5}

\begin{ficha}
\fichabase{5 — Decisiones de inversión, riesgo y costo del capital}
\fila{Texto base}{R. Brealey, S. Myers y F. Allen, \emph{Principios de finanzas corporativas}, McGraw-Hill: cap.~5 «Valor presente neto y otros criterios de inversión» (resumen).}
\fila{}{S. Ross, R. Westerfield y J. Jaffe, \emph{Finanzas corporativas}, McGraw-Hill: esquemas de repaso «Rendimiento y riesgo: el CAPM» y «Una visión alternativa del riesgo y el rendimiento: el APT».}
\fila{}{G.\,L. Dumrauf, \emph{Finanzas corporativas: un enfoque latinoamericano}, 3.ª ed., Alfaomega: cap.~7 «Riesgo y rendimiento», cap.~8 «El CAPM, el costo de capital en mercados emergentes y el APM» y cap.~12 «El costo del capital».}
\fila{}{Apuntes de clase «Gestión de riesgos financieros» (FCA): valor y riesgo de un portafolio de acciones, retornos, volatilidad y VaR, portafolios de dos y de $N$ activos, frontera eficiente, \emph{bootstrap}, ratio de Sharpe, betas de mercado y riesgo sistémico.}
\fila{Complementaria}{Brealey, Myers y Allen, ob. cit.: resúmenes «Toma de decisiones de inversión con la regla del valor presente neto» y «El riesgo y el costo de capital».}
\fila{}{Ross, Westerfield y Jaffe, ob. cit.: resúmenes «Algunas reglas de inversión alternativas», «VPN y presupuesto de capital», «Estrategia y análisis en el uso del VPN», «Teoría del mercado de capitales», «Riesgo, costo de capital y presupuesto de capital» y «Decisiones de financiamiento corporativo y mercado de capitales eficiente»; preguntas de repaso sobre presupuesto de capital.}
\fila{}{Apuntes de cátedra: evaluación de proyectos y flujos de caja (flujos incrementales, valor terminal, EVA, MVA y ROIC), determinantes del beta y reducción del costo de capital.}
\fila{Fórmulas previas}{VAN y valor actual \vease{sec:u1-va}; FCFF \vease{sec:u1-fcff}; FCF, ECF y CCF \vease{sec:u2-flujos}; WACC, APV y betas apalancadas \vease{sec:u3-ccpp}; escenarios y punto de equilibrio \vease{sec:u3-escenarios}; leverage \vease{sec:u3-leverage}.}
\end{ficha}

%% ---------------------------------------------------------------------
\section{Valor presente neto y otros criterios de inversión}\label{sec:u5-vpn}

\subsection{Repaso de lo básico}

Evaluar una inversión propuesta exige cuatro pasos: pronosticar los flujos de fondos que el proyecto generará a lo largo de su vida; determinar el costo de oportunidad del capital $r$, que debe reflejar el valor tiempo del dinero y el riesgo; descontar con él los flujos futuros para obtener su valor presente, y restar la inversión neta para obtener el valor presente neto \vease{sec:u1-va}.

\begin{definicion}{Propiedades de la regla del VPN}
\begin{itemize}
\item Reconoce que un peso hoy vale más que un peso mañana; ninguna regla que ignore el valor tiempo del dinero puede ser sensata.
\item Depende únicamente de los flujos pronosticados del proyecto y del costo de oportunidad del capital. Las reglas que dependen de otros factores (gustos de los directivos, método contable, rentabilidad de otros negocios) generan conflictos de intereses y decisiones inferiores.
\item Los VPN se suman: $VPN(A+B) = VPN(A) + VPN(B)$. Por eso «empaquetar» un proyecto malo con uno bueno no lo vuelve aceptable, aunque el paquete pueda inducir a los directivos a aceptar malas inversiones.
\end{itemize}
Según Ross, el VPN tiene tres atributos que las reglas alternativas no reúnen: usa flujos de efectivo y no ganancias, usa todos los flujos del proyecto y los descuenta correctamente.
\end{definicion}

\begin{figure}[H]
  \centering
  \includegraphics[width=0.5\linewidth,height=0.42\textheight,keepaspectratio]{bma process-11}
\end{figure}
\begin{figure}[H]
  \centering
  \includegraphics[width=0.5\linewidth,height=0.42\textheight,keepaspectratio]{bma process-12}
\end{figure}
\begin{figure}[H]
  \centering
  \includegraphics[width=0.5\linewidth,height=0.42\textheight,keepaspectratio]{bma process-13}
\end{figure}
\begin{figure}[H]
  \centering
  \includegraphics[width=0.5\linewidth,height=0.42\textheight,keepaspectratio]{bma process-14}
  \caption{La regla del VPN y los criterios de inversión: ilustraciones del capítulo de Brealey, Myers y Allen.}
\end{figure}

Al informar a los accionistas se presentan también ganancias y activos contables, y allí importa si un egreso de caja se registra como \textbf{inversión de capital} —figura como activo y se deprecia— o como \textbf{gasto operativo}, que se deduce de inmediato de las ganancias. Para el VPN, en cambio, solo cuenta cuándo sale y entra el efectivo.

\subsection{Período de recupero}

\begin{definicion}{Período de recupero (\emph{payback})}
Es el número de años que tarda el flujo de fondos acumulado en igualar la inversión inicial. La regla acepta el proyecto si ese plazo es menor que un período de corte prefijado. El \textbf{período de recupero descontado} aplica la misma regla a los flujos descontados.
\end{definicion}

Las empresas lo usan porque es la forma más sencilla de comunicar una idea de la rentabilidad, porque el éxito de corto plazo acelera el ascenso de los directivos y porque un acceso limitado al capital restringe la capacidad de financiarse en el futuro. Pero da respuestas engañosas: ignora todos los flujos posteriores a la fecha de corte, da el mismo peso a todos los anteriores (no descuenta) y el plazo de corte es un estándar arbitrario. El recupero descontado corrige solo el segundo defecto.

\subsection{Tasa de rendimiento contable promedio}

\begin{formulaBox}[Tasa de rendimiento contable]
\[
\begin{gathered}
\text{Rendimiento contable} = \frac{\text{Beneficio contable}}{\text{Activos contables}}
\\[4pt]
TRCP = \frac{\text{Utilidad promedio después de impuestos y depreciación}}{\text{Valor contable promedio de la inversión}}
\end{gathered}
\]
\end{formulaBox}

Tiene tres defectos: usa la utilidad neta en lugar del flujo de efectivo, no tiene en cuenta el momento en que ocurren los resultados y no orienta sobre cuál debería ser la tasa objetivo.

\subsection{Tasa interna de retorno}

\begin{definicion}{Tasa interna de retorno (TIR o tasa de flujo de fondos descontado)}
Es la tasa de descuento que hace cero el VPN del proyecto. La regla acepta el proyecto si el costo de oportunidad del capital es menor que la TIR; da la misma respuesta que el VPN siempre que este sea una función suavemente decreciente de la tasa de descuento, como ocurre en el caso normal de una inversión inicial seguida solo de ingresos. A mano, se calculan tres o cuatro pares VPN-tasa, se grafican y se busca la tasa con VPN nulo.
\end{definicion}

\begin{formulaBox}[Tasa interna de retorno]
\[
VPN = C_0 + \frac{C_1}{1+TIR} + \frac{C_2}{(1+TIR)^{2}} + \dots + \frac{C_T}{(1+TIR)^{T}} = 0
\]
\end{formulaBox}

La regla tiene cuatro trampas:
\begin{enumerate}
\item \textbf{Prestar o pedir prestado.} Al prestar se busca una tasa alta y al endeudarse una baja: si el proyecto tiene ingresos iniciales seguidos de egresos, la regla se invierte y hay que aceptar cuando la TIR es \emph{menor} que la tasa de descuento.
\item \textbf{Tasas múltiples.} Puede haber tantas TIR como cambios de signo en el flujo; en ese caso debe usarse el VPN.
\item \textbf{Proyectos mutuamente excluyentes} (formas alternativas de hacer el mismo trabajo o de usar una misma instalación). Por diferencias de tamaño o de calendario, el proyecto con mayor TIR no necesariamente tiene el mayor VPN; la TIR solo sirve aplicada a los flujos incrementales entre ambos.
\item \textbf{Más de un costo de oportunidad.} Si la tasa de descuento cambia según el plazo, habría que comparar la TIR con un promedio ponderado de esas tasas para obtener una cifra comparable.
\end{enumerate}

\subsection{Elección de inversiones cuando los recursos son limitados}

Muchas empresas no pueden aceptar todos los proyectos con VPN positivo. Ese problema es el \textbf{racionamiento de capital}: hay que elegir el paquete de proyectos que dé el mayor VPN posible dentro de los recursos disponibles.

\begin{definicion}{Índice de rentabilidad}
Ordena los proyectos según el valor que crean por peso invertido. En Ross se define como el valor presente de los flujos posteriores a la inversión dividido por la inversión inicial, y un proyecto independiente se acepta si supera 1; en Brealey, Myers y Allen, como el VPN dividido por la inversión, que debe ser positivo. Ambas versiones ordenan igual. Con fondos limitados no se ordena por VPN sino por este índice; para proyectos mutuamente excluyentes, en cambio, el índice falla porque ignora las diferencias de escala.
\end{definicion}

\begin{formulaBox}[Índice de rentabilidad]
\[
IR = \frac{VP \text{ de los flujos posteriores}}{\text{Inversión inicial}}\ (\text{aceptar si} > 1)
\qquad
IR = \frac{VPN}{\text{Inversión}}\ (\text{aceptar si} > 0)
\]
\end{formulaBox}

Hay dos modelos de racionamiento. En el \textbf{racionamiento blando}, la dirección fija límites presupuestarios para que cada división establezca sus propias prioridades: es un antídoto contra los pronósticos de flujos sesgados por el optimismo. En el \textbf{racionamiento duro}, imperfecciones del mercado levantan una barrera entre la empresa y los mercados de capitales; mientras esa barrera sea la única imperfección, el VPN sigue siendo el criterio válido.

\begin{figure}[H]
  \centering
  \includegraphics[width=0.5\linewidth,height=0.42\textheight,keepaspectratio]{bma process-15}
\end{figure}
\begin{figure}[H]
  \centering
  \includegraphics[width=0.5\linewidth,height=0.42\textheight,keepaspectratio]{bma process-16}
\end{figure}
\begin{figure}[H]
  \centering
  \includegraphics[width=0.5\linewidth,height=0.42\textheight,keepaspectratio]{bma process-17}
  \caption{Racionamiento de capital: selección del paquete de proyectos de mayor VPN (Brealey, Myers y Allen).}
\end{figure}

%% ---------------------------------------------------------------------
\subsection{Toma de decisiones de inversión con la regla del VPN}

Además de instalaciones productivas, las empresas invierten en activos intangibles —conocimiento, investigación y desarrollo, tecnología de la información—, cuyos flujos son difíciles de estimar. En todos los casos valen tres reglas generales.

\begin{definicion}{Las tres reglas para estimar los flujos de un proyecto}
\begin{enumerate}
\item \textbf{Solo importa el flujo de efectivo}, la diferencia entre el efectivo cobrado y el pagado, y no la utilidad contable. Un edificio comprado hoy en \$\,100.000 es una salida inmediata por ese importe, aunque con depreciación lineal a veinte años solo se registren \$\,5.000 de gasto por año.
\item \textbf{Los flujos deben ser incrementales}: interesa la diferencia entre los flujos de la empresa con el proyecto y sin él, incluidos los efectos sobre el resto del negocio (por ejemplo, ingresos posventa), el capital de trabajo mínimo requerido y los valores de liquidación al final; los costos hundidos no cuentan y los gastos generales solo se imputan si el proyecto los aumenta.
\item \textbf{La inflación debe tratarse con coherencia}: flujos nominales con tasa nominal y flujos reales con tasa real; nunca se mezclan.
\end{enumerate}
\end{definicion}

\paragraph{El capital de trabajo en los proyectos.} Sus tres componentes principales son los inventarios (incluido el efectivo), las cuentas por cobrar y las cuentas por pagar. Surge porque se compran insumos antes de vender, se mantiene efectivo como amortiguador ante gastos imprevistos y se vende a crédito; es una salida de fondos porque inmoviliza en el proyecto efectivo generado en otra parte de la empresa. Al estimar flujos se cometen tres errores típicos: olvidarlo, olvidar que cambia a lo largo de la vida del proyecto y olvidar que se recupera al final.

\paragraph{Momento óptimo de la inversión.} Un VPN positivo no obliga a invertir de inmediato. Con flujos ciertos, se calcula el valor futuro neto de cada fecha posible de inicio $t$ y se lo descuenta al presente; conviene la fecha que aporta más valor hoy.

\begin{formulaBox}[VPN según la fecha de inicio y costo anual equivalente]
\[
\begin{gathered}
VPN \text{ de invertir en } t = \frac{VFN_t}{(1+r)^{t}}
\\[4pt]
\text{Flujo anual equivalente} = \frac{VP \text{ de la inversión (o de los costos)}}{\text{Factor de anualidad a } T \text{ años y tasa } r}
\end{gathered}
\]
\end{formulaBox}

\begin{definicion}{Flujo de efectivo anual equivalente}
Es la anualidad que alcanza para recuperar una inversión de capital, incluido su costo de oportunidad, a lo largo de la vida económica del activo. Invierte el cálculo habitual: transforma una inversión actual en una corriente equivalente de flujos futuros. Sirve para comparar alternativas de vidas distintas: entre dos máquinas de vida desigual se elige con el enfoque del \textbf{ciclo de coincidencia} (repetir ambas hasta un horizonte común) o con el del \textbf{costo anual equivalente}, dos presentaciones de la misma información que deben preferir la misma máquina. El costo anual equivalente reaparece en el punto de equilibrio financiero \vease{sec:u3-escenarios}.
\end{definicion}

%% ---------------------------------------------------------------------
\subsection{Evaluación de proyectos y flujos de caja}

Evaluar una inversión de largo plazo implica cinco pasos: estimar los flujos incrementales, agregarles los efectos colaterales y los costos especiales, incorporar los costos de oportunidad, fijar la tasa de descuento según el riesgo y calcular el valor presente. Se trabaja en base caja porque el inversor solo dispone del dinero cuando lo cobra. Las finanzas corporativas usan flujos; la contabilidad, ganancias.

\paragraph{El enfoque «con y sin».} Hay que comparar la empresa \emph{con} el proyecto contra la empresa \emph{sin} él, nunca «antes» contra «después», porque esta última comparación ignora los costos de oportunidad.

\begin{ejemplo}[Costo de oportunidad de un terreno y reemplazo de una máquina]
Una empresa usará para expandirse un terreno que podría vender en \$\,100.000. Antes y después del proyecto sigue teniendo el terreno: el cambio en el flujo sería cero. Con el criterio correcto, con el proyecto usa el terreno y no cobra nada, y sin él lo vende y cobra 100.000: el proyecto debe cargar un costo de oportunidad de $0 - 100.000 = -100.000$. Del mismo modo, si se reemplaza una máquina vieja que puede revenderse en \$\,20 por una nueva de \$\,100, la inversión incremental es de \$\,80.
\end{ejemplo}

\begin{definicion}{Costos hundidos, costos de oportunidad y efectos secundarios}
\begin{itemize}
\item \textbf{Costo hundido}: erogación pasada e irreversible que ninguna decisión puede cambiar; se ignora. Si General Milk pagó el año pasado \$\,100.000 a una consultora por una prueba de mercado de su leche chocolatada, ese gasto es «leche derramada» y no pesa en la decisión actual.
\item \textbf{Costo de oportunidad}: lo que se deja de ganar al destinar al proyecto un activo que podría venderse, arrendarse o usarse en otra parte del negocio; se incluye aunque no implique una salida de caja.
\item \textbf{Efectos secundarios}: flujos que el proyecto provoca en el resto de la empresa. El más importante es la \textbf{erosión}, el flujo que el nuevo producto capta de clientes y ventas de otros productos de la misma empresa. También puede cambiar la demanda de productos complementarios, aliviar cuellos de botella —si dos procesos encadenados producen 300 y 600 unidades por hora, agregar máquinas al más lento destraba el conjunto— o mejorar la imagen de la empresa y la demanda de otras líneas.
\end{itemize}
\end{definicion}

También se excluyen las partidas que no son egresos reales de fondos (como las amortizaciones) y los gastos generales que no cambian por el proyecto. Los aumentos de capital de trabajo y de bienes de uso son egresos, y sus disminuciones, ingresos. Los impuestos se computan en base caja y según las reglas fiscales: hay dos juegos de libros —los impositivos, que siguen las normas del fisco, y los de los accionistas, que siguen las normas contables (que, por ejemplo, tratan distinto los intereses de bonos municipales)— y para el flujo interesan los primeros, porque el impuesto es una salida de caja.

\paragraph{Flujo del proyecto y flujo del accionista.} A la empresa —y al banco que le presta— le interesa ante todo si el proyecto es bueno, y trata su financiamiento como un problema separado. Por eso el \textbf{\emph{free cash flow} del proyecto} incluye los flujos operativos atribuibles a él, excluye las fuentes de financiamiento y resta los impuestos que corresponderían al resultado operativo, sin computar el ahorro fiscal de la deuda: en la práctica los flujos se calculan como si el proyecto se financiara solo con capital propio y el efecto de la deuda se refleja en la tasa de descuento \vease{sec:u3-metodos}. Si se evalúa desde el \textbf{\emph{equity cash flow}}, en cambio, se hacen dos ajustes: en la inversión inicial solo figura el aporte propio (se resta la parte financiada con deuda) y de los flujos siguientes se restan intereses y devoluciones de capital. La construcción del FCF desde el EBIT o desde el NOPAT, y su diferencia con la utilidad neta, se vieron en §\,\ref{sec:u1-fcff} y §\,\ref{sec:u2-flujos}.

\paragraph{Horizonte y valor continuo.} La proyección explícita debe extenderse hasta que la empresa alcance un \emph{estado estacionario}, es decir, hasta que la competencia iguale el rendimiento de las nuevas inversiones con el costo de capital. Después se recurre a un \textbf{valor continuo} (terminal), calculado como perpetuidad con o sin crecimiento, o a un valor de liquidación neto de impuestos.

\begin{formulaBox}[Valor de la firma con valor continuo]
\[
\begin{gathered}
V_C = \frac{FCF_{T+1}}{k_u}\quad\text{o}\quad V_C = \frac{FCF_{T+1}}{k_u - g}
\qquad\qquad
V = \sum_{t=1}^{T}\frac{FCF_t}{(1+k_u)^{t}} + \frac{FCF_{T+1}}{(k_u-g)(1+k_u)^{T}}
\\[4pt]
V = \underbrace{\frac{FCF_1}{1+WACC} + \dots + \frac{FCF_H}{(1+WACC)^{H}}}_{\text{VP del flujo libre}}
  + \underbrace{\frac{1}{(1+WACC)^{H}}\,\frac{FCF_{H+1}}{WACC-g}}_{\text{VP del valor horizonte}}
\end{gathered}
\]
$k_u$ es el costo del capital de la empresa sin deuda; con el WACC \vease{sec:u3-ccpp} la estructura es la misma (con $g = 0$, $V_C = FCF_{T+1}/WACC$).
\end{formulaBox}

\paragraph{Flujos reales y nominales.} Descontar flujos reales a la tasa real equivale a descontar flujos nominales a la tasa nominal, porque la inflación afecta al numerador y al denominador a la vez; la relación entre tasas es la ecuación de Fisher \vease{sec:u1-va}. Ignorar la inflación —por ejemplo, proyectar flujos nominales y descontarlos a una tasa real— infla artificialmente el VPN o la TIR.

\begin{formulaBox}[Coherencia entre flujos y tasas]
\[
\frac{FF_{\text{reales}}\,(1+\pi)}{(1+k_r)\,(1+\pi)} = \frac{FF_{\text{nominales}}}{1+k}
\qquad\qquad
k_r = \frac{1+k}{1+\pi} - 1 \;\approx\; k - \pi
\]
$\pi$ es la inflación, $k$ la tasa nominal y $k_r$ la real.
\end{formulaBox}

%% ---------------------------------------------------------------------
\subsection{EVA, MVA y los conductores del valor}

\begin{definicion}{NOPAT, capital invertido y ROIC}
El \textbf{NOPAT} (\emph{net operating profit after taxes}) es el resultado operativo después de impuestos: ventas menos gastos operativos menos impuestos ajustados. El \textbf{capital invertido} $C$ puede leerse por el lado operativo o por el financiero, y ambos coinciden. El \textbf{ROIC} es el rendimiento del capital invertido: NOPAT sobre $C$.
\end{definicion}

\begin{center}
\small
\begin{tabular}[t]{@{}ll@{}}
\toprule
\multicolumn{2}{@{}c@{}}{\textbf{NOPAT}}\\
\midrule
Desde el EBIT & Desde la utilidad neta\\
\midrule
EBIT & Utilidad neta\\
$-$ Impuestos sobre el EBIT & $+$ Intereses pagados $\times(1-t)$\\
 & $-$ Intereses ganados $\times(1-t)$\\
 & $+$ Egresos no operativos $\times(1-t)$\\
 & $-$ Ingresos no operativos $\times(1-t)$\\
\bottomrule
\end{tabular}
\hfill
\begin{tabular}[t]{@{}ll@{}}
\toprule
\multicolumn{2}{@{}c@{}}{\textbf{Capital invertido}}\\
\midrule
Lado operativo & Lado financiero\\
\midrule
Capital de trabajo & Deuda financiera\\
$+$ Activo fijo neto & $+$ Patrimonio neto\\
$-$ Otros activos operativos & \\
\bottomrule
\end{tabular}
\end{center}

\begin{definicion}{Valor económico agregado (EVA) y valor de mercado agregado (MVA)}
El \textbf{EVA} es el excedente del NOPAT sobre el costo del capital empleado; equivalentemente, el diferencial entre ROIC y WACC multiplicado por el capital invertido. El \textbf{MVA} es la riqueza creada para los accionistas: la diferencia entre el valor de mercado del capital y el capital aportado; es igual al valor actual de todo el EVA futuro, de modo que el valor de mercado de la empresa es el capital más el MVA, y el MVA refleja el EVA actual más la mejora esperada del EVA.
\end{definicion}

\begin{formulaBox}[EVA, MVA y valuación por EVA]
\[
\begin{gathered}
EVA = NOPAT - WACC\times C = (ROIC - WACC)\times C
\qquad\qquad
ROIC = \frac{NOPAT}{C}
\\[4pt]
MVA = E - E_c = \sum_{t=1}^{\infty}\frac{EVA_t}{(1+WACC)^{t}}
\\[4pt]
V = C + \sum_{j=1}^{T}\frac{C\,(ROIC-WACC)}{(1+WACC)^{j}} + \frac{C\,(ROIC-WACC)}{WACC\,(1+WACC)^{T}}
\\[4pt]
\text{con } EVA \text{ creciente a la tasa } g \text{ e inversión anual } I = g\,C\text{:}
\\[2pt]
V = C + \frac{C\,(ROIC-WACC)}{WACC-g} = \frac{C\,ROIC - g\,C}{WACC-g} = \frac{NOPAT - I}{WACC-g} = \frac{FCF}{WACC-g}
\end{gathered}
\]
$E$ es el valor de mercado del capital y $E_c$ el capital aportado por los inversores. La última igualdad muestra que la valuación por EVA coincide con el descuento del FCF.
\end{formulaBox}

\begin{ejemplo}[EVA y ROIC]
Con un capital invertido de \$\,10.000, ROIC del 15\,\% y WACC del 10\,\%: $EVA = (0{,}15-0{,}10)\times 10.000 = \$\,500$. Una empresa con NOPAT de 7.784 y capital invertido de 66.149 tiene un ROIC de $7.784/66.149 \approx 11{,}8\,\%$.
\end{ejemplo}

El EVA señala cuatro palancas para crear valor, porque $EVA = (NOPAT/C - WACC)\times C$: \textbf{operar}, mejorando el rendimiento del capital existente; \textbf{financiar}, reduciendo el costo de capital; \textbf{invertir}, mientras el rendimiento marginal supere el costo, y \textbf{racionalizar}, desinvirtiendo el capital que no rinde lo suficiente.

\begin{formulaBox}[Descomposición del valor de la firma]
\[
V = E + D = \underbrace{\frac{FCF}{k_u}}_{\text{operaciones corrientes}}
  + \underbrace{\frac{D\,k_d\,t}{k_u}}_{\text{ahorro fiscal de la deuda}}
  + \underbrace{\frac{I\,(TIR - WACC)\,T}{WACC\,(1+WACC)}}_{\text{plan futuro de inversiones}}
\]
El tercer término solo es positivo si las inversiones futuras $I$, realizadas durante $T$ años, rinden más que el WACC.
\end{formulaBox}

%% =====================================================================
%%  UNIDAD 5 (parte B): riesgo y rendimiento (Dumrauf, cap. 7; Ross)
%% =====================================================================
\section{Riesgo y rendimiento}\label{sec:u5-riesgo}

Para valuar un bono o una acción se descuentan sus flujos a la tasa de las inversiones de riesgo similar \vease{sec:u4-valuacion}, pero falta precisar cómo se mide el riesgo y cuánto rendimiento adicional exige. Debería existir un «premio» que compense al inversor por el mayor riesgo que acepta. La sección muestra primero la evidencia de los mercados, repasa la estadística necesaria y construye la teoría del portafolio, que distingue el riesgo único del riesgo de mercado y mide este último con el coeficiente beta; es la base del CAPM \vease{sec:u5-capm}.

\subsection{¿La Bolsa recompensa a los inversores?}

El rendimiento de un activo financiero combina las ganancias (o pérdidas) de capital con los dividendos o intereses, y se expresa en porcentaje para que no dependa del monto invertido y pueda compararse. \emph{Ex post} es una magnitud conocida con certeza; \emph{ex ante} es una variable aleatoria, que toma distintos valores con ciertas probabilidades.

\paragraph{Estados Unidos, 1925-2009.} El estudio de Ibbotson y Sinquefield (1982), actualizado periódicamente, sigue la inversión de un dólar desde fines de 1925 en letras del Tesoro a tres meses (\emph{T-bills}), bonos del gobierno a veinte años (\emph{T-bonds}), acciones de compañías grandes y de compañías pequeñas (el 20\,\% más pequeño de la Bolsa de Nueva York), comparada con la inflación, reinvirtiendo dividendos e intereses y sin ajustar por inflación (rendimientos nominales antes de impuestos). Letras y bonos apenas superaron la inflación; los bonos rindieron algo más que las letras pero con precios más variables, porque al ser \emph{bullet} son más sensibles a las tasas de interés. Las acciones rindieron más, con más volatilidad: un dólar en acciones pequeñas se convirtió en U\$S 6.545. Los rendimientos se expresan como media geométrica (tasa equivalente compuesta), que siempre es menor que la media aritmética salvo que todos los rendimientos sean iguales: $r = (6.545/1)^{1/84} - 1 = 11{,}0\,\%$ anual. Los bonos del gobierno de Estados Unidos se consideran libres de riesgo, porque el gobierno siempre puede subir impuestos para pagar; algunos analistas reservan esa condición para las letras, de vencimiento más corto.

\begin{figure}[H]
  \centering
  \includegraphics[width=0.5\linewidth,height=0.42\textheight,keepaspectratio]{Finanzas20corporativas203ed20Dumrauf-86}
  \caption{Evolución de un dólar invertido a fines de 1925 en letras, bonos, acciones grandes y pequeñas de EE.\,UU., e inflación (hasta 2009).}
\end{figure}

\paragraph{Latinoamérica, 1989-2009.} Los mercados emergentes suelen ser más volátiles. Medidos en dólares (series de Economatica, sin dividendos), casi todos los mercados de la región rindieron más del 10\,\% anual, con dos casos extremos: Lima, 30,8\,\%, y Caracas, 3,94\,\%, un mercado poco profundo. Para la Argentina el rendimiento por dividendos rondaría el 2\,\% anual, lo que llevaría el Merval a alrededor del 12\,\%.

\begin{figure}[H]
  \centering
  \includegraphics[width=0.5\linewidth,height=0.42\textheight,keepaspectratio]{Finanzas20corporativas203ed20Dumrauf-87}
\end{figure}
\begin{figure}[H]
  \centering
  \includegraphics[width=0.5\linewidth,height=0.42\textheight,keepaspectratio]{Finanzas20corporativas203ed20Dumrauf-88}
\end{figure}
\begin{figure}[H]
  \centering
  \includegraphics[width=0.5\linewidth,height=0.42\textheight,keepaspectratio]{Finanzas20corporativas203ed20Dumrauf-89}
  \caption{Rendimientos en dólares de las bolsas latinoamericanas y de EE.\,UU., 1989-2009 (figuras 7.2 a 7.4; la 7.3 muestra Colombia, México, EE.\,UU. y Venezuela y la 7.4, Perú).}
\end{figure}

La Bolsa es un \textbf{casino al revés}: quien va todos los días al casino puede ganar alguna vez, pero a la larga pierde, porque las probabilidades favorecen a la casa. En la Bolsa ocurre lo contrario: invirtiendo en una cartera diversificada —un índice, no una sola acción, que puede quebrar— y reinvirtiendo las ganancias, la evidencia muestra que a largo plazo se recompensa al inversor paciente, siempre que no necesite liquidar en el corto o mediano plazo.

\begin{definicion}{Prima por el riesgo de mercado (\emph{market risk premium})}
Es la diferencia entre el rendimiento esperado de un índice de acciones y el rendimiento libre de riesgo. Es una expectativa basada en la evidencia: el mercado no promete esa recompensa, pero quien invierte con paciencia en una cartera bien diversificada puede esperar obtener, en un período largo, ese diferencial sobre la tasa libre de riesgo. En la práctica se mide contra los bonos del gobierno.
\end{definicion}

\begin{center}
\small
\begin{tabular}{@{}p{0.52\linewidth}>{\centering\arraybackslash}p{0.18\linewidth}>{\centering\arraybackslash}p{0.2\linewidth}@{}}
\toprule
Activo (EE.\,UU. 1925-2009; Latinoamérica 1989-2009) & Rendimiento geométrico & Diferencia sobre los \emph{T-bonds}\\
\midrule
Letras del Tesoro de EE.\,UU. & 3,7\,\% & $-1{,}7$\,\%\\
Bonos del Tesoro de EE.\,UU. & 5,4\,\% & 0\,\%\\
Acciones de compañías grandes & 9,8\,\% & 4,4\,\%\\
Acciones de compañías pequeñas & 11,0\,\% & 5,6\,\%\\
Acciones latinoamericanas (promedio de Brasil, Argentina, Colombia, Chile y México, sin dividendos) & 12,5\,\% & 7,1\,\%\\
\bottomrule
\end{tabular}
\end{center}

La evidencia es contundente: a mayor riesgo, mayor rentabilidad, y la prima crece en los índices de compañías más pequeñas, tal como había postulado el CAPM.

%% ---------------------------------------------------------------------
\subsection{Repaso de las categorías de estadística utilizadas en finanzas}

\paragraph{La distribución normal.} Describe muchos fenómenos y se usa para representar cómo se distribuyen los rendimientos alrededor de su media $\bar{x}$: con forma de campana, alrededor del 68\,\% de los resultados cae a un desvío estándar $\sigma$ de la media, el 95\,\% a dos y el 99,7\,\% a tres. Es una aproximación razonable para los rendimientos de las acciones —los precios siguen más bien una distribución lognormal—: si el mercado estadounidense rindió en promedio 11,76\,\% con un desvío de 20,4\,\%, en el 95\,\% de los años su rendimiento debería ubicarse entre $-29{,}04$\,\% y 52,56\,\%. Pero la normal subestima los eventos de las colas: las grandes ganancias y pérdidas son más probables de lo que predice, el tema de \emph{El cisne negro} de Taleb (2007).

\begin{figure}[H]
  \centering
  \includegraphics[width=0.5\linewidth,height=0.42\textheight,keepaspectratio]{Finanzas20corporativas203ed20Dumrauf-90}
\end{figure}
\begin{figure}[H]
  \centering
  \includegraphics[width=0.5\linewidth,height=0.42\textheight,keepaspectratio]{Finanzas20corporativas203ed20Dumrauf-91}
  \caption{La distribución normal y la distribución de frecuencias de los rendimientos anuales del mercado accionario de EE.\,UU. (figuras 7.5 y 7.6).}
\end{figure}

\paragraph{Media y media ponderada.} La media aritmética no considera la importancia de cada valor; la media ponderada sí, porque pondera cada rendimiento por su probabilidad (o, en un portafolio, por la proporción invertida). Con rendimientos de 10\,\% y 20\,\%, la media es siempre 15\,\%, pero con probabilidades de 40\,\% y 60\,\% la media ponderada es $0{,}40\times 10\,\% + 0{,}60\times 20\,\% = 16\,\%$. Para las acciones suele tomarse la media histórica como referencia, porque la media muestral es un estimador insesgado de la poblacional; suponer que el futuro será igual al pasado es heroico, y por eso estimar rendimientos tiene algo de arte.

\paragraph{Variabilidad de los rendimientos.} Los rendimientos de las acciones fluctúan mucho más que los de los bonos: en EE.\,UU. el peor año fue 1931 ($-43{,}8$\,\%) y el mejor, 1954 (52,6\,\%). En Latinoamérica la volatilidad es aún mayor: casi todos los mercados tuvieron años con subas de más del 100\,\% y caídas de más del 70\,\%, y en 2008 la crisis internacional de las hipotecas \emph{subprime} los golpeó a todos.

\begin{figure}[H]
  \centering
  \includegraphics[width=0.5\linewidth,height=0.42\textheight,keepaspectratio]{Finanzas20corporativas203ed20Dumrauf-92}
\end{figure}
\begin{figure}[H]
  \centering
  \includegraphics[width=0.5\linewidth,height=0.42\textheight,keepaspectratio]{Finanzas20corporativas203ed20Dumrauf-93}
\end{figure}
\begin{figure}[H]
  \centering
  \includegraphics[width=0.5\linewidth,height=0.42\textheight,keepaspectratio]{Finanzas20corporativas203ed20Dumrauf-94}
\end{figure}
\begin{figure}[H]
  \centering
  \includegraphics[width=0.5\linewidth,height=0.42\textheight,keepaspectratio]{Finanzas20corporativas203ed20Dumrauf-95}
\end{figure}
\begin{figure}[H]
  \centering
  \includegraphics[width=0.5\linewidth,height=0.42\textheight,keepaspectratio]{Finanzas20corporativas203ed20Dumrauf-96}
\end{figure}
\begin{figure}[H]
  \centering
  \includegraphics[width=0.5\linewidth,height=0.42\textheight,keepaspectratio]{Finanzas20corporativas203ed20Dumrauf-97}
\end{figure}
\begin{figure}[H]
  \centering
  \includegraphics[width=0.5\linewidth,height=0.42\textheight,keepaspectratio]{Finanzas20corporativas203ed20Dumrauf-98}
\end{figure}
\begin{figure}[H]
  \centering
  \includegraphics[width=0.5\linewidth,height=0.42\textheight,keepaspectratio]{Finanzas20corporativas203ed20Dumrauf-99}
\end{figure}
\begin{figure}[H]
  \centering
  \includegraphics[width=0.5\linewidth,height=0.42\textheight,keepaspectratio]{Finanzas20corporativas203ed20Dumrauf-100}
  \caption{Rendimientos anuales de acciones y bonos del Tesoro en EE.\,UU. y de los principales mercados latinoamericanos, con su mayor suba, su mayor baja y el año 2008 (figuras 7.7 a 7.14).}
\end{figure}

\begin{definicion}{Varianza y desvío estándar}
Miden la dispersión de los resultados alrededor del promedio: cuanto mayores, más dispersos los rendimientos. La \textbf{varianza} es el promedio de los desvíos al cuadrado (se elevan al cuadrado porque los desvíos positivos y negativos se cancelarían); como se expresa en «unidades cuadradas», se interpreta mejor su raíz, el \textbf{desvío estándar} o desviación típica, que está en la misma unidad que los rendimientos y es la medida usual de la \textbf{volatilidad}.
\end{definicion}

\begin{formulaBox}[Media, varianza y desvío estándar]
\[
\bar{x} = \frac{1}{n}\sum_{t=1}^{n} r_t
\qquad\qquad
s^{2} = \frac{\sum (r_t-\bar{x})^{2}}{n-1}
\qquad\qquad
s = \sqrt{s^{2}}
\]
Con $n-1$ se obtiene un estimador insesgado de la varianza poblacional; con muestras grandes da prácticamente lo mismo que dividir por $n$.
\end{formulaBox}

\begin{ejemplo}[Varianza histórica de una acción]
Una acción rindió 10\,\%, 15\,\% y $-4$\,\% en tres años: $\bar{x} = 7\,\%$. Los desvíos son 0,03, 0,08 y $-0{,}11$ (suman cero) y sus cuadrados 0,0009, 0,0064 y 0,0121, que suman 0,0194. Entonces $s^{2} = 0{,}0194/2 = 0{,}0097$ y $s = \sqrt{0{,}0097} = 9{,}8\,\%$. Si los rendimientos fueran normales, en el 68\,\% de los casos estarían entre $-2{,}8$\,\% y 16,8\,\%. (Con rendimientos en puntos porcentuales la varianza daría 97, diez mil veces más, y el desvío 9,8: la unidad no cambia la interpretación.)
\end{ejemplo}

\begin{figure}[H]
  \centering
  \includegraphics[width=0.5\linewidth,height=0.42\textheight,keepaspectratio]{Finanzas20corporativas203ed20Dumrauf-102}
  \caption{Cálculo del desvío estándar a partir de la varianza.}
\end{figure}

La práctica de los administradores de carteras, que trabajan con series largas (por ejemplo, de dos o tres años de datos diarios), es dividir por $n$: para Tenaris entre 2006 y 2009 la varianza da 0,001224 con $n-1$ y 0,001222 con $n$, y los desvíos coinciden hasta el cuarto decimal (0,03498 y 0,03496). Para dividir por $n$ hay además un argumento formal: corresponde cuando la media se conoce \emph{a priori} en lugar de estimarse con la misma muestra (Flannery y Teukolsky, citados por Benninga).

\begin{formulaBox}[Volatilidad anual]
\[
\sigma^{2}_{\text{anual}} = \sigma^{2}_{\text{diaria}}\times N
\qquad\qquad
\sigma_{\text{anual}} = \sigma_{\text{diaria}}\times\sqrt{N}
\]
$N$ es la cantidad de ruedas hábiles del año.
\end{formulaBox}

\begin{figure}[H]
  \centering
  \includegraphics[width=0.5\linewidth,height=0.42\textheight,keepaspectratio]{Finanzas20corporativas203ed20Dumrauf-103}
\end{figure}
\begin{figure}[H]
  \centering
  \includegraphics[width=0.5\linewidth,height=0.42\textheight,keepaspectratio]{Finanzas20corporativas203ed20Dumrauf-104}
  \caption{Varianza anual a partir de la varianza diaria y volatilidad anual.}
\end{figure}

\begin{definicion}{Covarianza y coeficiente de correlación}
La \textbf{covarianza} mide cómo dos variables aleatorias tienden a moverse juntas: es positiva si, cuando una supera su media, la otra también tiende a hacerlo; negativa si se mueven en sentidos opuestos, y nula si no hay relación regular. Como la varianza, es una medida «cuadrada» difícil de interpretar. El \textbf{coeficiente de correlación} $\rho$ tiene su mismo signo, pero está acotado entre $-1$ y 1 y mide la fuerza de la asociación: 1 es correlación positiva perfecta, entre 0 y 1 positiva imperfecta, $-1$ negativa perfecta, entre 0 y $-1$ negativa imperfecta y cerca de 0, ausencia de correlación. Entre acciones reales lo normal son valores entre 0 y 1; una correlación negativa fuerte es muy rara.
\end{definicion}

\begin{formulaBox}[Covarianza y correlación]
\[
\sigma_{xy} = \frac{1}{n}\sum_{t=1}^{n}(r_{x,t}-\bar{x})(r_{y,t}-\bar{y})
\qquad\qquad
\rho_{xy} = \frac{\sigma_{xy}}{\sigma_x\,\sigma_y}
\qquad\qquad
\sigma_{xy} = \rho_{xy}\,\sigma_x\,\sigma_y
\]
\end{formulaBox}

\begin{ejemplo}[Covarianza de dos acciones]
X rindió 0,10, 0,15 y $-0{,}04$ (media 0,07, desvío 0,080) e Y, $-0{,}03$, $-0{,}04$ y 0,12 (media 0,017, desvío 0,073). Los productos de los desvíos suman $-0{,}0173$, de modo que $\sigma_{xy} = -0{,}0173/3 = -0{,}00576$ (Excel, que divide por $n$, da lo mismo) y $\rho_{xy} = -0{,}00576/(0{,}080\times 0{,}073) \approx -0{,}98$, una correlación negativa casi perfecta.
\end{ejemplo}

\begin{figure}[H]
  \centering
  \includegraphics[width=0.5\linewidth,height=0.42\textheight,keepaspectratio]{Finanzas20corporativas203ed20Dumrauf-106}
\end{figure}
\begin{figure}[H]
  \centering
  \includegraphics[width=0.5\linewidth,height=0.42\textheight,keepaspectratio]{Finanzas20corporativas203ed20Dumrauf-108}
\end{figure}
\begin{figure}[H]
  \centering
  \includegraphics[width=0.5\linewidth,height=0.42\textheight,keepaspectratio]{Finanzas20corporativas203ed20Dumrauf-109}
\end{figure}
\begin{figure}[H]
  \centering
  \includegraphics[width=0.5\linewidth,height=0.42\textheight,keepaspectratio]{Finanzas20corporativas203ed20Dumrauf-110}
  \caption{Interpretación gráfica del coeficiente de correlación (figuras 7.18 a 7.22).}
\end{figure}

%% ---------------------------------------------------------------------
\subsection{El portafolio o cartera de inversiones}

Los inversores «no ponen todos los huevos en la misma canasta»: mantienen acciones de varias empresas, bonos, inmuebles o monedas, y una empresa que invierte en varios negocios también tiene una cartera. La teoría del portafolio de Harry Markowitz (1952), premio Nobel en 1990 junto con William Sharpe y Merton Miller, enseña que el riesgo de un activo no debe juzgarse por la variabilidad de su propio rendimiento, sino por su \textbf{contribución marginal al riesgo de la cartera}. El rendimiento esperado de la cartera es el promedio ponderado de los rendimientos de sus activos, pero su riesgo es menor que el promedio ponderado de sus desvíos.

\begin{figure}[H]
  \centering
  \includegraphics[width=0.5\linewidth,height=0.42\textheight,keepaspectratio]{Finanzas20corporativas203ed20Dumrauf-111}
  \caption{La cartera de inversiones: diversificar entre varios activos.}
\end{figure}

\paragraph{Rendimiento esperado.} El rendimiento esperado es una expectativa matemática: por la ley de los grandes números, en muchas inversiones los buenos y los malos resultados se compensan y el promedio se acerca a la media del conjunto. En la práctica los administradores estiman una media histórica de dos o tres años y la corrigen, por ejemplo, por rebotes de corto plazo después de una caída.

\begin{formulaBox}[Rendimiento y varianza de un portafolio de dos activos]
\[
\begin{gathered}
E(r_p) = w_1\,E(r_1) + w_2\,E(r_2) + \dots + w_n\,E(r_n)
\\[4pt]
\sigma_p^{2} = w_A^{2}\,\sigma_A^{2} + w_B^{2}\,\sigma_B^{2} + 2\,w_A\,w_B\,\rho_{AB}\,\sigma_A\,\sigma_B
\qquad\qquad
\sigma_p = \sqrt{\sigma_p^{2}}
\end{gathered}
\]
Con dos activos hay tantas varianzas como covarianzas: la de A, la de B, la de A con B y la de B con A.
\end{formulaBox}

\sinejemplos{Las figuras siguientes aplican estas fórmulas a dos activos: A, con un rendimiento esperado del 21\,\% y un desvío del 40\,\%, y B, con un rendimiento esperado del 15\,\% y un desvío del 20\,\%.}

\begin{ejemplo}[Portafolio de dos activos]
Se invierte el 20\,\% en A ($E(r) = 21\,\%$, $\sigma = 40\,\%$) y el 80\,\% en B ($E(r) = 15\,\%$, $\sigma = 20\,\%$), con $\rho_{AB} = 0{,}50$. El rendimiento esperado es $0{,}20\times 21\,\% + 0{,}80\times 15\,\% = 16{,}2\,\%$. El promedio ponderado de los desvíos daría 24\,\%, pero solo sería correcto con correlación perfecta:
\[
\begin{aligned}
\sigma_p^{2} &= 0{,}20^{2}\times 40^{2} + 0{,}80^{2}\times 20^{2} + 2\times 0{,}20\times 0{,}80\times 0{,}50\times 40\times 20\\
&= 64 + 256 + 128 = 448
\qquad\Longrightarrow\qquad \sigma_p = \sqrt{448} = 21{,}16\,\%
\end{aligned}
\]
(Con desvíos en tanto por uno la varianza es 0,0448, y su raíz, 0,2116.) Para las mismas proporciones, el desvío baja al reducirse la correlación:
\begin{center}
\begin{tabular}{@{}lccccc@{}}
\toprule
$\rho_{AB}$ & $-1$ & $-0{,}5$ & 0 & 0,5 & 1\\
\midrule
$\sigma_p$ & 8,0\,\% & 13,9\,\% & 17,9\,\% & 21,2\,\% & 24,0\,\%\\
\bottomrule
\end{tabular}
\end{center}
\end{ejemplo}

\begin{figure}[H]
  \centering
  \includegraphics[width=0.5\linewidth,height=0.42\textheight,keepaspectratio]{Finanzas20corporativas203ed20Dumrauf-115}
\end{figure}
\begin{figure}[H]
  \centering
  \includegraphics[width=0.5\linewidth,height=0.42\textheight,keepaspectratio]{Finanzas20corporativas203ed20Dumrauf-116}
  \caption{Varianza y desvío estándar del portafolio de los activos A y B.}
\end{figure}

El riesgo del portafolio depende de la proporción invertida en cada activo, del desvío de cada uno y de la covarianza entre ellos. Las varianzas son siempre positivas, pero las covarianzas pueden ser negativas, y eso es lo que permite diversificar. Mientras la correlación sea menor que 1, el desvío del portafolio es menor que el promedio ponderado de los desvíos individuales; solo con correlación 1 la relación entre riesgo y rendimiento es una recta y la diversificación no ayuda. Con correlación 0,5 la curva de combinaciones se desplaza hacia el origen, y con $-1$ el riesgo baja al principio hasta casi el 4\,\% al aumentar la participación de A antes de volver a crecer. El conjunto de oportunidades (conjunto factible) se curva hacia atrás siempre que $\rho \le 0$ y puede hacerlo o no si $\rho > 0$; el efecto diversificación crece a medida que baja la correlación.

\begin{figure}[H]
  \centering
  \includegraphics[width=0.5\linewidth,height=0.42\textheight,keepaspectratio]{Finanzas20corporativas203ed20Dumrauf-117}
  \caption{Combinaciones de riesgo y rendimiento de A y B para distintos coeficientes de correlación (figura 7.24).}
\end{figure}

\paragraph{Correlación negativa perfecta.} En teoría —solo en teoría— puede eliminarse el riesgo combinando dos activos con $\rho = -1$ en las proporciones exactas. Correlación negativa perfecta no significa igual variación proporcional: si B varía siempre el doble que A en sentido contrario, la cartera mitad y mitad sigue teniendo riesgo; con C, que varía lo mismo que A en sentido opuesto, la cartera por mitades tiene desvío nulo.

\begin{center}
\small
\begin{tabular}{@{}lccccc@{}}
\toprule
 & A & B & C & 50\,\% A $+$ 50\,\% B & 50\,\% A $+$ 50\,\% C\\
\midrule
Período 1 & $-10$\,\% & 20\,\% & 10\,\% & 5\,\% & 0\,\%\\
Período 2 & 5\,\% & $-10$\,\% & $-5$\,\% & $-2{,}5$\,\% & 0\,\%\\
Período 3 & $-10$\,\% & 20\,\% & 10\,\% & 5\,\% & 0\,\%\\
Período 4 & $-2$\,\% & 4\,\% & 2\,\% & 1\,\% & 0\,\%\\
\midrule
Desvío estándar & 7,23\,\% & 14,46\,\% & 7,23\,\% & 3,61\,\% & 0\,\%\\
\bottomrule
\end{tabular}
\end{center}

\begin{formulaBox}[Portafolio de riesgo nulo con correlación $-1$]
\[
\sigma_p = w\,\sigma_A - (1-w)\,\sigma_B = 0
\qquad\Longrightarrow\qquad
w = \frac{\sigma_B}{\sigma_A + \sigma_B} = \frac{20}{40+20} = 0{,}333
\]
Con $\rho = -1$ la varianza es un cuadrado perfecto, $\sigma_p^{2} = [w\,\sigma_A - (1-w)\,\sigma_B]^{2}$; para los activos A y B del portafolio de dos activos (desvíos del 40 y del 20\,\%) se invierte un tercio en A y dos tercios en B.
\end{formulaBox}

\begin{figure}[H]
  \centering
  \includegraphics[width=0.5\linewidth,height=0.42\textheight,keepaspectratio]{Finanzas20corporativas203ed20Dumrauf-121}
  \caption{Desvío del portafolio con correlación negativa perfecta.}
\end{figure}

\paragraph{Portafolios eficientes.} Con muchos títulos, cada combinación de rendimiento esperado y riesgo es un punto de una figura parecida a un paraguas sin mango. Como el inversor quiere más rendimiento y menos riesgo, solo le interesan las carteras de su borde superior.

\begin{definicion}{Frontera eficiente}
Es el conjunto de \textbf{portafolios eficientes}: los que ofrecen el mayor rendimiento esperado para un nivel de riesgo dado o, lo que es equivalente, el menor riesgo para un rendimiento dado. Es la parte del conjunto factible situada por encima de la cartera de mínima varianza. Qué cartera de la frontera elegir depende de la preferencia por el riesgo, pero nunca conviene una cartera situada por debajo de ella.
\end{definicion}

\paragraph{El teorema de la separación y la línea del mercado de capitales.} Tobin (1958) amplió el modelo combinando acciones con un activo libre de riesgo: se puede invertir una parte en acciones y prestar el resto a la tasa libre de riesgo, o invertir en acciones más de lo que se tiene, endeudándose a esa tasa. De todas las rectas que parten de $r_f$, la que toca la frontera en la cartera M domina a las demás: para cualquier riesgo ofrece más rendimiento.

\begin{figure}[H]
  \centering
  \includegraphics[width=0.5\linewidth,height=0.42\textheight,keepaspectratio]{Finanzas20corporativas203ed20Dumrauf-123}
\end{figure}
\begin{figure}[H]
  \centering
  \includegraphics[width=0.5\linewidth,height=0.42\textheight,keepaspectratio]{Finanzas20corporativas203ed20Dumrauf-124}
  \caption{Frontera eficiente de carteras riesgosas y combinaciones con el activo libre de riesgo: la línea que une $r_f$ con M (figuras 7.25 y 7.26).}
\end{figure}

\begin{ejemplo}[Prestar y endeudarse a la tasa libre de riesgo]
Con $r_f = 5\,\%$ y una cartera M con $E(r_M) = 18\,\%$ y $\sigma_M = 20\,\%$: invirtiendo la mitad en cada una, $E(r) = 0{,}50\times 5\,\% + 0{,}50\times 18\,\% = 11{,}5\,\%$ y $\sigma = 0{,}50\times 20\,\% = 10\,\%$ (el activo libre de riesgo no tiene desvío ni correlación con M). Si en cambio se invierte en M el doble del patrimonio, tomando prestado otro tanto a la tasa libre de riesgo, $E(r) = 2\times 18\,\% - 1\times 5\,\% = 31\,\%$ y $\sigma = 2\times 20\,\% = 40\,\%$.
\end{ejemplo}

\begin{figure}[H]
  \centering
  \includegraphics[width=0.5\linewidth,height=0.42\textheight,keepaspectratio]{Finanzas20corporativas203ed20Dumrauf-126}
  \caption{Desvío estándar de una cartera que combina el activo libre de riesgo con la cartera M.}
\end{figure}

\begin{definicion}{Línea del mercado de capitales (CML) y teorema de la separación}
La \textbf{CML} es la recta que une $r_f$ con la cartera M y es la nueva frontera eficiente cuando se puede prestar y tomar prestado a la tasa libre de riesgo. El \textbf{teorema de la separación} dice que la inversión se decide en dos etapas: primero se elige la cartera eficiente de activos riesgosos M —la misma para todos— y después se la combina con préstamos o endeudamiento según la preferencia de cada uno por el riesgo. Como todos quieren la misma combinación óptima, en equilibrio M es la \textbf{cartera de mercado}: incluye todas las acciones en proporción a su capitalización bursátil.
\end{definicion}

\begin{formulaBox}[Línea del mercado de capitales y precio de mercado del riesgo]
\[
E(r_p) = r_f + \frac{E(r_M) - r_f}{\sigma_M}\,\sigma_p
\qquad\qquad
\text{Precio de mercado del riesgo} = \frac{r_M - r_f}{\sigma_M} = \frac{E(r_p) - r_f}{\sigma_p}
\]
La prima de riesgo de una cartera eficiente es proporcional a su desvío estándar: en equilibrio, el cociente entre prima y riesgo es constante.
\end{formulaBox}

\begin{figure}[H]
  \centering
  \includegraphics[width=0.15\linewidth,height=0.42\textheight,keepaspectratio]{Finanzas20corporativas203ed20Dumrauf-128}
  \caption{Relación constante entre la prima de riesgo y el desvío estándar en la línea del mercado de capitales.}
\end{figure}

Las carteras sobre la CML están perfectamente correlacionadas entre sí y con la cartera de mercado, y el arbitraje empuja hacia la línea a las carteras del interior del paraguas. La cartera de mercado es, en rigor, una ficción: lo que hay es una sucesión de carteras de equilibrio a lo largo de la CML. Cuanto más explique el riesgo el rendimiento de las distintas carteras, más cerca del equilibrio estará el mercado. La idea supone \textbf{expectativas homogéneas}: todos comparten las mismas creencias sobre rendimientos, varianzas y covarianzas, aunque cada uno tenga distinta aversión al riesgo.

\begin{ejemplo}[Precio de mercado del riesgo]
En un período largo el S\&P 500 rindió cerca del 11\,\%, los bonos del Tesoro el 5\,\% y $\sigma_M = 20\,\%$: $(0{,}11-0{,}05)/0{,}20 = 0{,}30$, es decir, 0,30 puntos de rendimiento adicional por cada punto porcentual de riesgo asumido. Con la volatilidad de 40 ruedas al 24 de noviembre de 2009 y rendimientos geométricos sin dividendos, el cociente fue de 0,41 en Brasil (Ibovespa), 0,47 en Colombia (IGBC), 1,22 en Perú (IGBVL), 0,81 en Chile (IGPA), $-0{,}08$ en Venezuela, 0,16 en México y en la Argentina (Merval) y 0,32 en EE.\,UU.: salvo en Venezuela, todos los mercados recompensaron el riesgo.
\end{ejemplo}

\paragraph{Límites a los beneficios de la diversificación.} El riesgo propio de cada activo, que afecta a uno solo o a un grupo pequeño (una huelga en una industria, por ejemplo), es el \textbf{riesgo no sistemático}, \emph{específico} o \emph{único}, y se reduce manteniendo una cartera diversificada. El que afecta a un gran grupo de activos —cambios en la tasa de interés, en la actividad económica, en el tipo de cambio o en la inflación, o hechos imponderables como un golpe de Estado— es el \textbf{riesgo sistemático} o \emph{de mercado}, y no se elimina diversificando; suele manifestarse como un «efecto sorpresa» cuando la información no estaba anticipada. El riesgo de la cartera cae a medida que se agregan títulos, pero con rendimientos marginales decrecientes: con unas diez acciones ya se logró la mayor parte del efecto. Con acciones argentinas, una diversificación ingenua en partes iguales reduce el riesgo del 42\,\% a menos del 15\,\% con 23 acciones, incluso por debajo del Merval, que concentra su ponderación en pocas acciones.

\begin{figure}[H]
  \centering
  \includegraphics[width=0.65\linewidth,height=0.42\textheight,keepaspectratio]{Finanzas20corporativas203ed20Dumrauf-132}
  \caption{El riesgo de la cartera disminuye a medida que se agregan acciones, hasta el límite del riesgo de mercado (figura 7.29).}
\end{figure}

\paragraph{El riesgo sistemático está representado por la covarianza media.} Al agregar acciones, cada varianza pesa menos (se multiplica por la proporción al cuadrado) y se multiplican las covarianzas, porque cada título covaría con todos los demás. Con $N$ títulos el modelo de Markowitz requiere $N$ rendimientos esperados, $N$ varianzas y $(N^{2}-N)/2$ covarianzas distintas —las que quedan por encima de la diagonal de la matriz de covarianzas, que es simétrica—, aunque en la varianza de la cartera figuran $N^{2}-N$ términos de covarianza. Con 10 acciones hay 45 covarianzas distintas (90 términos), 10 rendimientos y 10 varianzas: 65 estimaciones; una acción más agrega 12 (una media, una varianza y 10 covarianzas).

\begin{formulaBox}[Varianza de un portafolio de $N$ activos]
\[
\begin{gathered}
\sigma_p^{2} = \sum_{i=1}^{N}\sum_{j=1}^{N} w_i\,w_j\,\sigma_{ij}
\qquad (\sigma_{ii} = \sigma_i^{2})
\\[4pt]
\text{con } w_i = \tfrac{1}{N}:\qquad
\sigma_p^{2} = \frac{1}{N}\,\overline{\text{varianza}} + \left(1-\frac{1}{N}\right)\overline{\text{covarianza}}
\end{gathered}
\]
Cuando $N$ crece, la varianza de la cartera tiende a la covarianza media.
\end{formulaBox}

\begin{figure}[H]
  \centering
  \includegraphics[width=0.5\linewidth,height=0.42\textheight,keepaspectratio]{Finanzas20corporativas203ed20Dumrauf-134}
\end{figure}
\begin{figure}[H]
  \centering
  \includegraphics[width=0.5\linewidth,height=0.42\textheight,keepaspectratio]{Finanzas20corporativas203ed20Dumrauf-135}
\end{figure}
\begin{figure}[H]
  \centering
  \includegraphics[width=0.5\linewidth,height=0.42\textheight,keepaspectratio]{Finanzas20corporativas203ed20Dumrauf-137}
\end{figure}
\begin{figure}[H]
  \centering
  \includegraphics[width=0.5\linewidth,height=0.42\textheight,keepaspectratio]{Finanzas20corporativas203ed20Dumrauf-138}
\end{figure}
\begin{figure}[H]
  \centering
  \includegraphics[width=0.5\linewidth,height=0.42\textheight,keepaspectratio]{Finanzas20corporativas203ed20Dumrauf-140}
  \caption{Matriz de varianzas y covarianzas de $N$ títulos y varianza del portafolio.}
\end{figure}

\begin{center}
\small
\begin{tabular}{@{}lccccc@{}}
\toprule
Activos en el portafolio & 2 & 5 & 10 & 20 & 100\\
\midrule
Términos de covarianza ($N^{2}-N$) & 2 & 20 & 90 & 380 & 9.900\\
\bottomrule
\end{tabular}
\end{center}

Si la covarianza media fuese cero, acumulando títulos se eliminaría todo el riesgo; pero las acciones se mueven juntas, ligadas por una red de covarianzas positivas que pone un límite a la diversificación. Por eso la varianza de una cartera con muchos títulos depende más de las covarianzas que de las varianzas individuales. Estas estimaciones se hacen con datos históricos, suponiendo que el futuro prolonga el pasado; en una planilla, la matriz de covarianzas se obtiene con la herramienta «Análisis de datos».

%% ---------------------------------------------------------------------
\subsection{El riesgo de mercado de la acción y su contribución al riesgo del portafolio}

Si el riesgo único puede eliminarse diversificando, el mercado no lo recompensa: el único premio que el inversor puede esperar es por el riesgo sistemático. Quien tiene un solo título debe medir su rendimiento por el rendimiento esperado y su riesgo por la varianza o el desvío; quien tiene una cartera diversificada se preocupa, en cambio, por la contribución de cada título al rendimiento y al riesgo de la cartera, y esa contribución la mide el beta.

\begin{definicion}{Coeficiente beta}
La acción $j$ aporta al riesgo del portafolio $w_j\,\sigma_{jp}$, y su contribución proporcional es $w_j\,\sigma_{jp}/\sigma_p^{2}$. El cociente $\sigma_{jp}/\sigma_p^{2}$ mide la sensibilidad de la acción a los cambios del rendimiento del portafolio: si es mayor que 1, aumentar su participación aumenta el riesgo. Cuando el portafolio es la cartera de mercado, ese cociente es el \textbf{beta}: la sensibilidad del rendimiento de la acción a los cambios del rendimiento del mercado y, por lo tanto, la medida relevante de su riesgo de mercado. Un beta de 1 indica que la acción varía como el mercado; de 0,5, la mitad; mayor que 1, más que el mercado. Prácticamente ninguna acción tiene beta negativo.
\end{definicion}

\begin{formulaBox}[Beta de una acción]
\[
\beta_j = \frac{\operatorname{Cov}(r_j, r_M)}{\operatorname{Var}(r_M)} = \frac{\sigma_{jM}}{\sigma_M^{2}} = \frac{\rho_{jM}\,\sigma_j\,\sigma_M}{\sigma_M^{2}} = \rho_{jM}\,\frac{\sigma_j}{\sigma_M}
\]
El beta es la fracción del desvío del activo que contribuye al riesgo de la cartera de mercado; si $w_j$ es su peso en esa cartera, $w_j\,\beta_j$ es su contribución relativa.
\end{formulaBox}

\begin{figure}[H]
  \centering
  \includegraphics[width=0.17\linewidth,height=0.42\textheight,keepaspectratio]{Finanzas20corporativas203ed20Dumrauf-144}
  \caption{Contribución proporcional de una acción al riesgo del portafolio.}
\end{figure}

A veces aparecen betas negativos durante algún período (Celulosa, en la Argentina de los noventa): incluir esa acción reduciría mucho el riesgo, pero antes conviene verificar su liquidez, porque la falta de negociación distorsiona el riesgo de mercado medido (Ibbotson propone un método para corregirlo). El desarrollo matemático que lleva de Markowitz (1952) al CAPM de Sharpe (1964) se omite aquí; el modelo se presenta en §\,\ref{sec:u5-capm}.

\paragraph{Estimación del beta por regresión.} Se grafican los rendimientos de la acción contra los del índice y se ajusta una línea de tendencia lineal, pidiendo en la planilla la ecuación y el $R^{2}$. Para Tenaris contra el Merval resulta $y = 1{,}1913\,x + 0{,}0027$: la ordenada al origen es casi nula —si el Merval no cambia, Tenaris tampoco— y la pendiente, 1,19, es su beta: por cada 1\,\% adicional que varía el mercado, Tenaris varía en promedio un 1,2\,\%. La recta minimiza la suma de los desvíos cuadrados de los puntos observados, y su bondad de ajuste es el $R^{2}$: 0,7244 indica que el 72,4\,\% de las variaciones de Tenaris se explica por el Merval (riesgo de mercado) y el 27,6\,\% restante es riesgo único. Las consultoras publican betas, que difieren según el período y la técnica de estimación.

\begin{formulaBox}[Línea característica]
\[
r_j = \alpha + \beta\,r_M + \varepsilon
\]
$\alpha$ y $\beta\,r_M$ son la parte del rendimiento ligada al riesgo sistemático; el error $\varepsilon$ es la parte no relacionada con el mercado: el riesgo único, que se elimina diversificando.
\end{formulaBox}

\begin{figure}[H]
  \centering
  \includegraphics[width=0.75\linewidth,height=0.42\textheight,keepaspectratio]{Finanzas20corporativas203ed20Dumrauf-148}
  \caption{Recta de regresión de los rendimientos de Tenaris contra los del Merval: beta 1,19 y $R^{2}$ de 0,7244 (figura 7.32).}
\end{figure}

\paragraph{Ideas fundamentales.} El riesgo de un portafolio bien diversificado depende del riesgo de mercado de sus títulos, no del riesgo de cada uno por separado. Los inversores prefieren rentabilidad alta y desviación baja, y las carteras que ofrecen la mayor rentabilidad para cada riesgo son las eficientes. Si pueden prestar y endeudarse a la tasa libre de riesgo, deberían combinar ese activo con una cartera de acciones; sin información confidencial, todos tendrían la misma, la cartera de mercado. Y el beta mide la contribución marginal de una acción al riesgo de esa cartera: el rendimiento esperado de la acción varía en proporción a su beta.

\subsection{Medidas de desempeño en la administración de portafolios}

Para juzgar si vale la pena pagar a un administrador profesional se compara su cartera con una no administrada, como un fondo que replica un índice; la comparación también pone a prueba la eficiencia del mercado en su forma fuerte \vease{sec:u5-capm}. El desempeño se describe con varios indicadores, no con uno solo.

\begin{definicion}{Índices de Sharpe, Treynor y Jensen}
Los índices de \textbf{Sharpe} y de \textbf{Treynor} miden la prima que paga el mercado por cada unidad de riesgo: Sharpe usa el riesgo total (el desvío de la cartera) y Treynor el riesgo sistemático (el beta); una cartera es mejor cuanto más alto es su índice. El índice de \textbf{Jensen} mide la distancia vertical entre el rendimiento de la cartera y el que le corresponde según la línea del mercado de títulos (SML): las carteras son superiores, inferiores o neutras según sea positivo, negativo o nulo.
\end{definicion}

\begin{formulaBox}[Medidas de desempeño]
\[
S_p = \frac{r_p - r_f}{\sigma_p}
\qquad\qquad
T_p = \frac{r_p - r_f}{\beta_p}
\qquad\qquad
J_p = r_p - \left[r_f + (r_M - r_f)\,\beta_p\right]
\]
\end{formulaBox}

%% =====================================================================
%%  UNIDAD 5 (parte C): gestión de riesgos financieros (apuntes FCA)
%% =====================================================================
\section{Gestión de riesgos financieros}\label{sec:u5-fca}

Los apuntes de clase llevan la teoría del portafolio \vease{sec:u5-riesgo} a la práctica de una institución que administra carteras: cómo valuarlas, cómo medir su riesgo con datos históricos en una planilla o en R, cómo optimizarlas y qué confianza merecen esas estimaciones. Las definiciones de media, varianza, covarianza, correlación, frontera eficiente, beta e índice de Sharpe son las de la sección anterior; aquí se agregan las herramientas de cálculo y sus límites.

\subsection{Introducción: la gestión del riesgo de un portafolio}

Un banco mantiene acciones (para invertir en empresas y generar ingresos), bonos (deuda para financiarse y diversificar), hipotecas y préstamos (ingresos por intereses), \emph{swaps} y opciones (derivados para cubrirse o invertir), \emph{commodities} (diversificación y protección contra la inflación) y monedas (transacciones internacionales y cobertura cambiaria). Gestionar ese portafolio exige responder tres preguntas.

\paragraph{¿Cuánto valen mis activos?} Su valor de mercado $MV$ es el valor presente de sus flujos $CF_t$ descontados a la tasa $r$; el factor $(1+r)^{t}$ lleva cada flujo futuro a valor de hoy \vease{sec:u1-va}.

\begin{figure}[H]
  \centering
  \includegraphics[width=0.5\linewidth,height=0.42\textheight,keepaspectratio]{FCA-1}
\end{figure}
\begin{figure}[H]
  \centering
  \includegraphics[width=0.5\linewidth,height=0.42\textheight,keepaspectratio]{FCA-2}
  \caption{Valor de mercado de los activos como valor presente de sus flujos: $MV = \sum_t CF_t/(1+r)^{t}$.}
\end{figure}

\paragraph{¿Cómo afectan los distintos riesgos a mi portafolio?}

\begin{definicion}{Riesgos de mercado, de crédito y operacional}
\begin{itemize}
\item \textbf{Riesgo de mercado}: fluctuaciones de precios, tipos de cambio y tasas de interés que alteran el valor de los activos; por ejemplo, una suba de tasas baja el precio de los bonos y una depreciación de la moneda de una inversión reduce su valor en la moneda base. Es lo que ocurre todos los días y lo que más se vigila.
\item \textbf{Riesgo de crédito}: posibilidad de que un emisor de deuda de la cartera incumpla sus pagos (\emph{default}) o sufra una degradación de su calificación, con pérdidas directas (pagos no cobrados) o indirectas (caída del valor de la deuda).
\item \textbf{Riesgo operacional}: pérdidas por errores humanos, problemas legales, fallas de sistemas o procesos ineficientes, como una operación mal ejecutada o un modelo mal evaluado.
\end{itemize}
\end{definicion}

\begin{figure}[H]
  \centering
  \includegraphics[width=0.5\linewidth,height=0.42\textheight,keepaspectratio]{FCA-3}
  \caption{Riesgos que afectan a un portafolio.}
\end{figure}

\paragraph{¿Cómo cuantifico el riesgo de mis portafolios?} Cada clase de activo tiene sus medidas de sensibilidad. En acciones, el \textbf{beta} (sensibilidad al mercado; mayor que 1, más volátil que el mercado) y, como indicio de valuación y de expectativas, la razón \textbf{P/E} \vease{sec:u4-acciones}. En renta fija, la \textbf{\emph{duration}} (sensibilidad del precio a la tasa: cuanto mayor, más sensible) y la \textbf{convexidad}, que mide la curvatura de la relación precio-tasa y gana importancia ante movimientos grandes. En derivados, la \textbf{delta} (cuánto cambia el precio de una opción por cada peso de cambio del subyacente: con delta 0,5, cincuenta centavos) y la \textbf{gamma} (cuánto cambia la delta cuando cambia el subyacente).

\begin{figure}[H]
  \centering
  \includegraphics[width=0.5\linewidth,height=0.42\textheight,keepaspectratio]{FCA-4}
  \caption{Medidas de riesgo según la clase de activo.}
\end{figure}

\begin{definicion}{Gestión del riesgo}
Es la práctica integral de \textbf{identificar} los factores de riesgo, \textbf{medir} cómo afectan al portafolio, \textbf{evaluar} para decidir y \textbf{mitigar} lo que se quiere proteger. La incertidumbre, expresada en factores de riesgo, obliga a resolver la valuación de los activos, la medición y consolidación del riesgo y su mitigación.
\end{definicion}
\end{document}